% concatenated from products-RiemConf-MOFM.tex
\documentclass[a4paper,10.9pt]{amsart}

\DeclareRobustCommand{\SkipTocEntry}[5]{}
%\setcounter{tocdepth}{3}
%\setcounter{secnumdepth}{3}

\usepackage{shadow}


%\usepackage{refcheck}

%\usepackage{showkeys}
\usepackage{accents}
\usepackage{marvosym}

\usepackage[T1]{fontenc}
\usepackage{empheq}
\usepackage{relsize}
\usepackage{amsmath}
\newcommand\bmmax{2}
\usepackage{bm}

\usepackage{upgreek}
\usepackage{ esint }
\usepackage{color}
\usepackage{comment}
\usepackage{amssymb}
\usepackage{amsfonts}
\usepackage{graphicx}

\usepackage{amscd}

\usepackage{slashed}
\usepackage{pdflscape}
\usepackage{tikz}
\usepackage{enumerate}
\usepackage{multirow}
\usepackage{stmaryrd}
\usepackage{cancel}
\usepackage{t1enc}

\usepackage{amssymb}
\usepackage{amscd}
\usepackage{scalerel,stackengine}

\usepackage{ifthen}

\usepackage{pdfsync}

\usepackage{graphicx} 
\usepackage{amsmath} 
\usepackage{amsfonts}
\usepackage{amssymb}


\usepackage{bbold}
\usepackage[linkcolor=blue]{hyperref}
\usepackage{mathrsfs}
\usepackage[colorinlistoftodos]{todonotes}
\usepackage{ytableau}

\setlength{\marginparwidth}{2.5cm}

\definecolor{blue}{rgb}{.255,.41,.884} % RoyalBlue of svgnames
\definecolor{red}{rgb}{1, 0, 0} % Red of svgnames
\definecolor{green}{rgb}{.196,.804,.196} % LimeGreen of svgnames
\definecolor{yellow}{rgb}{1,.648,0} % Orange of svgnames
\definecolor{pink}{rgb}{1,0.5,0.5}


\newcommand{\sam}[1]{\todo[size=\tiny,color=red!30]{#1 \\ \hfill --- Sam}}
\newcommand{\Sam}[1]{\todo[size=\tiny,inline,color=red!30]{#1
     \\ \hfill --- Sam}}


\setlength{\textwidth}{418pt}
\setlength{\oddsidemargin}{17.5pt}
\setlength{\evensidemargin}{17.5pt}

%\newcommand{\bm}[1]{\mbox{\boldmath~$ #1~$}}

\newtheorem{theorem}{Theorem}[section]
\newtheorem{lemma}[theorem]{Lemma}
  
\newtheorem{proposition}[theorem]{Proposition}
\newtheorem{corollary}[theorem]{Corollary}
\newtheorem{conjecture}[theorem]{Conjecture}

\theoremstyle{definition}
\newtheorem{definition}[theorem]{Definition}
\newtheorem{example}[theorem]{Example}
\newtheorem{xca}[theorem]{Exercise}

\theoremstyle{remark}
\newtheorem{remark}[theorem]{Remark}

\newtheorem{problem}[theorem]{Problem}



\newcommand{\boldnabla}{\mbox{\boldmath$ \nabla$}}
\newcommand{\al}{\mbox{\boldmath$\Delta$}}


\newcommand{\II}{{\rm  I\hspace{-.2mm}I}}
\newcommand{\IIo}{\hspace{0.4mm}\mathring{\rm{ I\hspace{-.2mm} I}}{\hspace{.0mm}}}
\newcommand{\iio}{\hspace{0.4mm}\mathring{{ \iota\hspace{-.41mm} \iota}}{\hspace{.0mm}}}
\newcommand{\Fo}{ {\hspace{.6mm}} \mathring{\!{ F}}{\hspace{.2mm}}}
\newcommand{\III}{{{{\bf\rm I\hspace{-.2mm} I \hspace{-.2mm} I}}{\hspace{.2mm}}}{}}
\newcommand{\IIIo}{{\mathring{{\bf\rm I\hspace{-.2mm} I \hspace{-.2mm} I}}{\hspace{.2mm}}}{}}
\newcommand{\IV}{{{{\bf\rm I\hspace{-.2mm} V}}{\hspace{.2mm}}}{}}
\newcommand{\IVo}{{\mathring{{\bf\rm I\hspace{-.2mm} V}}{\hspace{.2mm}}}{}}
\newcommand{\IVnn}{{{\bf\rm I\hspace{-.2mm} V}{\hspace{.2mm}}}{}}
\newcommand{\V}{{{\underline{\overline{\bf\rm V}}}}{}}
\newcommand{\Vo}{{\mathring{{\bf\rm V}}}{}}
\newcommand{\VIo}{{\mathring{{\bf\rm V\hspace{-.2mm}I}}{\hspace{.2mm}}}{}}
\newcommand{\VIIo}{{\mathring{{\bf\rm  V\hspace{-.2mm}I\hspace{-.2mm}I}}{\hspace{.2mm}}}{}}


\newcommand{\otop}{\mathring{\top}}
\newcommand{\ID}{{{{\bf\rm I\hspace{-.3mm} D}}{\hspace{0.1mm}}}}
\newcommand{\IDh}{{{{\bf\rm I\hspace{-.3mm} \hat{D}}}{\hspace{0.1mm}}}}


\hyphenation{Pa-wel}
\hyphenation{Nu-row-ski}
\hyphenation{And-rzej}
\hyphenation{Traut-man}
\hyphenation{Je-rzy}
\hyphenation{Le-wan-dow-ski}
\hyphenation{Car-tan}
\hyphenation{Car-tan-Pet-rov-Pen-rose}
\hyphenation{Pen-rose}
\hyphenation{or-tho-go-nal}
\hyphenation{comp-lex}
\hyphenation{Pet-rov}
\hyphenation{Euc-lid-ean}
\hyphenation{ge-om-etry}
\hyphenation{Rie-man-nian}
\hyphenation{Ein-stein}
\hyphenation{Ka-te-dra} 
\hyphenation{Me-tod} 
\hyphenation{Ma-te-ma-tycz-nych}
\hyphenation{Fi-zy-ki}
\hyphenation{Uni-wer-sy-tet} 
\hyphenation{War-szaw-ski} 
\hyphenation{War-sza-wa}


\newcommand{\aR}{\mbox{\boldmath{$ R$}}}
\newcommand{\aS}{\mbox{\boldmath{$ S$}}}
\newcommand{\aD}{\mbox{\boldmath{$ D$}}\hspace{-.2mm}}
\newcommand{\aRic}{\mbox{\boldmath{$ \Ric$}}}
\newcommand{\aLap}{\mbox{\boldmath{$ \Delta$}}}

\renewcommand{\colon}{\scalebox{1.2}{:}}



% sideremark
\def\sideremark#1{\ifvmode\leavevmode\fi\vadjust{\vbox to0pt{\vss
 \hbox to 0pt{\hskip\hsize\hskip1em
 \vbox{\hsize2cm\tiny\raggedright\pretolerance10000
  \noindent #1\hfill}\hss}\vbox to8pt{\vfil}\vss}}}

\newcommand{\edz}[1]{\sideremark{#1}}
\def\idx#1{{\em #1\/}} %


\numberwithin{equation}{section}


\newcommand{\hh}{{\hspace{.3mm}}}

\renewcommand\colon{\scalebox{1.3}{$:$}}

\newcommand{\nablab}{\bar\nabla}
\newcommand{\cc}{\boldsymbol{c}}
\newcommand{\pdot}{{\boldsymbol{\cdot}}}

%equal along Sigma
\renewcommand{\=}{\stackrel \Sigma =}
\newcommand{\eqU}{\stackrel{U_{\Lambda}}{=}}

\newcommand{\Jb}{\bar \J}
\newcommand{\Pb}{\bar P}
\newcommand{\Rb}{\bar R}
\newcommand{\bn}{\boldsymbol{n}}
\newcommand{\bG}{\boldsymbol{G}}
\newcommand{\bg}{\boldsymbol{g}}
\newcommand{\bLap}{\bar{\Delta}}

%tangential derivative
\newcommand{\nablat}{\nabla^\top}
%normal derivatives
\newcommand{\nablan}{\nabla_n}
\newcommand{\nablanh}{\nabla_{\nh}}

%trace
\DeclareMathOperator{\tr}{tr}
\DeclareMathOperator{\divergence}{div}

\renewcommand\geq{\geqslant}
\renewcommand\leq{\leqslant}


\newcommand{\mh}{{}\hspace{-.3mm}}

\stackMath
\newcommand\reallywidehat[1]{%
\savestack{\tmpbox}{\stretchto{%
  \scaleto{%
    \scalerel*[\widthof{\ensuremath{#1}}]{\kern-.6pt\bigwedge\kern-.6pt}%
    {\rule[-\textheight/2]{1ex}{\textheight}}%WIDTH-LIMITED BIG WEDGE
  }{\textheight}% 
}{0.5ex}}%
\stackon[1pt]{#1}{\tmpbox}%
}

\newcommand\eqSig{ \mathrel{\overset{\makebox[0pt]{\mbox{\normalfont\tiny\sffamily~$\Sigma$}}}{=}} }
\renewcommand\eqSig{\mathrel{\stackrel{\Sigma\hh}{=}} }
\newcommand\eqLam{\mathrel{\stackrel{\Lambda\hh}{=}}}
\newcommand{\hd }{\hat{D}}
\newcommand{\hdb}{\hat{\bar{D}}}
\newcommand{\csdot}{\hspace{-0.75mm} \cdot \hspace{-0.75mm}}
\newcommand{\IdD}{(I \csdot \hd)}
\newcommand{\bdot }{\mathop{\lower0.33ex\hbox{\LARGE$\cdot$}}}

\definecolor{ao}{rgb}{0.0,0.0,1.0}
\definecolor{forest}{rgb}{0.0,0.3,0.0}
\definecolor{red}{rgb}{0.8, 0.0, 0.0}

\newcommand{\APE}[1]{{\rm APE}_{#1}}
\newcommand{\PE}{{\rm PE}}
\newcommand{\FFo}[1]{\mathring{\underline{\overline{\rm{#1}}}}}
\newcommand{\FF}[1]{\underline{\overline{\rm{#1}}}}
\newcommand{\ltots}[1]{{\rm ltots}_{#1}}
\newcommand{\tsim}[1]{\stackrel{#1}{\sim}}
\newcommand{\ce}{\mathcal{E}}
\newcommand{\ct}{\mathcal{T}}
%\color{ao}

\DeclareMathOperator{\Gram}{Gram}
\DeclareMathOperator{\Id}{Id}




\begin{document}



\subjclass[2010]{
53C18, 53A55, 53C21, 58J32.
}

\renewcommand{\today}{}
\title{
{Higher Fundamental Forms and Warped Product Hypersurfaces
}}
%\author{ Cesar Arias${}^\flat$, A. Rod Gover${}^\sharp$ \&  Andrew Waldron${}^\natural$}
%
%\address{${}^\sharp$Department of Mathematics\\
%  The University of Auckland\\
%  Private Bag 92019\\
%  Auckland 1142\\
%  New Zealand,  and\\
%  Mathematical Sciences Institute, Australian National University, ACT 
%  0200, Australia} \email{gover@math.auckland.ac.nz}
%  
%  \address{${}^{\natural}$
%  Center for Quantum Mathematics and Physics (QMAP)\\
%  Department of Mathematics\\ 
%  University of California\\
%  Davis, CA95616, USA} \email{wally@math.ucdavis.edu}
%  
% \Andrew{Better title?} 

\author{ Samuel Blitz${}^\flat$ \& Josef \v{S}ilhan${}^\sharp$}

\address{${}^\flat$
 Department of Mathematics and Statistics \\
 Masaryk University\\
 Building 08, Kotl\'a\v{r}sk\'a 2 \\
 Brno, CZ 61137} 
   \email{blitz@math.muni.cz}

\address{${}^\sharp$
 Department of Mathematics and Statistics \\
 Masaryk University\\
 Building 08, Kotl\'a\v{r}sk\'a 2 \\
 Brno, CZ 61137} 
   \email{silhan@math.muni.cz}
 
\vspace{10pt}

\begin{abstract}
Warped products are one of the simplest families of Riemannian manifolds that can have non-trivial geometries. In this article, we characterize the geometry of hypersurface embeddings arising from warped product manifolds using the language of higher (Riemannian) fundamental forms. In a similar vein, we also study the geometry of conformal manifolds with embedded hypersurfaces that admits a trivialization of the conformal metric to a product metric, with base manifold given by the embedded hypersurface. We show that the higher conformal fundamental forms play a critical role in their characterization.

\vspace{0.5cm}

\noindent
%\begin{keywords}
{\sf \tiny Keywords: 
Riemannian geometry, conformal geometry, submanifold embeddings, holography}
 %\end{keywords}

\vspace{2cm}

\end{abstract}

%58E30  	Variational principles
%49S05  	Variational principles of physics
%51P05  	Geometry and physics
%53A30  	Conformal differential geometry
%53A55  	Differential invariants (local theory), geometric objects
%53C80  	Applications to physics (Global differential geoemtry)
%53B50  	Applications to physics (Local differential geometry)
%53A07  	Higher-dimensional and -codimensional surfaces in Euclidean~$n$-space
%53C21  	Methods of Riemannian geometry, including PDE methods; curvature restrictions
%53B15  	Other connections
%83C99  	None of the above, but in this section (General relativity)

\maketitle

\pagestyle{myheadings} \markboth{S. Blitz \& J. \v{S}ilhan}{On Fundamental Forms and Riemannian and Conformal Warped Products}

%\newpage


\tableofcontents

\section*{Acknowledgements}
SB is supported by the Operational Programme Research Development and Education Project No. CZ.02.01.01/00/22-010/0007541. J\v{S} is supported by the Czech Science Foundation (GACR) grant GA22-00091S.
%SB and J\v{S} would also like to acknowledge Gavin Pandya for valuable insight in proving XYZ.

\section*{Data Availability}
We do not analyse or generate any datasets, because our work proceeds within a theoretical and mathematical approach. One can obtain the relevant materials from the references below.

%\section{Main}
%
%OUTLINE
%\begin{itemize}
%\item Definitions of higher fundamental forms
%\item Warm up - geometric meaning of second fundamental form
%\item Geometric meaning of third fundamental form
%\item Geometric meaning of fourth fundamental form
%\item Generalization: Obstruction flat
%\item Generalization: Higher fundamental forms obstruct normal form
%\end{itemize}

\section{Introduction}
The notion of a \textit{warped product} manifold was introduced by Bishop and O'Neill in 1969~\cite{bishoponeill}. Given two (pseudo-)Riemannian manifolds $(F,g_B)$ and $(B,g_B)$, a warped product manifold $F \times_f B$ is a manifold $F \times B$ equipped with a (pseudo-)Riemannian metric
\begin{equation} \label{warped}
g = f^2 g_F + g_B\,,
\end{equation}
where $f  \in C^\infty_+(B)$. Here, $B$ is called the \textit{base} manifold and $F$ is called the \textit{fiber} manifold. When the warping function $f$ is constant, the manifold is called a \textit{product manifold}. In general, warped product manifolds are emblematic of some of the simplest non-trivial geometries, and as such find broad application in general relativity. For example, both the Schwarzschild metric and the FLRW metric are warped product manifolds and are physically relevant solutions to the Einstein field equations. Indeed, partial differential equations are separable on these manifolds, and hence provide a fertile ground for solving partial differential equations in non-trivial settings.

The distinction between warped product and product manifolds is lost in the conformal geometric setting. A conformal manifold $(M,\cc)$ is a smooth manifold equipped with an equivalence class of metrics $\cc = [g]$, where equivalence between two metrics $g \sim g'$ holds when there exists some $\Omega \in C^\infty_+ (M)$ such that $g' = \Omega^2 g$. Thus, it is easy to see that if a warped product metric is an element of the conformal class, then so is a product manifold, as $f^{-2} \in C^\infty_+ (M)$.

Now, as (warped) product manifolds distinguish a family of submanifolds, it is clear that these distinguished submanifold embeddings should be have distinct geometric features, both intrinsically and extrinsically. To capture this behavior, we specifically study geometric tensor invariants of the embeddings.

The first example of such a tensor invariant was the second fundamental form, introduced by Gauss in the 19th century. In 1867, Bonnet proved the fundamental theorem of surfaces, which, in modern language, asserts that up to rigid Euclidean motions, a surface embedded in $\mathbb{R}^3$ is uniquely determined by its first and second (Riemannian) fundamental forms. This result easily generalizes to hypersurfaces embedded in $\mathbb{R}^d$. The naming of these tensors begs for explanation: is there a third fundamental form? A fourth? Indeed there are---but in the context of Euclidean embeddings, they either vanish or are easily expressible in terms of the first and second fundamental forms~\cite[Volume 3]{Spivak} (depending on the definition provided). However, recent work~\cite{BGW1} found that higher fundamental forms in a conformal setting, which capture information about the hypersurface embedding, have interesting geometric interpretations. It is the goal of this article to show that the Riemannian and conformal fundamental forms play special roles in characterizing (warped) product manifolds that arise from hypersurface embeddings in the Riemannian and conformal settings, respectively.

In Section~\ref{sec:riem}, we introduce and investigate the Riemannian fundamental forms in Riemannian product and warped product manifolds. In Section~\ref{sec:conf}, we remind the reader of conformal fundamental forms and study their behavior when a conformal hypersurface embedding admits a product manifold in its conformal class.

\subsection{Notations and Conventions}

In this article, we denote by $M$ a smooth oriented manifold of dimension $d \geq 3$ and by $\Sigma$ a $d-1$ dimensional manifold, smoothly embedded via $\iota$ into $M$. We use the overline notation $\bar{\bullet}$ to denote quantities associated with the hypersurface $\Sigma$. After equipping $M$ with a metric $g$, there exists a unique Levi-Civita connection $\nabla$ on the tangent bundle $TM$ (and its dual and tensor products). The Riemann curvature tensor of the connection is denoted by $R$ and is defined according to
$$R(x,y)z = (\nabla_x \nabla_y - \nabla_y \nabla_x - \nabla_{[x,y]}) z\,,$$
where $x,y,z \in \Gamma(TM)$ are sections of the tangent bundle and $[x,y]$ is their Lie bracket.

Throughout this manuscript, we use abstract index notation to denote sections of the tangent bundle, the cotangent bundle, and their tensor products; for example, we may write $t^a{}_{bc} \in \Gamma(TM \otimes T^* M \otimes T^* M)$, and the Einstein summation convention is used to denote contraction. We use $\wedge^k$ to denote the antisymmetric $k$th tensor power of a vector bundle and square brackets $[]$ around indices to denote antisymmetrization, $\odot^k$ to denote symmetric $k$th power of the same and round brackets $()$ around indices to denote symmetrization, and often use the subscript $\circ$ (as in $t_{(ab)_{\circ}} \in \Gamma(\odot^2_{\circ} T^* M)$) to denote the trace-free part of a section or a bundle. Furthermore, we often use subscripts to denote contraction; for example, given a vector $n \in \Gamma(TM)$ and a tensor $t_{abc} \in \Gamma(\otimes^3 T^* M)$, we denote the contraction $n^a t_{abc} =: t_{nbc} \in \Gamma(\otimes^2 T^* M)$.

Using these notations, we have that
$$x^a y^b R_{ab}{}^c{}_d z^d =x^a y^b [\nabla_a, \nabla_b] z^c\,.$$
The trace of the Riemann curvature tensor, the Ricci tensor, is denoted by $Ric_{ab} := R_{ca}{}^c{}_b$ and the scalar curvature is denoted by $Sc := g^{ab} Ric_{ab}$. The trace-free part of the Riemann curvature tensor is the Weyl tensor and is defined according to
$$W_{abcd} := R_{abcd} - g_{ac} P_{bd} + g_{ad} P_{bc} + g_{bc} P_{ad} - g_{bd} P_{ac}\,,$$
where the Schouten tensor $P$ is defined according to
$$P_{ab} := \tfrac{1}{d-2} \left(Ric_{ab} - \tfrac{1}{2(d-1)} Sc \, g_{ab} \right)\,,$$
with $J := g^{ab} P_{ab}$.
For later use, we also provide the formula for two well-known higher-order curvature quantities: the Cotton tensor and the Bach tensor, respectively:
\begin{align*}
C_{abc} :=& \nabla_a P_{bc} - \nabla_b P_{ac} \\
B_{ab} :=& \Delta P_{ab} - \nabla^c \nabla_a P_{bc} + P^{cd} W_{acbd}\,,
\end{align*}
where $\Delta :=\nabla^a \nabla_a $ is the negative Laplace operator.

\section{Riemannian Fundamental Forms and Warped Products} \label{sec:riem}
\subsection{Riemannian Fundamental Forms and Hypersurface Embeddings}
We begin by defining Riemannian fundamental forms:
\begin{definition} \label{ff-def}
Let $\iota : \Sigma \hookrightarrow (M,g)$ be a smooth hypersurface embedding and let $n \in \Gamma(T M)$ be any vector orthogonal $T\Sigma$ such that $g(n,n)|_{\Sigma} = 1$. The \textit{first fundamental form} $\bar{g} := \iota^* g$ is the pullback of $g$ to the hypersurface; this is also called the \textit{induced metric}. The \textit{second fundamental form} is defined according to $\II := \iota^* \nabla g(n,\cdot)$. We define the \textit{third fundamental form} according to $\III := \iota^* g(R(n,\cdot) n, \cdot)$, and finally, for $k \geq 4$, we define the \textit{$k$th fundamental form} using abstract index notation:
$$\FF{k}_{ab} := \iota^* (n^{a_1} \cdots n^{a_{k-3}} n^c n^d \nabla_{a_1} \cdots \nabla_{a_{k-3}} R_{cabd})\,.$$
\end{definition}
As is shown in Lemma~\ref{normal-derivs}, this definition provides a natural extension of the second fundamental form, insofar as they are related to normal derivatives of $\nabla n$. Observe that a Riemannian fundamental form $\FF{k}$ has \textit{transverse order} $k-1$, meaning that it depends on no fewer and no more than $k-1$ normal derivatives of the metric.

Generally speaking, the unique feature of the Riemannian fundamental forms is their leading behavior in terms of normal derivatives of the metric---one may always add powers of the second fundamental form at lower order without ruining this behavior. As such, there are alternative forms of the Riemannian fundamental forms that are useful in other contexts. To show leading-order equivalence of these alternative forms, we first need a special set of coordinates near a hypersurface.

In general, for any hypersurface embedding $\Sigma \hookrightarrow (M,g)$, we may express the metric (in a collar neighborhood) in the \textit{normal form}, given by 
$$g = dt^2 + \bar{g}(t,x^k)_{ij} dx^i dx^j\,,$$
where $\bar{g}(t_0,x^k)$ is the induced metric on the hypersurface defined by $t = t_0$ and coordinates $\{x^i\}$ are defined by parallel transporting coordinates on $\Sigma$ along geodesics generated by $\partial_t$. 
The 1-form $dt$ is called the \textit{unit conormal 1-form} and the dual vector field will be usually denoted by $n$; 
it provides a canonical extension of 
the unit normal vector along $\Sigma$ to $M$.
This normal form is quite useful, as it allows us to easily establish alternative formulas for the fundamental forms. Indeed, we have the following technical lemmas.
\begin{lemma} \label{normal-derivs}
Let $\iota : \Sigma \hookrightarrow (M,g)$ be a smooth embedded hypersurface with the vector field $n$ dual to 
the unit conormal 1-form. Then, for every $k \geq 2$, we have that
$$\FF{k} = \iota^* \nabla_n^{k-2} \nabla n + \text{products of lower order fundamental forms.}$$
\end{lemma}
\begin{proof}
We prove this lemma by induction. Clearly, when $k = 2$, the identity is true by explicit computation. Now, suppose that the identity holds for all $2 \leq n \leq k$. Now, we compute:
\begin{align*}
\iota^* \nabla_{n}^{k-1} \nabla_a n_b =& \iota^*  \nabla_n^{k-2} n^c \nabla_a \nabla_c n_b + \iota^* \nabla_n^{k-2} n^c R_{cabd} n^d \\
=& \iota^*\nabla_n^{k-2} \nabla_a \nabla_n n_b - \iota^*\nabla_n^{k-2} [(\nabla_a n^c)(\nabla_c n_b)] + \iota^*\nabla_n^{k-2} R_{nabn} \\
=& \iota^*\nabla_n^{k-2} R_{nabn} - \iota^*\nabla_n^{k-2} [(\nabla_a n^c)(\nabla_c n_b)] \\
=& \iota^* n^c n^d n^{a_1} \cdots n^{a_{k-2}} \nabla_{a_1} \cdots \nabla_{a_{k-2}} R_{cabd} - \iota^*\nabla_n^{k-2} [(\nabla_a n^c)(\nabla_c n_b)] \\
=& \FF{k+1}{}_{ab} -  \iota^*\nabla_n^{k-2} [(\nabla_a n^c)(\nabla_c n_b)]\,.
\end{align*}
Note that the third and fourth identities follow because, from the existence of the normal form, there exists a function $t$ such that $g(n,\cdot) = dt$, and so
$$\nabla_n n_a = n^c \nabla_c \nabla_a t = n^c \nabla_a n_c = \tfrac{1}{2} \nabla_a n^2 = 0\,.$$
Now by the Leibniz property and the induction hypothesis, it follows that $\iota^*\nabla_n^{k-2} [(\nabla_a n^c)(\nabla_c n_b)]$ can be expressed along $\Sigma$ as a sum of products of lower order fundamental forms. The lemma follows.
\end{proof}
\begin{remark}
Note that the definition of the third fundamental form given in~\cite[Volume 3]{Spivak} is $-\iota^* \nabla_n \nabla_a n_b$, and so Lemma~\ref{normal-derivs} shows that Definition~\ref{ff-def} is a reasonable extension to higher fundamental forms.
\end{remark}

Using this result, we have another straightforward lemma.  Here $\mathcal{L}$ denotes the Lie derivative.
\begin{lemma} \label{lie-derivs}
Let $\iota : \Sigma \hookrightarrow (M,g)$ be a smooth embedded hypersurface with the vector field $n$ dual to 
the unit conormal 1-form. Then, for every $k \geq 2$, we have that
$$\FF{k} = \tfrac{1}{2} \iota^* \mathcal{L}_n^{k-1} g+ \text{products of lower order fundamental forms.}$$
\end{lemma}
\begin{proof}
We prove by induction. When $k = 2$, the result holds by definition.
Next, suppose that the identity holds for all $2 \leq n \leq k$. Then,
\begin{align*}
\tfrac{1}{2} \iota^* \mathcal{L}_n^{k} g_{ab} =& \tfrac{1}{2} \iota^* \mathcal{L}_n \mathcal{L}_n^{k-1} g_{ab} \\
=& \tfrac{1}{2} \iota^* \nabla_n \mathcal{L}_n^{k-1} g_{ab} + \iota^* (\mathcal{L}^{k-1}_n g)_{c(a} (\nabla_{b)} n^c) \\
=& \tfrac{1}{2} \iota^* \nabla_n \mathcal{L}_n^{k-1} g_{ab} + \FF{k}_{c(a} \II_{b)}^c + \text{products of lower order fundamental forms} \\
=& \iota^* \nabla_n^{k-1} \nabla_a n_b + \text{products of lower order fundamental forms.}
\end{align*}
But from Lemma~\ref{normal-derivs}, the last line can be expressed in terms of $\FF{k+1}$, and so induction follows.
\end{proof}

\begin{remark}\label{weights}
Note that in all of these constructions involving equality modulo lower order terms, we may ascribe weights to better control which fundamental forms appear in each identity. Indeed, assigning a weight of $3-k$ to $\FF{k}$ and a weight of $-2$ for every contraction, it is clear that the weights of each summand in the above identities are always equal.

Additionally, observe that when the lower fundamental forms vanish, the original definition of the fundamental forms is equivalent to the constructions provided in Lemmas~\ref{normal-derivs} and~\ref{lie-derivs}.
\end{remark}


Using Lemma~\ref{lie-derivs}, for any fundamental form, one may expres
 it as a sum of Lie derivatives of the metric and their contractions:
\begin{corollary} \label{cor-lies}
Let $\iota : \Sigma \hookrightarrow (M,g)$ be a smooth embedded hypersurface with the vector field $n$ dual to 
the unit conormal 1-form. Then for every $2 \leq k$, there exists a homogeneous formula for $\FF{k}$ given by
$$\FF{k} =\tfrac{1}{2} \iota^* \mathcal{L}_n^{k-1} g +  F_k(\bar{g}^{-1}, \iota^* \mathcal{L}_n g, \ldots, \iota^* \mathcal{L}_n^{k-2} g)\,.$$
\end{corollary}
These formulas can be explicitly computed for any choice of $k$ simply by a recursive algorithm.

Lemma~\ref{lie-derivs} also implies that the metric in a sufficiently small neighborhood around a point in $\Sigma$ 
(where $(M,g)$ is analytic and $\Sigma$ is analytically embedded) is uniquely determined by the induced metric on $\Sigma$ and the infinite family of fundamental forms for the embedding:
\begin{proposition} \label{unique}
Let $\Sigma \hookrightarrow (M,g)$ be an analytic hypersurface embedding. Then, for any point $p \in \Sigma$, there exists a sufficiently small neighborhood $U_p$ around $p$ such that $g|_U$ is uniquely determined (up to diffeomorphism) by the induced metric and the infinite family of fundamental forms.
\end{proposition}
\begin{proof}
Using Corollary~\ref{cor-lies}, one may invert the formula for each $\FF{m+1}$ to find a formula for each Lie derivative of $\iota^* \mathcal{L}_n^m g$ in terms of the $\{\bar{g}^{-1}, \II, \ldots, \FF{m+1}\}$. Now, such an embedding may always be expressed in terms of a normal form (via an isometry to a collar neighborhood), so that
$$g = dt^2 + \bar{g}(t,x^k)_{ij} dx^i dx^j\,,$$
where the component functions $\bar{g}(t,x^k)_{ij}$ are analytic. But then $\bar{g}(t,x^k)$ is uniquely determined (in a sufficiently small neighborhood around a point $p \in \Sigma$) by its Taylor series expansion about $t = 0$. But because $n = \partial_t$ and $\partial_a \partial_t = 0$, it follows that $\iota^* \mathcal{L}_{n}^m g = (\partial_t^{m} \bar{g})(0,x^k)_{ij} dx^i dx^j$, and so $\bar{g}(t,x^k)$ is determined uniquely (in a sufficiently small neighborhood) by the family of tensors $\{\iota^* \mathcal{L}_n^m g\}_{m=0}^{\infty}$. However, as established, there exists a formula for each of these Lie derivatives in terms of $\bar{g}^{-1}$ and the fundamental forms. The proposition follows.
\end{proof}

This uniqueness result is quite powerful, as it allows us to compare the fundamental forms of hypersurface embeddings to determine whether those hypersurface embeddings are equivalent.
\begin{theorem} \label{equal-metrics}
Let $\iota : \Sigma \hookrightarrow (M,g)$ and $\hat{\iota} : \Sigma \hookrightarrow (M,\hat{g})$ be two analytic hypersurface embeddings such that, for every $p \in \Sigma$,  $\iota(p) = \hat{\iota}(p)$. Then, there is an analytic diffeomorphism $\phi$ defined in a collar neighborhood around $\Sigma$ that fixes $\Sigma$ pointwise such that $\phi^* \hat{g} = g$ near $\Sigma$ if and only if $\iota^* g = \hat{\iota}^* \hat{g}$ and, for every $k$, we have that $\FF{k} = \hat{\FF{k}}$.
\end{theorem}
\begin{proof}
First, suppose that a diffeomorphism $\phi$ fixing $\Sigma$ pointwise exists so that $\phi^* \hat{g} = g$. Then it follows trivially that all geometric quantities are equal, including the induced metric and the fundamental forms. On the other hand, the other direction follows directly from the uniqueness of Proposition~\ref{unique} and the existence of the normal form.
\end{proof}

As a consequence of Theorem~\ref{equal-metrics}, we can determine when a hypersurface embedding is actually a product manifold.
\begin{corollary} \label{prod-metric}
Let $\Sigma \hookrightarrow (M,g)$ be an analytic hypersurface embedding. Then $\FF{k} = 0$ for every $k \geq 2$ if and only if in a collar neighborhood $U := (-\epsilon,\epsilon) \times \Sigma$ of $\Sigma$, $(U,g|_U)$ is a product manifold.
\end{corollary}
\begin{proof}
We define a product manifold on $(-\epsilon,\epsilon) \times \Sigma$ by the metric $g = dt^2 + \iota^* g$. Clearly, this product manifold has vanishing fundamental forms and the induced metric on $\Sigma$ agrees with that of $\Sigma \hookrightarrow (M,g)$. Hence by Theorem~\ref{equal-metrics}, these metrics agree. The other direction is trivial.
\end{proof}

\subsection{Fiber-like Hypersurface Embeddings} \label{fiber}
Now, when a hypersurface $\Sigma \hookrightarrow (M,g)$ can be viewed as the fiber in a warped product, there exists a choice of local coordinates on $M$ such that the metric takes the form
$$g = dt^2 + h(t) \bar{g}_{ij} dx^i dx^j\,,$$
where $\Sigma = \mathcal{Z}(t)$ is zero locus of $t$, 
coordinates $\{x^i\}$ are coordinates on $\Sigma$ parallel transported by $\partial_t$, 
$\bar{g}_{ij}$ depends only on $x^i$,
and $h \in C^\infty_+ (\mathbb{R})$. We call such a hypersurface embedding a \textit{fiber-like hypersurface embedding}.
Note this corresponds to  \eqref{warped} with $h = f^2$ and 1-dimensional base manifold.

We can then use the results of the previous section to determine when a hypersurface embedding is indeed a fiber-like hypersurface embedding. We thus have one of the main theorems of this article.
\begin{theorem} \label{fiber-theorem}
Let $\Sigma \hookrightarrow (M,g)$ be an analytic hypersurface embedding with $\bar{g} := \iota^* g$. Then, a collar neighborhood around $\Sigma$ is isometric to a fiber-like hypersurface embedding if and only if there exists analytic $h \in C^\infty_+(\mathbb{R})$ such that for every $k \geq 2$,
$$\FF{k} = \tfrac{1}{2} h^{(k-1)}(0) \, \bar{g} + F_k(\bar{g}^{-1}, h'(0) \,\bar{g}, \ldots, h^{(k-2)}(0)\, \bar{g})\,$$
for $F_k$ from Lemma \ref{cor-lies}. That is, $\FF{k} = H\bar{g}$ where the constant $H$ is given by
$$ H = \tfrac{1}{2} h^{(k-1)}(0) + \widetilde{F}_k(1, h'(0), \ldots, h^{(k-2)}(0))  $$
where $\widetilde{F}_k$ is obtained from ${F}_k$ by replacing contraction of 2-tensor by products of corresponding constants.
\end{theorem}
\begin{proof}
First, suppose that $\Sigma \hookrightarrow (M,g)$ is an analytic fiber-like hypersurface embedding, meaning there exists some analytic $h \in C^\infty_+(\mathbb{R})$ such that, in the normal form,
$$g = dt^2 + h(t) \bar{g}_{ij} dx^i dx^j\,.$$
Clearly, because $n = \partial_t$, it follows that $\iota^* \mathcal{L}_n^{m} g = \iota^* (\partial_t^{m} h) \bar{g}$.
But from Corollary~\ref{cor-lies}, it follows that
$$\FF{k} =\tfrac{1}{2} \iota^* \mathcal{L}_n^{k-1} g +  F_k(\bar{g}^{-1}, \iota^* \mathcal{L}_n g, \ldots, \iota^* \mathcal{L}_n^{k-2} g) = \tfrac{1}{2} h^{(k-1)}(0)\, \bar{g} +  F_k(\bar{g}^{-1}, h'(0)\, \bar{g}, \ldots, h^{(k-2)}(0) \, \bar{g})\,.$$
This gives us one direction.

On the other hand, suppose that there exists an analytic $h \in C^\infty_+(\mathbb{R})$ such that for every $k \geq 2$, 
$$\FF{k} = \tfrac{1}{2} h^{(k-1)}(0) \, \bar{g} + F_k(\bar{g}^{-1}, h'(0) \,\bar{g}, \ldots, h^{(k-2)}(0)\, \bar{g})\,.$$
Now, the metric $g$ is expressible in normal form as $g = dt^2 + \bar{g}(t,x^k)_{ij} dx^i dx^j$. So,  define a metric on a collar neighborhood of $\Sigma$ by $\hat{g} = dt^2 + h(t) \bar{g}(0,x^k)_{ij} dx^i dx^j$. However, from the above considerations, this metric $\hat{g}$ induces fundamental forms that agree exactly with those of the original embedding. Hence by Theorem~\ref{equal-metrics}, $\hat{g} = g$, and so $\Sigma \hookrightarrow (M,g)$ is isometric to a fiber-like hypersurface embedding in a collar neighborhood.
%On the other hand, suppose that there exists an analytic $f \in C^\infty_+(\mathbb{R})$ such that for every $k \geq 2$, 
%$$\FF{k} = \tfrac{1}{2} f^{(k-1)}(0) \, \bar{g} + F_k(\bar{g}^{-1}, f'(0) \,\bar{g}, \ldots, f^{(k-2)}(0)\, \bar{g})\,.$$
%Now, about every points in $\Sigma$, there exists a sufficiently small neighborhood equipped with a warped product metric $g = dt^2 + f(t) \bar{g}(0,x^k)_{ij} dx^i dx^j$ that shares the same fundamental forms as the given embedding by Corollary~\ref{cor-lies}. However, because both $f$ and the embedding are analytic, this metric must be unique from Proposition~\ref{unique}, and hence these two metrics must agree in a neighborhood about the point. But this holds for a neighborhood around any points in $\Sigma$, and hence holds for a collar neighborhood around $\Sigma$. The theorem follows.
\end{proof}

\begin{remark}
Explicitly, one may show that, in the language of the above theorem,
$$\III_{ab} = \tfrac{1}{2} h''(0) \bar{g}_{ab} - \tfrac{1}{4} \bar{g}^{cd} (h'(0) \bar{g}_{ac}) (h'(0) \bar{g}_{bd}) = \left[\tfrac{1}{2} h''(0) - \tfrac{1}{4} (h'(0))^2)\right] \bar{g}_{ab}\,. $$
\end{remark}





%%%In this section, we discuss higher fundamental forms more generally. For example, as mentioned in the introduction, depending on the definition used, fundamental forms are always expressible as powers of the seocnd fundamental form when the bulk manifold is flat. Using our definition of the fundamental forms, they trivially vanish. However, the result can be re-expressed as describing the behavior of normal derivatives of $\nabla n$, which would (typically) have terms including higher fundamental forms.
%%%\begin{proposition}
%%%Let $\Sigma \hookrightarrow (M^d,g)$ be an embedded hypersurface and let $(M,g)$ be a spaceform. Furthermore, let $n = ds$ be a unit conormal for $s$ a defining function of $\Sigma$. Then, for every positive integer $k$, $\nabla_n^k \nabla n$ is expressible as a sum of powers of the second fundamental form.
%%%\end{proposition}
%%%\begin{proof}
%%%We perform a simple calculation:
%%%\begin{align*}
%%%\nabla_n^k \nabla_a n_b =& \nabla_n^{k-1} n^c \nabla_a \nabla_c n_b + \nabla_n^{k-1} n^c R_{cabd} n^d \\
%%%=& \nabla_n^{k-1} \nabla_a \nabla_n n_b - \nabla_n^{k-1} [(\nabla_a n^c)(\nabla_c n_b)] + \nabla_n^{k-1} R_{nabn} \\
%%%\eqSig& \nabla_n^{k-1} R_{nabn} + \text{ products of lower order terms.}
%%%\end{align*}
%%%Because $(M,g)$ is a spaceform, it follows by induction that all such normal derivatives are expressible as products of the second fundamental form $\II$.
%%%%Next, observe that
%%%%\begin{align*}
%%%%\nabla_n^{k-3} R_{nbcn} \eqSig& n^a n^d n^e \nabla_e \nabla^{k-4}_n R_{abcd} + \text{ lower order terms containing $\nabla n$}\\
%%%%\eqSig& n^d (g^{ae} - \bar{g}^{ae}) \nabla_e \nabla^{k-4}_n R_{abcd} + \text{ lower order terms containing $\nabla n$} \\
%%%%\eqSig& n^d \nabla^a \nabla^{k-4}_n R_{abcd} - n^d \nabla^{\top a} \nabla^{k-4}_n R_{abcd} + \text{ lower order terms containing $\nabla n$} \\
%%%%\eqSig& n^d \nabla^{k-4}_n \nabla^a  R_{abcd} - n^d \nabla^{\top a} \nabla^{k-4}_n R_{abcd} + \text{ lower order terms containing $\nabla n$} \\
%%%%\eqSig& -n^d \nabla^{k-4}_n (\nabla_d R_{ab}{}^a{}_c + \nabla_c R_{abd}{}^a) - n^d \nabla^{\top a} \nabla^{k-4}_n R_{abcd} + \text{ lower order terms containing $\nabla n$} \\
%%%%\eqSig&- n^d \nabla^{\top a} \nabla^{k-4}_n R_{abcd} + \text{ lower order terms containing $\nabla n$} \\
%%%%\end{align*}
%%%
%%%\end{proof}
%%%
%%%\begin{remark} \label{weights}
%%%More generally, when $n = ds$ is a unit conormal for $s$ a defining function of the hypersurface, a similar calculation as that above shows that $\iota^* \nabla_n^k \nabla_a n_b \= \FF{k+2} + \text{subleading}$, where the subleading terms are contractions of lower fundamental forms. More specifically, assigning a weight of $3-m$ to $\FF{m}$ and a weight of $-2$ for every contraction, we find that each summand in an expression for $\iota^* \nabla_n^k \nabla n$ has weight $1-k$.
%%%\end{remark}
%%%
%%%Furthermore, the fundamental forms can be viewed as an obstruction to a manifold with embedded hypersurface being isomorphic to a product manifold:
%%%\begin{proposition} \label{prod-vanishingFFs}
%%%Let $\iota : \Sigma \hookrightarrow (M,g)$ be a smooth hypersurface embedding. Then, $\FF{k} = 0$ for every $k \geq 2$ if and only if in a collar neighborhood around $\Sigma$, $(M,g)$ is locally isometric to a product manifold $((-\epsilon,\epsilon) \times \Sigma, ds^2 + \iota^* g)$ for some $\epsilon > 0$.
%%%\end{proposition}
%%%\begin{proof}
%%%In a collar neighborhood of $\Sigma$, we may always choose coordinates $(s,\vec{x})$ such that $g = ds^2 + \bar{g}(s,\vec{x})$ and $\bar{g}$ is a metric on slices of constant $s$. Now, because the embedding is smooth, we may Taylor series expand $\bar{g}(s,\vec{x})$ about $s = 0$, so that 
%%%$$g = ds^2 + \bar{g}^{(0)}(\vec{x}) + s \bar{g}^{(1)}(\vec{x}) + \cdots\,.$$
%%%
%%%Now suppose that, for each $k \geq 2$, we have that $\FF{k} = 0$. Then, it follows that $\iota^* \mathcal{L}_{\partial_s}^k g \= 0$ \edz{need to prove this... maybe put it earlier?} for all $k \geq 1$. It therefore follows that $\bar{g}^{(k)}(\vec{x}) =0$ for each $k \geq 1$. Hence, we have that
%%%$$g = ds^2 + \bar{g}^{(0)}(\vec{x}) = ds^2 + \iota^* g\,.$$
%%%Thus, the claim of isometry follows.
%%%
%%%On the other hand, if $g = ds^2 + \bar{g}(\vec{x})$, then for the same reason, it follows that $\II = \tfrac{1}{2} \iota^* \mathcal{L}_n g = 0$. Furthermore, because the manifold is a product and the $\mathbb{R}$ prodand is flat, it follows that the curvature tensor $R = \bar{R}$. Thus, it follows that $R_{nabn} = \III = 0$. Finally, because the manifold is a product, the connection satisfies $\nabla_a = n_a \partial_s + \bar{\nabla}_a$. But $\partial_s R = 0$, so it follows that all higher fundamental forms vanish as well.
%%%\end{proof}
%%%
%%%\begin{remark}
%%%It trivially follows that, if each fundamental form vanishes, then in a collar neighborhood of $\Sigma \hookrightarrow (M,g)$, all geodesics are determined by geodesics on $\Sigma$. In particular, every geodesic on $M$ is given by $t \mapsto (\alpha t, \bar{\gamma}(t)) \in M$, where $\bar{\gamma}$\edz{do we need to make this arclength parametrization?} is a geodesic on $\Sigma$ and $\alpha$ is some constant. \edz{make this bidirectional}
%%%\end{remark}



%
%\color{red}
%Willmore surfaces?
%\color{black}
%
%LOOK FOR OTHER INTERESTING CLASSES OF TOTALLY GEODESIC HYPERSURFACES... zero locus of Killing vector fields? What else? Geodesic hypersurfaces in spaceforms have non-vanishing third fundamental form but vanishing second fundamental form?
%
%Redo calculation with conformally Einstein---right everything in terms of $\IIo$, $H$, $\mathring{P}$, and $J$. Also, what about Bach-flat to Bach-flat?
%
%Conformally Einstein:
%$$\nabla_{a} \nabla_{b} s + P_{ab} s + g_{ab} \rho = 0\,.$$
%Pullback to the hypersurface with $\II = \III = 0$:
%\begin{align*}
%0 \=& \bar{\nabla}_a \bar{\nabla}_b \bar{s} + P^\top_{ab} \bar{s} + \bar{\rho} \bar{g}_{ab} \\
%\=& \bar{\nabla}_a \bar{\nabla}_b \bar{s} + \tfrac{1}{d-2}\left(Ric^\top_{ab} - \frac{Sc}{2(d-1)} \bar{g}\right) \bar{s}  + \bar{\rho} \bar{g}_{ab} \\
%\=& \bar{\nabla}_a \bar{\nabla}_b \bar{s} + \tfrac{1}{d-2} \left(\bar{Ric}_{ab} - \frac{2 Ric_{nn} + \bar{Sc}}{2(d-1)} \bar{g} \right) \bar{s}   + \bar{\rho} \bar{g}_{ab}\\
%\=&\bar{\nabla}_a \bar{\nabla}_b \bar{s} + \tfrac{1}{d-2} \left((d-3) \bar{P}_{ab} + (\bar{J} - \frac{\bar{Sc}}{2(d-1)}) \bar{g} \right) \bar{s}  + \bar{\rho} \bar{g}_{ab} \\
%\=&\bar{\nabla}_a \bar{\nabla}_b \bar{s} + \tfrac{1}{d-2} \left((d-3) \bar{P}_{ab} +\tfrac{1}{2(d-1)(d-2)}  \bar{Sc} \bar{g} \right) \bar{s}  + \bar{\rho} \bar{g}_{ab} \\
%\=&\bar{\nabla}_a \bar{\nabla}_b \bar{s} + \tfrac{1}{d-2} \left((d-3) \bar{P}_{ab} +\tfrac{1}{d-1}  \bar{J} \bar{g} \right) \bar{s}  + \bar{\rho} \bar{g}_{ab} \\
%\end{align*}
%Note that because $\III = 0$, we have that $Ric_{nn} = 0$. Taking the trace-free piece, we see that the conformally-Einstein condition is not satisfied.
%
%
%When does symmetric space pullback to symmetric space?
%




\subsection{Base-like Hypersurface Embeddings} \label{base}

Warped product manifolds with a hypersurface embedding $\Sigma \hookrightarrow (M,g)$ as the base manifold have a much richer structure. In this case, there exists some function $f \in C^\infty_+ (\Sigma)$ and some coordinates $(t,x^i)$ such that
$$g = f(x^i)^2 dt^2 + \bar{g}_{ij} dx^i dx^j\,.$$
Here $\Sigma = \mathcal{Z}(t)$ is the zero locus of $t$,
coordinates $\{x^i\}$ are coordinates on $\Sigma$ parallel transported by $\partial_t$, and $\bar{g}_{ij}$ 
depends only on $x^i$.
We call these manifolds \textit{base-like hypersurface embeddings}. Such embeddings can also be characterized in terms of fundamental forms, but their behavior is much more subtle. Note this corresponds to  \eqref{warped} with  a $(d-1)$-dimensional base manifold.


First, we prove that for any base-like hypersurface embedding, the even fundamental forms vanish.
\begin{theorem} \label{vanishing-evens}
Let $\Sigma \hookrightarrow (M,g)$ be a base-like hypersurface embedding with warping function $f$. Then, $\FF{k} = 0$ for all $k = 2n$ for positive integers $n$.
\end{theorem}
\begin{proof}
As $(M,g)$ is a warped product, in a coordinate system adapted to the warped product $(t, x^i)$, the metric is expressible as
$$g = f(x^1, \ldots, x^{d-1})^2 dt^2 + \bar{g}_{ij}(x^1, \ldots, x^{d-1}) dx^i dx^j\,.$$
%Now, we may explicitly compute the second fundamental form in this metric, noting that $n = \sqrt{f} dt$:
%$$\iota^* \nabla_a n_b = \iota^* [(dt)_b \partial_a \sqrt{f}] - \sqrt{f} \iota^* \Gamma^t_{ab}\,.$$
%An explicit computation of the required Christoffel symbols show that they vanish, so we have that $\II = 0$.

Now we wish to change coordinates so that $g$ is asymptotically in normal form, i.e. so that
$$g = F^{(n+2)}(s,y^j) ds^2 + G_i^{(n+1)}(s,y^j) ds dy^i + \gamma_{ij}(s,y^1,\ldots,y^{d-1}) dy^i dy^j\,,$$
where $n \geq 0$ is even and $F^{(n+2)}(s,y^j) = 1 + \mathcal{O}(s^{n+2})$ is an even function of $s$. Similarly, we will require that $G_i^{(n+1)} = \mathcal{O}(s^{n+1})$ is an odd function of $s$.

To find such a coordinate transformation, we solve order-by-order inductively. The base case is fairly easy, with $n=0$: define $t = \frac{s}{f(y^i)}$ and $x^i = y^i$. Then, we have that
$$g = ds^2 - \frac{2s \partial_{y^i} f}{f}  dy^i ds + \gamma_{ij} dy^i dy^j \,,$$
for some $\gamma_{ij}$. Clearly, here $F^{(2)}(s,y^j) = 1$ is even and $G_i^{(1)}(s,y^j) = -\frac{2s \partial_{y^i} f}{f}$ is odd in $s$.
We now use induction to proceed. 

Suppose that $n \geq 0$ is even, that there exists $T_{n+1}(s,y^j) = \frac{s}{f(y^j)} + \cdots + Q_{n+1}(y^j) s^{n+1}$ an odd polynomial in $s$, and $X^i_{n}(s,y^j) = y^i + \cdots + P^i_{n}(y^j) s^{n}$ an even function of $s$ such that
$$g = F^{(n+2)}(s,y^j) ds^2 + G_i^{(n+1)}(s,y^j) ds dy^i + \gamma_{ij} dy^i dy^j.$$
Here we use coordinates $(s,y^j)$ related to the original coordinates $(t,x^i)$ by
$t = T_{n+1}(s,y^j)$ and $x^i = X^i_{n}(s,y^j)$, where $F^{(n+2)}$ is an even function of $s$ (with first coefficient 1)
and $G_i^{(n+1)}$ is an odd function of $s$. When $n=0$, this is the base case.

We now increment $n \rightarrow n+2$. Let $t = T_{n+1}(s,y^j) + s^{n+3} A_{n+3}(y^j)$ and $x^i = X^i_{n}(s,y^j) + s^{n+2} B^i_{n+2}(y^j)$, with $A_{n+3}$ and $B^i_{n+2}$ arbitrary functions of $\vec{y}$. We would like to show that there exist unique $A_{n+3}$ and $B^i_{n+2}$ such that 
$g_{ss} = F^{(n+4)} = 1 + \mathcal{O}(s^{n+4})$ and $g_{si} = G^{(n+3)}_i = \mathcal{O}(s^{n+3})$. 
To do so, we may compute components of the metric:
\begin{align*}
g_{si} =& G^{(n+1)}_i +  (n+2)  s^{n+1} \bar{g}_{jk} (\partial_{y^i} X_n^j) B^k_{n+2} \\
&+ s^{n+2} \left((n+3) g_{tt} A_{n+3} \partial_{y^i} T_{n+1} + \bar{g}_{jk} (\partial_s X^j_n)\partial_{y^i} B^k_{n+2} \right) \\
&+ s^{n+3} g_{tt} (\partial_s T_{n+1}) \partial_{y^i} A_{n+3} \\
&+ (n+2) s^{2n+3} \bar{g}_{jk} B^j_{n+2} \partial_{y^i} B^k_{n+2} + (n+3) s^{2n+5} g_{tt} A_{n+3}  \partial_{y^i} A_{n+3} \,.
\end{align*}
Because $\bar{g}_{jk}$ is invertible, we may fix $B^i_{n+2}$ so that the second term cancels with the leading term of $G^{(n+1)}_i$. Because $G^{(n+1)}_i$ is odd and $X^j_n$ is even (satisfying $\partial_{y^i} X^j_n = \delta_i^j + \mathcal{O}(s^2)$), the remaining higher-order terms are shunted to order $n+3$ and remain odd. Then, note that because both $T_{n+1}$ and $\partial_s X^j_n$ are odd in $s$, we see that the second line is actually odd and is of order $s^{n+3}$. Indeed, it follows that $g_{si}$ is odd and of order $s^{n+3}$. Now having fixed $B^i_{n+2}$, we may consider $g_{ss}$:
\begin{align*}
g_{ss} =& 1 + 2(n+2) s^{n+1} \bar{g}_{jk} B^j_{n+2} \partial_s X^k_n \\
&+ (F^{(n+2)} - 1)  + 2(n+3) s^{n+2} g_{tt} A_{n+3} \partial_s T_{n+1} \\
&+ (n+2)^2 s^{2n+2} \bar{g}_{jk} B^j_{n+2} B^k_{n+2} + (n+3)^2 s^{2n+4} g_{tt} A_{n+3}^2 \,.
\end{align*}
Now because $B^i_{n+2}$ is fixed and $\partial_s X^k_n$ is odd, the second term on the first line is 
even and has the factor $s^{n+2}$. The same holds for $F^{(n+2)} - 1$ and also for the second
term on the middle lide of the previous display since $T_{n+1}$ is odd. Since moreover $\partial_s T_{n+1}|_{s = 0} \neq 0$, 
we may choose $A_{n+3}$ uniquely so that the sum of first two lines of the previous display is of the form
$1 + \mathcal{O}(s^{n+4})$. 
But again, $\partial_s T_{n+1}$ is even as is $F^{(n+2)}$, so all higher-order terms remain even. Thus, the induction holds. 
Note $\gamma_{jk}$ changes in an induction step but since we did not impose any particular condition
on these functions, we keep the same notation $\gamma_{jk}$ for all induction steps.


Given such a coordinate transformation, we can now determine the parity of $\gamma_{ij}$. However, this is trivially the case: $\frac{dt}{dy^i}$ is odd, so $\frac{dt}{dy^i} \frac{dt}{dy^j}$ is even. Similarly, $\frac{dx^a}{dy^i} \frac{dx^b}{dy^j}$ is even. Therefore, we find that $\gamma_{ij}$ is an even function of $s$. Taking stock of this coordinate transform, we have that
$$g = ds^2 + (\gamma^{(0)}_{ij} + s^2 \gamma^{(2)}_{ij} + \cdots) dy^i dy^j + \mathcal{O}(s^m)$$
for any choice of $m$. Therefore, we have that in these coordinates, $n^a =  \partial_s + \mathcal{O}(s^m)$, and so $\mathcal{L}_{n} T = \partial_s T + \mathcal{O}(s^{m-1})$ for any tensor $T$. So, for any $k \leq m$, we have that $\iota^* \mathcal{L}_n^k g = \iota^* \partial_s^k g$. Notably, for any $k$ odd, we have that we may always find a coordinate transform so that
$$\iota^* \partial_s^k g = 0\,.$$

Now recall from Lemma~\ref{lie-derivs} that $\iota^* \mathcal{L}_n^k g$ is proportional to the $(k+1)$th fundamental form, modulo lower order fundamental forms. So it follows directly from  Lemma~\ref{lie-derivs} that when $k$ is odd,
$$ \FF{k+1} = \text{products of lower order fundamental forms.}$$
Here, as per Remark~\ref{weights}, each summand on the right hand side of the above display has odd weight $2-k$. But odd fundamental forms always have even weight, as does $\bar{g}^{-1}$. It thus follows that no summand on the right hand side of the above display is composed entirely of odd fundamental forms. Hence, it follows that if $\FF{m}$ vanishes for each $2 \leq m \leq k-1$ even, then it also follows that $\FF{k+1}$ vanishes. This completes the induction and hence the proof.

\end{proof}

This theorem shows that the first potentially non-vanishing fundamental form for a base-like hypersurface embedding is $\III$. It follows from the definition of $\III$ and~\cite[Proposition 42]{oneill1983} that
$$\III = f^{-1} \bar{\nabla}^2 f\,.$$
This observation is essential for the results that follow.

%\begin{proposition} \label{warped-product-isometry}
%Let $\Sigma \hookrightarrow (M,g)$ be a base-like hypersurface embedding where $(\Sigma,\bar{g})$ is a connected complete Riemannian manifold and the fiber manifold is $\Lambda$. Furthermore, let $\III = 0$. Then, $M$ is isometric to $\Lambda \times_f \left(\mathbb{R} \times \bar{\Sigma} \right)$ for some manifold $\bar{\Sigma}^{d-2}$, and there exists a choice of coordinates such that
%$$g = f^2 dt^2 + df^2 + \bar{\gamma}_{ij} dx^i dx^j\,.$$
%\end{proposition}
%\begin{proof}
%First, note that if $\III = 0$, then $\bar{\nabla}^2 f = 0$. Then, the isometry between $\Sigma$ and $\mathbb{R} \times \bar{\Sigma}$ follows directly from~\cite[Theorem 6.2]{wu2014}. The remainder of the result follows from the definition of $M$. \color{red} CHECK THIS \color{black}
%\end{proof}

%\begin{proposition} \label{warped-product-isometry}
%Let $\Sigma \hookrightarrow (M,g)$ be a base-like hypersurface embedding. Furthermore, let $\III = 0$. Then, in a small neighborhood around a point $p \in \Sigma$ there exists a choice of coordinates such that
%$$g = f^2 dt^2 + df^2 + \bar{\gamma}_{mn} dy^m dy^n\,,$$
%where $(f,y^m)$ are coordinates on $\Sigma$ in that neighborhood.
%\end{proposition}
%\begin{proof}
%First, note that if $\III = 0$, then $\bar{\nabla}^2 f = 0$. It therefore follows that both $\bar{\nabla} f$ and $\iota^* g - df \otimes df$ are parallel with respect to $\bar{\nabla}$. Therefore, for any point $p \in \Sigma$, there exists a neighborhood $\bar{U} \subset \Sigma$ around $p$ with coordinates $(f,y^m)$ such that the induced metric satisfies
%$$\iota^* g|_{\bar{U}} = df^2 + \bar{\gamma}_{mn} dy^m dy^n\,.$$
%Now, by parallel transporting these coordinates transverse to $\Sigma$, these coordinates are defined on a small neighborhood $U \subset M$ around $p$ such that the metric becomes
%$$g|_U = f^2 dt^2 + df^2 + \bar{\gamma}_{mn} dy^m dy^n\,.$$
%This completes the proof.
%\end{proof}


\begin{proposition} \label{warped-product-isometry}
Let $\Sigma \hookrightarrow (M,g)$ be a base-like hypersurface embedding. Furthermore, let $\III = 0$. Then, in a neighborhood around any point $p \in \Sigma$, $(M,g)$ is isometric to a product manifold.
\end{proposition}
\begin{proof}
If $f$ is a (nonzero) constant, the warped product is actually a product metric. If $f$ is not a constant and 
$df$ is parallel then $df$ is nonzero everywhere. Thus we can further assume  $\III = 0$, i.e.\ $\bar{\nabla}^2 f = 0$,
and $df$ is nonzero everywhere. Then both $\bar{\nabla} f$ and $\iota^* g - df \otimes df$ are parallel with respect 
to $\bar{\nabla}$. Therefore, for any point $p \in \Sigma$, there exists a neighborhood $\bar{U} \subset \Sigma$ around $p$ with coordinates $(f,y^m)$ such that the induced metric satisfies
$$\iota^* g|_{\bar{U}} = df^2 + \bar{\gamma}_{mn} dy^m dy^n\,.$$
Extending these coordinates on $\bar{U}$ to coordinates on some neighbourhood $U  \subseteq \mathbb{R} \times \bar{U}$ 
of $p$ in $M$, we obtain
$$g|_U = f^2 dt^2 + df^2 + \bar{\gamma}_{mn} dy^m dy^n\,.$$

But then the standard coordinate transformation to cartesian coordinates $f = \sqrt{u^2 + v^2}$ and $t = \arctan \tfrac{v}{u}$ yields a product metric (at least locally):
$$g = dv^2 + du^2 + \bar{\gamma}_{mn} dy^m dy^n\,,$$
and here $\Sigma = \mathcal{Z}(v)$ is the zero locus of $v$. The proposition follows.
\end{proof}

In fact, we have indentified $(M,g)$ from the previous proposition as a Riemannian product 
with a 2-dimensional Euclidean factor. But of course, it can be also considered as a product metric with a 1-dimensional
factor (as the original warped product metric).


%Note that from this proposition we easily get the following corollary.
%\begin{corollary}
%Let $\Sigma^2 \hookrightarrow (M^3,g)$ be a base-like hypersurface embedding whose warping function $f$ vanishing Hessian on $\Sigma$, a connected complete Riemannian manifold. Then, $(M,g)$ is isometric to one of the following:
%\begin{itemize}
%\item $(\mathbb{R}^3,g_{flat})$
%\item $(\mathbb{R}^2 \times S^1, g_{flat})$
%\item $(S^1 \times \mathbb{R}^2, g_{flat})$
%\item $(S^1 \times \mathbb{R} \times S^1,g_{flat})$.
%\end{itemize}
%\end{corollary}
%\begin{proof}
%From Proposition~\ref{warped-product-isometry}, it follows that $(M,g)$ is isometric to $\Lambda \times \mathbb{R} \times \bar{\Sigma}$. However, because $\Sigma$ is connected and complete, $\bar{\Sigma}$ is either topologically a circle or the real line, as is $\Lambda$. Either way, in some small neighborhood, the isometric condition implies that there exists a choice of coordinates such that
%$$g = f^2 dt^2 + df^2 + dy^2\,.$$
%However, this metric is just the flat metric in cylindrical coordinates, so the corollary follows.
%\end{proof}

%Next, it follows from Proposition~\ref{warped-product-isometry} that a vanishing third fundamental form means that a warped product is a product:
%\begin{proposition} \label{III0-prod}
%Let $\Sigma \hookrightarrow (M,g)$ be a base-like hypersurface embedding where $\Sigma$ is a connected complete Riemannian manifold and $\III = 0$. Then, $(M,g)$ is a product manifold $\Lambda \times \Sigma$ where $\Lambda$ is the fiber manifold of $(M,g)$.
%\end{proposition}
%\begin{proof}
%From Proposition~\ref{warped-product-isometry}, we may write
%$$g = f^2 dt^2 + df^2 + \bar{\gamma}_{ij} dx^i dx^j\,.$$
%But then the standard coordinate transformation to cartesian coordinates $f = \sqrt{u^2 + v^2}$ and $t = \arctan \tfrac{v}{u}$ yields a product metric (at least locally):
%$$g = dv^2 + du^2 + \bar{\gamma}_{ij} dx^i dx^j\,,$$
%and here $\Sigma = \mathcal{Z}(v)$. The proposition follows.
%\end{proof}

But from this proposition and from Corollary~\ref{prod-metric}, it follows that each fundamental form vanishes, which suggests that the odd fundamental forms must be expressible as homogeneous operators on $\III$. Indeed, we have the following proposition:
\begin{proposition} \label{diff-op-FFs}
Let $\Sigma \hookrightarrow (M,g)$ be a base-like hypersurface. Then, there exists an infinite family $\{O^f_n\}_{n=1}^{\infty}$ of differential operators on $\Gamma(\odot^2 T^* \Sigma)$ depending on $f$ such that
$$\FF{2n+3} = O^f_n (\III)\,.$$
Furthermore, this operator satisfies $O^f_n(0) = 0$.
\end{proposition}
\begin{proof}
First, note that $\FF{2n+3}$ has a local, differential formula expressible entirely in terms of $f$, $\bar{g}$, their partial derivatives, and $\bar{g}^{-1}$. Now, as such objects are tensorial, we  can instead write such objects using formulas constructed entirely from $f$, its (hypersurface) covariant derivatives, the metric $\bar{g}$ and its inverse, and the curvature tensor $\bar{R}$.

Now in the case where $f$ is constant, we trivially have from Corollary~\ref{prod-metric} and Proposition~\ref{warped-product-isometry} that $\FF{2n+3} = 0$. Therefore, any such formula for $\FF{2n+3}$ cannot contain terms that only depend on $\bar{R}$ and $f$ --- i.e. all summands in such a formula must contain one or more derivatives of $f$. Thus, we may decompose any such formula into the following pieces:
$$\FF{2n+3} = O(\bar{\nabla}^2 f) + P(\bar{\nabla} f)\,,$$
where $O$ is a differential operator depending on $f$ which is  homogeneous (with respect to rescaling $f$ by a constant)  and 
$P$ is a (potentially non-linear) homogeneous algebraic operator composed of the curvature tensor $\bar{R}$, its derivatives, the metric, its inverse, and $f$.

Now consider the case where $\III = 0$, and thus $\bar{\nabla}^2 f = 0$. Because $\FF{2n+3} = 0$ by Proposition
\ref{warped-product-isometry}, it therefore follows that
$$P(\bar{\nabla} f) = 0\,.$$
However, $\bar{\nabla} f$ does not necessarily vanish (as it is independent of $\bar{\nabla}^2 f$). But because $P$ is an algebraic operator that vanishes for \textit{any} choice of $f$ satisfying $\bar{\nabla}^2 f = 0$, we must have that $P = 0$ identically. It thus follows that $\FF{2n+3} = O(\bar{\nabla}^2 f) =: O^f_n(\III)$. As $O$ is homogeneous, it trivially follows that $O^f_n(0) = 0$.
%
%Now, we suppose for contradiction that there exists some odd fundamental form $\FF{2n+3}$ such that there exists no differential operator $O^f_n$ such that $\FF{2n+3} = O^f_n(\III)$. But $\III = f^{-1} \bar{\nabla} \bar{\nabla} f$, and so in particular any differential formula for  $\FF{2n+3}$ that contains in each term at least one instance of $\bar{\nabla}^2 f$ necessarily can be written as an operator on $\III$ that depends on $f$. So, there must exist terms in this formula for $\FF{2n+3}$ that are expressible just in terms of $\bar{R}$ (and its covariant derivatives) and at most one derivative of $f$.
%
%Now consider the case where $\III = 0$. In that case, from Propositions~\ref{prod-vanishingFFs} and~\ref{III0-prod}, it follows that $\FF{2n+3} = 0$. However, $\FF{2n+3}$ depends only  on (derivatives of) $\bar{R}$ and at most a single derivative of $f$ (in a way where at least one term cannot be written in terms of $\bar{\nabla}^2 f$). But, $f$ and $\bar{R}$ are independent, so it must be the case that $\FF{2n+3}$ vanishes identically. Hence, it follows that $\FF{2n+3}$ is expressible as a differential operator on $\III$. The second claim follows as well.
\end{proof}

\begin{remark}
As a first example, one may show, by direct computation, that for base-like hypersurface embeddings,
$$\V_{ab} =  f^{-1} (\bar{\nabla}^c f )(-3 \bar{\nabla}_c \III_{ab} + 2 \bar{\nabla}_{(a} \III_{b)c})\,.$$
Surprisingly, the operator $O^f_1$ is quite simple.
\end{remark}

From Proposition~\ref{diff-op-FFs}, we finally have the following theorem.
\begin{theorem} \label{isometric-base}
Let $\Sigma \hookrightarrow (M,g)$ be an analytic hypersurface embedding. Further, suppose that $\FF{2n} = 0$ for each $n$, that there exists a positive analytic function $f : \Sigma \hookrightarrow \mathbb{R}$ satisfying $\bar{\nabla}^2 f = f \III$, and that $\FF{2n+3} = O^f_n (\III)$ for each $n$. Then, in a neighborhood $U$ of a point in $\Sigma$, $(U,g|_U)$ is isometric to a warped product manifold with base manifold $U \cap \Sigma$.
\end{theorem}
\begin{proof}
Given the function $f$ and following Theorem~\ref{vanishing-evens} and Proposition~\ref{diff-op-FFs}, we may construct a base-like hypersurface embedding with induced metric and fundamental forms that match those given in the theorem. Hence it follows from Theorem~\ref{equal-metrics}, that, in a neighborhood of a point in $\Sigma$, we have that the given embedding restricts to a base-like hypersurface embedding.

\end{proof}

\section{Conformal Fundamental Forms and Products} \label{sec:conf}
In the same way that a Riemannian fundamental form is a natural Riemannian invariant, meaning that the quantity transforms covariantly with respect to coordinate transformations and can be constructed from natural geometric objects arising from the embedding (such as the metric, curvatures, and the unit normal vector), conformal fundamental forms are conformal invariants, meaning they transform covariantly under both coordinate transformations and conformal rescalings and are similarly constructible. Conformal fundamental forms, then, can be considered curvature corrections to the Riemannian fundamental forms such that they respect conformal rescalings. More generally, conformal fundamental forms are defined as below~\cite{BGW1}:
\begin{definition}
Let $2 \leq k \in \mathbb{N}$. A $k$th \textit{conformal fundamental form} is any natural section $\FFo{k}$ of $\odot^2_{\circ} T^* \Sigma$ with transverse order $k-1$ such that, for an embedding $\Sigma \hookrightarrow M$ with $M$ equipped with two metrics $g$ and $\Omega^2 g$, the section obey
$$\FFo{k}^{\Omega^2 g} = \Omega^{3-k} \FFo{k}^g\,.$$
\end{definition}

Observe that the trace-free second fundamental form $\IIo$ satisfies this definition, as does, for example, a third conformal fundamental form
$$\IIIo := W(n,\cdot, \cdot, n)|_{\Sigma}\,.$$
Conformal fundamental forms are examples of sections of weighted density bundles.

To elaborate on this fact, note that a conformal manifold $(M,\cc)$ can be identified with a ray subbundle $\mathcal{G} \in \odot^2 T^* M$ where the fiber at $p \in M$ is the set of all possible metrics $g \in \cc$ evaluated at $p$. Note that this subbundle is a principal bundle with structure group $\mathbb{R}_+$. A \textit{density bundle of weight} $w \in \mathbb{R}$ denoted by $\ce M[w]$, then, is a line bundle associated to $\mathcal{G}$ via the irreducible representation $\mathbb{R}_+ \ni t \mapsto t^{-w/2} \in \operatorname{End}(\mathbb{R})$. Note that this bundle can also be constructed as a real power of the volume forms
$$\ce M[w] \cong [(\wedge^d TM)^2]^{w/2d}\,,$$
and a section can be viewed as a double equivalence class $\phi = [g; f] = [\Omega^2 g; \Omega^w f] \in \Gamma(\ce M[w])$.
Then, for any vector bundle $\mathcal{V}$, we may define a weighted vector bundle $\mathcal{V}[w]$ as the tensor product $\mathcal{V}[w]:= \mathcal{V} \otimes \ce M[w]$. In that case,
$$\FFo{k} \in \Gamma(\odot^2_{\circ} T^* \Sigma[3-k])\,.$$

There are several constructions of conformal fundamental forms~\cite{BGW1,Blitz1}, and it is not known to what extent conformal fundamental forms can always be constructed. Indeed, it is known that when the conformal manifold $(M^d,\cc)$ has even dimension, conformal fundamental forms up to transverse order $d-2$ may be constructed for any conformal hypersurface embedding~\cite{Blitz1}. On the other hand, when the conformal manifold has odd dimension, conformal fundamental forms are constructible up to transverse order $\tfrac{d-1}{2}$~\cite{BGW1}. Going forward, when results involving those conformal fundamental forms $\FFo{k}$ with $k \geq 4$, we use the canonical constructions offered in~\cite{BGW1} when available or otherwise use the construction in~\cite{Blitz1}.

Several of the results in the section that follows assume that the dimension $d$ of the conformal manifold $(M^d,\cc)$ is even. However, all of these results may be reproduced in the odd dimensional case, except they only characterize or rely upon conformal fundamental forms up to transverse order $\tfrac{d-1}{2}$.

\medskip

Now we may study the relationship between conformal fundamental forms and the existence of a product metric $g \in \cc$ such that for some metric $\bar{g} \in \bar{\cc} := \iota^* \cc$, $g$ is isometric to $ds^2 + \bar{g}$, as we did for the relationship between Riemannian fundamental forms and warped product manifolds. We call such conformal hypersurface embeddings \textit{conformal product manifolds}. As we do not necessarily have access to a ``complete'' family of conformal fundamental forms (as perhaps one might hope for an infinite family while only a finite number are defined), we may also define an \textit{asymptotic conformal product manifold}, which is a conformal hypersurface embedding $\Sigma \hookrightarrow (M,\cc)$ for which there exists  $g \in \cc$ and $\bar{g} \in \bar{\cc}$ such that there exists an isometry $\Phi$
$$\Phi^* g = ds^2 +\bar{g} + \mathcal{O}(s^{d-1})\,.$$

\medskip

Having established the conformal equivalent of (warped) product manifolds, we begin with a conformally-invariant version of Theorem~\ref{vanishing-evens}.
\begin{lemma} \label{even-conf-FFs}
Let $\Sigma \hookrightarrow (M^d,\cc)$ be a conformal product manifold with $d$ even. For each integer $k \leq \tfrac{d-2}{2}$, $\FFo{2k} = 0$.
%, and when $d$ is odd, for each integer $k \leq \tfrac{d-1}{4}$, $\FFo{2k} = 0$.\edz{FIX ME!}
\end{lemma}
%\edz{CHECK IF THESE ARE DEFINED IN ODD DIMENSION}
\begin{proof}
As $\Sigma \hookrightarrow (M,\cc)$ is a conformal product manifold, by definition, there exists a representative $g \in \cc$ and $\bar{g} \in \bar{\cc}$ such that 
%$M$ is diffeomorphic to $\mathbb{R} \times \Sigma$ and 
there exists an isometry $\Phi$ such that  $\Phi^* g = ds^2 + \bar{g}$. We work in this metric representative.

Now note that all conformal fundamental forms are natural conformal hypersurface invariants, meaning they are expressible in terms of the metric $g$, the unit normal vector $n$, the covariant derivative $\nabla$, and the curvature tensor $R$~\cite{Blitz2}. Note that in any such formula, the only derivatives of $n$ that may appear are along the hypersurface, i.e. $(\nabla^\top)^m n$. Next, observe by direct computation of $\mathcal{L}_{n} g|_{\Sigma}$ that $\II = 0$. Thus, derivatives of $n$ may not contribute in a non-vanishing way to any conformal fundamental form. It follows that all natural conformal hypersurface invariants are expressible in terms of an undifferentiated $n$, the metric $g$, the curvature $R$, and its derivatives.

However, as the embedding is a product manifold, $\nabla = \partial_s + \bar{\nabla}$. Furthermore, $R^\top = \bar{R}$, $R_{abc n}^\top = 0$, $R_{nabn} = 0$, and $Ric = \bar{Ric}$, and there is no $s$ dependence in any curvature quantities. Thus, every natural conformal hypersurface invariant is expressible in terms of $\bar{g}$, $\bar{\nabla}$, and $\bar{R}$, i.e. these natural conformal hypersurface invariants are in fact expressible as natural conformal invariants of $(\Sigma, \bar{g})$. Note that $n$ cannot appear as natural conformal hypersurface invariants are tensors along $\Sigma$.

Now, observe that all pointwise-conformal invariants of $(\Sigma,\bar{g})$ (that do not involve the volume form) have even conformal weight, as $\bar{g}$ has homogeneity $2$, $\bar{\nabla}$ has homogeneity $0$, and $\bar{R}$ has homogeneity $0$ (with the natural index placement), where a Riemannian invariant $I$ with homogeneity $w$ is one that scales to $\lambda^w I$ when $g \mapsto \lambda^2 g$ for constant $\lambda$. But, $\FFo{2k} \in \Gamma(\odot^2_\circ T^* \Sigma[3-2k])$ has odd weight. As no composition of the intrinsic geometric quantities can produce an odd weight, there is only one section of this space that is built from geometric quantities: the zero section. This completes the proof.
\end{proof}

Like in the base-like hypersurface case for Riemannian warped product metrics, characterizing the odd conformal fundamental forms in the case of a conformal product manifold takes more finesse. As such, we first introduce some technical results.

The first of these technical results relates the third conformal fundamental form to the trace-free component of the Schouten tensor induced on the hypersurface.
\begin{proposition} \label{IIIo-Pb}
Let $\Sigma^{d-1} \hookrightarrow (\mathbb{R} \times \Sigma^{d-1},ds^2 + \bar{g})$. Then, $\IIIo = \tfrac{d-3}{d-2} \mathring{\bar{P}}$.
\end{proposition}
\begin{proof}
By definition, $\IIIo \= W_{nabn}$. First, observe from the relationship between the Weyl tensor and the Riemann curvature tensor, we have that
$$\otop R_{nabn} = W_{nabn} - \otop P_{ab}\,,$$
where $\otop$ is defined as the projection of a tensor the hypersurface (co)tangent bundle (or tensor products thereof) and subsequent projection to the trace-free part.
On a product manifold, however, it follows from Corollary~\ref{prod-metric} that $R_{nabn} = 0$, and so $W_{nabn} = \otop P_{ab}$. But by applying a standard hypersurface identity~\cite{Will1}
$$\IIo^2_{(ab)_{\circ}} - W_{nabn} = (d-3) \left(\otop P_{ab} - \mathring{\bar{P}}_{ab} + H \IIo_{ab} \right)\,,$$
and recalling from Proposition~\ref{even-conf-FFs} that $\IIo = 0$ for a product manifold, it follows that by direct computation that $\IIIo = \tfrac{d-3}{d-2} \mathring{\bar{P}}$. 
\end{proof}
The corollary then trivially follows:
\begin{corollary} \label{Einstein-IIIo}
Let $\Sigma \hookrightarrow (\mathbb{R} \times \Sigma,ds^2 + \bar{g})$. Then, $\IIIo = 0$ if and only if $(\Sigma,\bar{g})$ is Einstein.
\end{corollary}
%\begin{proof}
%First, suppose that $(\Sigma,\bar{g})$ is Einstein. Then, $\bar{Sc}$ is constant. Furthermore, $\mathring{\bar{Ric}} = 0$. However, because $\bar{Ric} = Ric$ for this product manifold, we may trace this equation to find that $Sc = \bar{Sc}$. It then implies that
%$$Ric_{ab} - \tfrac{1}{d-1} \bar{g}_{ab} Sc = 0\,.$$
%Expressing this equation in terms of the Schouten tensor and taking the tangential part (and recalling that $\mathring{\bar{P}} = 0$), we have that
%$$P^\top - \bar{P} = 0\,.$$
%But for a product manifold, this implies that $\IIIo = 0$. \edz{prove...}
%
%On the other hand, if $\IIIo = 0$, then $W_{nabn} = 0$. Expressing this in terms of the Schouten tensor and the Riemann tensor and recalling that $R_{nabn} = 0$ for a product manifold, we have that
%$$P^\top - \bar{g} P_{nn} = 0\,.$$ \edz{need to put standard hypersurface identities somewhere... maybe in-line in proofs, and then reference later}
%But also recall that because $W_{nabn} = 0$, we have that $P^\top - \bar{P} = 0$ for geodesic manifolds (of which a product manifold is an example). So, substituting this in to the above equation and tracing, we find that
%$$\bar{P} - \tfrac{1}{d-1} \bar{J} \bar{g} = 0\,.$$
%That is, $\mathring{\bar{P}} = 0$, so $(\Sigma,\bar{g})$ is Einstein.
%\end{proof}

Next, we are concerned with a solution to the singular Yamabe problem~\cite{Loewner,Aviles,ACF}, which is the problem of 
finding a defining density $\sigma \in \Gamma(\ce M[1])$ of $\Sigma = \mathcal{Z}(\sigma)$ such that 
$|I_{\sigma}|^2 := |d\sigma|^2 - \tfrac{2\sigma}{d}(\Delta \sigma + J \sigma) = 1$. 
(Note $\sigma$ is defining density if it is a defining function in any scale in $\cc$.)
It is known that there always exists such a density, and furthermore there always exists~\cite{Will1} a smooth solution to $|I_{\sigma}|^2 = 1 + \sigma^d \psi_d$, 
where $\psi_d \in \Gamma(\ce \Sigma[-d])$ is the celebrated Willmore invariant. 
We shall call the defining density $\sigma$ \idx{canonical} if it is such solution. Note
canonical defining density is given uniquelly modulo $\mathcal{O}(\tilde{\sigma}^{d+1})$ for any defining density 
$\tilde{\sigma}$.


\begin{theorem} \label{Yamabe-general}
Let $(\Sigma,\bar{g})$ has dimension greater than 2 and let 
$\Sigma \hookrightarrow (\mathbb{R} \times \Sigma,g:= ds^2 + \bar{g}_{ij} dx^i dx^j)$. Then, 
the canonical defining density that solves the singular Yamabe problem $\sigma$ is given an by
odd power series in $s$, i.e.
$$
\sigma = s + \sum_{\ell=1}^r \varphi_{2\ell+1} s^{2\ell+1}.
$$
More precisely, if $d$ is odd then $|I_\sigma|^2=1$ can be solved to all orders,
i.e.\ $r=\infty$, hence the Willmore invariant is zero.
If $d$ is even then the power series is deterermined only to the order $d-1$, i.e.\ for $r=\tfrac{d}{2}-1$.
Then $|I_\sigma|^2=1 + \psi_{d}s^d + O(s^{d+1})$ where $\psi_d$ is the Willmore invariant.
Moreover, all coefficients $\varphi_{2\ell+1}$  and also $\psi_{d}$
are polynomials in $J$ and its covariant derivatives.
\end{theorem}

The curvature quantity $J$ and the covariant derivative correspond to the metric $g$. But since
there is a very simple relation between covariant derivatives of $g$ and $\bar{g}$, the previous theorem
means coefficients $\varphi_{2\ell+1}$ and $\psi_d$ can be expressed 
as polynomials in $\bar{J}$ of $\bar{g}$ and its covariant derivatives with respect to  $\overline{\nabla}$.

\begin{proof}
We shall prove by induction (with respect to $\ell_0$) that we can choose
$\varphi_{2\ell+1}$, $0 \leq \ell \leq \ell_0 \leq r$ 
(as polynomials in $J$ and its covariant derivatives) in such a way that 
$|I_\sigma|^2 = 1 + \psi_{2(\ell_0+1)} s^{2(\ell_0+1)} + O(s^{2(\ell_0+1)+1})$. 
The case $\ell_0=0$---where the previous display is understood as
$\sigma=s$---can be easily  verified by a direct computation with $\psi_2 = -\tfrac{2J}{d}$.
Now assume we have 
$$
\sigma = s + \sum_{\ell=1}^{\ell_0} \varphi_{2\ell+1} s^{2\ell+1}
\quad \text{such that} \quad
|I_\sigma|^2 = 1 + \psi_{2(\ell_0+1)} s^{2(\ell_0+1)} + O(s^{2(\ell_0+1)+1})
$$
and put $\tilde{\sigma} = \sigma + \varphi_{2\ell_0+3} s^{2\ell_0+3}$ with $\varphi_{2\ell_0+3}$
to be determined. Using the notation $n_a := \nabla_a s$ (which satisfies $\nabla_a n_b =0$ and $|n|^2=1$), 
we compute
\begin{align*}
& \nabla_a \tilde{\sigma} = \nabla_a \sigma + (2\ell_0+3) \varphi_{2\ell_0+3} n_a s^{2\ell_0+2}
+ O(s^{2\ell_0+3}), \\
& \Delta \tilde{\sigma} = \Delta \sigma + (2\ell_0+2) (2\ell_0+3) \varphi_{2\ell_0+3} s^{2\ell_0+1}
+ O(s^{2\ell_0+2}).
\end{align*}
Since $\nabla_a \sigma = n_a + O(s^2)$, we further compute
\begin{align*}
& (\nabla^a \tilde{\sigma})(\nabla_a \tilde{\sigma}) = (\nabla_a \sigma)( \nabla_a \sigma) + 2(2\ell_0+3) \varphi_{2\ell_0+3} s^{2\ell_0+2} + O(s^{2\ell_0+3}), \\
& \tilde{\sigma} \Delta \tilde{\sigma} = \sigma \Delta \sigma + (2\ell_0+2) (2\ell_0+3) \varphi_{2\ell_0+3} s^{2\ell_0+2} + O(s^{2\ell_0+3}), \\
& \tilde{\sigma}^2 = \sigma^2 + O(s^{2\ell_0+4}).
\end{align*}
Summarizing, we obtain
\begin{equation} \label{Isigma}
\begin{split}
|I_{\tilde{\sigma}}|^2 &= |I_{\sigma}|^2 + \bigl[ 2(2\ell_0+3) - \tfrac{2}{d} (2\ell_0+2) (2\ell_0+3) \bigr]
 \varphi_{2\ell_0+3} s^{2\ell_0+2} + O(s^{2\ell_0+3}) = \\
& = 1 + \bigl[ \psi_{2(\ell_0+1)} + \tfrac{2(2\ell_0+3)(d-2\ell_0-2)}{d}  \varphi_{2\ell_0+3} \bigr]
s^{2\ell_0+2} + O(s^{2\ell_0+3})
\end{split}
\end{equation}
hence
$$
\varphi_{2\ell_0+3} = - \tfrac{d}{2(2\ell_0+3)(d-2\ell_0-2)} \psi_{2(\ell_0+1)}.
$$
Considering the denominator, this is always possible for $d$ odd. If $d$ is even, this is possible
for small $\ell$ up to $\ell = \tfrac{d}{2}-2$ but when $\ell = \tfrac{d}{2}-1$, the square bracket
in \eqref{Isigma} is independent on the choice of $\varphi_{2\ell_0+3}$; this bracket is the Willmore invariant
$\psi_{d}$~\cite{Will1}.

Finally note that if the square bracket in \eqref{Isigma} is zero, we have 
$|I_{\tilde{\sigma}}|^2 = 1 + O(s^{2\ell_0+3})$. However, a quick analysis of possible terms
in $|I_{\tilde{\sigma}}|^2$ reveals that this is necessarily an even power series and can only depend on $J$ and its derivatives. Hence
$|I_{\tilde{\sigma}}|^2 = 1 + O(s^{2(\ell_0+2)})$ which completes the proof.
\end{proof}

First a few coefficients of the power series for $\sigma$ can be computed by a direct combination. We obtain
\begin{align*}
\sigma =& s + \tfrac{1}{3(d-2)} Js^3 
+ \tfrac{1}{15(d-2)} \bigl[ \tfrac{1}{2(d-2)}J^2 + \tfrac{1}{d-4} \Delta J \bigr] s^5 + \\
& + \tfrac{1}{105(d-2)^2} \bigl[ \tfrac{1}{6(d-2)}J^3 - \tfrac{2(2d+3)}{3(d-4)(d-6)} J \Delta J 
- \tfrac{5d}{6(d-2)(d-6)} (\nabla^a J) (\nabla_a J) \bigr] s^7 + O(s^9).\\
\end{align*}
Similalry, we can evaluate the Willmore invariant $\psi_d$. We have $\psi_d=0$ for $d$ odd and further
$\psi_4 = -\tfrac{1}{12} \Delta J$ and 
$\psi_6 = \tfrac{1}{180} [J \Delta J + (\nabla^a J)(\nabla_a J) - \tfrac12 \Delta^2J]$.

In the special case of constant scalar curvature product manifolds, such $\sigma$ is constructible explicitly 
to all orders:
\begin{lemma} \label{Yamabe-soln}
Let $(\Sigma,\bar{g})$ be a constant scalar curvature Riemannian manifold with dimension greater than 2 and let $\Sigma \hookrightarrow (\mathbb{R} \times \Sigma,g:= ds^2 + \bar{g}_{ij} dx^i dx^j)$. Then, the canonical defining density that solves the singular Yamabe problem $\sigma$ is given by
\begin{align} \label{SY-const-sc}
\sigma := \left[g;  \begin{cases}
\sqrt{\tfrac{(d-1)(2-d)}{\bar{Sc}}} \sin\left(\sqrt{\tfrac{\bar{Sc}}{(d-1)(2-d)}} s\right) & \bar{Sc} < 0 \\
 s  & \bar{Sc}= 0 \\
 \sqrt{\tfrac{(d-1)(d-2)}{\bar{Sc}}} \sinh \left(\sqrt{\tfrac{\bar{Sc}}{(d-1)(d-2)}} s \right)  & \bar{Sc} > 0\,.
\end{cases} \right]
\end{align}
\end{lemma}
\begin{proof}
By the same arguments found in the proof of Lemma~\ref{even-conf-FFs}, we have that $(\mathbb{R} \times \Sigma,ds^2 + \bar{g})$ has $Sc = \bar{Sc}$. Now, we compute pointwise at $p \in \mathbb{R} \times \Sigma$.

Suppose first that $\bar{Sc}(x^i(p)) = 0$. Then, it trivially follows that $|ds|^2 - \tfrac{2}{d}(\Delta s + Js) = 1$, and so in the product metric representative, $\sigma = [g; s]$.

Next, suppose that $Sc < 0$. Then, suppose that $\sigma = [g; A \sin (rs)]$. The singular Yamabe equation becomes
$$A^2 r^2 \cos (rs)^2 + \frac{2A^2(r^2 - J)}{d} \sin(rs)^2 = 1\,.$$
Thus, we require that $A^2 r^2 = 1$ and that $2A^2 (r^2 - J)/d = 1$. Solving this system and substituting in $J = \tfrac{1}{2(d-1)} \bar{Sc}$, the result follows.

Finally, suppose that $Sc > 0$. Then, suppose that $\sigma = [g; A \sinh (rs)]$. The singular Yamabe equation becomes
$$A^2 r^2 \cosh (rs)^2 - \frac{2A^2(r^2 + J)}{d} \sinh(rs)^2 = 1\,.$$
Again, we require that $A^2 r^2 = 1$ and that $2A^2 (r^2 + J)/d = 1$. Solving again (for positive $A$ and $r$), the result for $\sigma$ follows.
\end{proof}

\begin{remark} \label{nonconst-Sc}
Theorem~\ref{Yamabe-general} and Lemma~\ref{Yamabe-soln} trivially imply that the canonical defining density for $\Sigma \hookrightarrow (\mathbb{R} \times \Sigma,g)$ with $(\Sigma, \bar{g})$ not having a constant scalar curvature differs only from the singular Yamabe density in Display~(\ref{SY-const-sc}) by terms involving derivatives of $\bar{Sc}$.
\end{remark}

%Trivially, then, we obtain the following corollary.
%
%\begin{corollary} \label{nonconst-Sc}
%Let $(\Sigma,\bar{g})$ be a Riemannian manifold with dimension greater than 2 and let $\Sigma \hookrightarrow (\mathbb{R} \times \Sigma, g := ds^2 + \bar{g}_{ij} dx^i dx^j)$. Then the canonical defining density that solves the singular Yamabe problem $\sigma$ differs from that of Equation~(\ref{SY-const-sc}) by terms only involving derivatives of $\bar{Sc}$. \edz{change this corollary to just an observation in the text...}
%\end{corollary}
From these technical results, we have the following technical theorem:
\begin{theorem} \label{E-PE}
Suppose that $(\Sigma, \bar{g})$ is Einstein and that $\Sigma \hookrightarrow (\mathbb{R} \times \Sigma,ds^2 + \bar{g})$. Then, $(\mathbb{R}_{\pm} \times \Sigma , g^+ := \frac{ds^2 + \bar{g}}{\sigma^2})$ is Poincar\'e--Einstein, where $\sigma$ is the singular Yamabe scale in the metric representative $ds^2 + \bar{g}$.
\end{theorem}
\begin{proof}
To prove this theorem, we rely on a result of Gover~\cite{Goal}, that if there exists a defining density $\sigma$ 
for the boundary $\Sigma$ such that
$$\nabla_{(a} \nabla_{b)\circ} \sigma + \sigma \mathring{P}_{ab} = 0\,,$$
i.e.\ $\sigma$ is an almost Einstein scale,
then the interior of the conformal manifold admits a Poincar\'e--Einstein metric representative in the scale $g^+$.
%To prove this theorem, we need to show that $\mathring{P}^{g^+} = 0$. Recall that in the metric representative $ds^2 + \bar{g}$, a conformally-invariant extension of the second fundamental form is given by
%%First, from Lemma~\ref{even-conf-FFs}, we have that $\IIo = \IVo = \cdots = 0$. Furthermore, from Proposition~\ref{Einstein-IIIo}, it follows that $\IIIo = 0$. We now need to show that all of the higher odd fundamental forms vanish as well. To do so, we will consider the canonical extension of the trace-free second fundamental form:
%$$\IIo^{\rm e}_{ab} = \nabla_{(a} \nabla_{b)\circ} \sigma + \sigma \mathring{P}_{ab}\,.$$
%Indeed, this is a conformal invariant of weight $1$, and hence in the asymptotically hyperbolic scale $g^+$, this tensor equals $\mathring{P}^{g^+}$. So it suffices to show that $\IIo^{\rm e}$ vanishes identically.

We first consider the case 
$\bar{Sc} < 0$, which is constant because $(\Sigma,\bar{g})$ is Einstein. We now calculate. First, observe that
$$\nabla \sigma = Ar \cos (rs) \, ds\,,$$
where $A$ and $r$ are as in Lemma~\ref{Yamabe-soln}. Next, we have that
$$\nabla^2 \sigma = -Ar^2 \sin (rs) ds \otimes ds\,.$$
Thus, the trace-free part of this tensor is
$$\nabla_{(a} \nabla_{b)\circ} \sigma = -Ar^2 \sin(rs) (n_a n_b - \tfrac{1}{d} g_{ab}) =  \tfrac{d-1}{d} Ar^2 \sin(rs) (\tfrac{1}{d-1} \bar{g}_{ab} - n_a n_b)$$

Now note that that
\begin{align*}
\mathring{P}_{ab} &= P_{ab} - \tfrac{1}{d} g_{ab} J \\
&= P^\top +n_a n_b P_{nn} - \tfrac{1}{d} {g}_{ab} J\\
&= \bar{P} + n_a n_b P_{nn} - \tfrac{1}{d} g_{ab} J  \\
&= \mathring{\bar{P}} + \tfrac{1}{d-1} \bar{g}_{ab} \bar{J} + n_a n_b P_{nn} - \tfrac{1}{d} g_{ab} J  \\
&=  \tfrac{1}{d-1} \bar{g}_{ab} \bar{J} + n_a n_b (J-\bar{J}) - \tfrac{1}{d} g_{ab} J \\
&= \tfrac{1}{d-1} \bar{g}_{ab} \bar{J} - \tfrac{1}{d-1} n_a n_b \bar{J} - \tfrac{d-2}{d(d-1)} g_{ab} \bar{J} \\
&= \tfrac{2}{d}\left(\tfrac{1}{d-1} \bar{g}_{ab} -  n_a n_b \right) \bar{J} \\
&= \tfrac{1}{d(d-2)}\left(\tfrac{1}{d-1} \bar{g}_{ab} -  n_a n_b \right) \bar{Sc}\,.
\end{align*}
Using that $r^2 = \frac{\bar{Sc}}{(d-1)(2-d)}$, one may easily check that
$$\nabla_{(a} \nabla_{b)\circ} \sigma + \sigma \mathring{P}_{ab} = 0\,,$$
as required.
%As this equation is conformally invariant, we may express it in the singular Yamabe scale, yielding $\mathring{P}_{ab}^{g^+} = 0$. That is, $(\mathbb{R} \times \Sigma,g^+)$ is Poincar\'e--Einstein.

When $\bar{Sc} = 0$, the identity trivially follows, as $\nabla^2 s = 0$. A similar calculation as the above follows for the case when $\bar{Sc} > 0$.
\end{proof}
As an aside, this theorem yields the following corollary:
\begin{corollary} \label{Einstein-vanishing-FFs2}
When $(\Sigma, \bar{g})$ is Einstein and $\Sigma \hookrightarrow (\mathbb{R}_{\pm} \times \Sigma, ds^2 + \bar{g})$ is a smooth embedding, $\FFo{k} = 0$ for each $k \in \{2, \ldots, d-1\}$, where $\FFo{k}$ is a canonical conformal fundamental form.
\end{corollary}
\begin{proof}
Because $(\Sigma,\bar{g})$ is Einstein, it follows that $(\mathbb{R}_{\pm} \times \Sigma,\tfrac{ds^2 + \bar{g}}{\sigma^2})$ is Poincar\'e--Einstein. But then the result follows trivially~\cite[Theorem 1.8]{BGW1}.
\end{proof}

Now, from the above technical results, we may derive the following relationship between conformal product manifolds and odd conformal fundamental forms.
\begin{proposition} \label{odd-conf-FFs}
Let $\Sigma^{d-1} \hookrightarrow (M^d := \mathbb{R} \times \Sigma, ds^2 + \bar{g})$ with $d \geq 4$ even. Then, $\IIIo = \tfrac{d-3}{d-2} \mathring{\bar{P}}$ and $\FFo{k} = \mathcal{O}_{k}(\IIIo)$ for each $k \in \{4 \ldots,d-1\}$, where $\mathcal{O}_{k}$ is some map on the jets of $\IIIo$ satisfying $\mathcal{O}_k(0) = 0$. Furthermore, $\IIo = 0 = \mathcal{O}_{2m}(\IIIo)$ for each $2 \leq m \leq \tfrac{d-2}{2}$.
\end{proposition}

\begin{proof}
First, observe from Proposition~\ref{IIIo-Pb}, it follows that $\IIIo = \tfrac{d-3}{d-2} \mathring{\bar{P}}$. So it suffices to show that all conformal fundamental forms are expressible in terms of a (potentially differential) operator on $\IIIo$ or equivalently $\mathring{\bar{P}}$.

Now, consider the case where $d = 4$. Then, from Lemma~\ref{even-conf-FFs} the proposition follows.

On the other hand, when $d \geq 6$, there  exists a conformally-invariant extension of the third conformal fundamental form~\cite{Blitz1}:
$$\IIIo^{\rm e}_{ab} := W_{\hat{n} ab \hat{n}} + 2 \sigma C_{\hat{n} (ab)} - \tfrac{\sigma^2}{d-4} B_{ab}\,,$$
where $\sigma$ is the singular Yamabe defining function described by Theorem~\ref{Yamabe-general} and $\hat{n} := d \sigma$. Now as the jets of $\IIIo^{\rm e}$ entirely capture the conformal fundamental forms from $\IIIo$ up to the $(d-1)$th conformal fundamental form~\cite{Blitz1}, to prove the proposition it is sufficient to show that $\IIIo^{\rm e}$ is expressible entirely in terms of $\IIIo$ and $\mathring{\bar{P}}$.

As $M$ is a product manifold, in the coordinates specified, $\nabla_n \mathcal{C} = 0$ for every tensor $C$ constructed from the metric and its derivatives. Furthermore, as $\sigma = \Omega s$, where
$$\Omega := \sum_{m = 0}^{\infty} \bar{\Omega}_{(m)} s^m\,,$$
with $\bar{\Omega}_{(m)}$ a coefficient that depends only on the induced scalar curvature and its derivatives (and $\bar{\Omega}_{(0)} = 1$), then a simple calculation yields
$$\nabla_n \sigma = 1 + \sum_{m=1}^{\infty} (m+1) \bar{\Omega}_{(m)} s^m$$
and
$$\nabla_n \hat{n}^a =  \sum_{m=1}^{\infty} \left[ (m+1) (\bar{\nabla}^a \bar{\Omega}_{(m)}) s + m(m+1) n^a \bar{\Omega}_{(m)} \right] s^{m-1} \,.$$
A consequence of this calculation is that, for $\ell \geq 1$,
$$\nabla_n^{\ell} \sigma \= \ell! \, \bar{\Omega}_{(\ell -1)}\,,$$
and
$$\nabla_n^{\ell} \hat{n}^a \= \ell !\, \bar{\nabla}^a \bar{\Omega}_{(\ell-1)} + (\ell + 1)! \, \bar{\Omega}_{(\ell)} n^a\,.$$

Now we may use these identities to evaluate the leading structure of the $(2k+3)$rd odd conformal fundamental form:
\begin{align*}
\otop \nabla_n^{2k} \IIIo^{\rm e}_{ab} &\=  \sum_{\ell +m = 2k} \Big [(\nabla_n^{\ell} \hat{n}^c)(\nabla_n^m \hat{n}^d ) W_{cabd} + 2 (\nabla_n^{\ell} \sigma)(\nabla_n^m \hat{n}^c) C_{c(ab)} - \tfrac{1}{d-4} (\nabla_n^{\ell} \sigma)(\nabla_n^m \sigma) B_{ab} \Big]^{\otop} \\
&\= \sum_{\ell +m = 2k} \Big [ \ell! m! (\bar{\nabla}^c \bar{\Omega}_{(\ell-1)})(\bar{\nabla}^d \bar{\Omega}_{(m-1)}) W_{cabd} + (\ell+1)!(m+1)! \bar{\Omega}_{(\ell)} \bar{\Omega}_{(m)} \IIIo_{ab} \\
 &\phantom{\= \sum_{\ell +m = 2k} \Big [}   + 2 \ell! m!\bar{\Omega}_{(\ell-1)} (\bar{\nabla}^c \bar{\Omega}_{(m-1)} ) \bar{C}_{c(ab)} + 2 \ell! (m+1)! \bar{\Omega}_{(\ell-1)} \bar{\Omega}_{(m)} C_{n(ab)}  \\
&\phantom{\= \sum_{\ell +m = 2k} \Big [} - \tfrac{1}{d-4} \ell! m! \bar{\Omega}_{(\ell-1)} \bar{\Omega}_{(m-1)} B_{ab} \Big]^{\otop}
\end{align*}
Because $\bar{\Omega}_{(m)}$ depends only on the induced scalar curvature and its derivatives (see the discussion after Theorem~\ref{Yamabe-general}) and
$$\bar{\nabla}_a \bar{Sc} = 2(d-1) \bar{\nabla}^b \mathring{\bar{P}}_{ab}\,,$$
it follows that all of the above terms containing $\bar{\nabla} \bar{\Omega}_{(m)}$ are expressible in terms of an operator on $\mathring{\bar{P}}$. Further recall that on product manifolds we have $\IIo = \IVo = 0$ according to 
Lemma~\ref{even-conf-FFs} and $H := \tfrac{1}{d-1} \II_a^a =  0$ according to Corrollary~\ref{prod-metric}. 
Thus it follows from the Codazzi-Mainardi equation \cite[(2.10)]{Will1} that $W_{abc\hat{n}}^\top=0$ hence the explicit
formula \cite[(1.3)]{BGW1} for $\IVo$ shows that $C_{n(ab)}^\top \= 0$. Finally, from an expression for the fifth conformal fundamental form in~\cite{BGW1} below the display (1.3), it follows that $B_{ab}^{\otop}$ is expressible in terms of the third fundamental form and the hypersurface Bach tensor $\bar{B}_{ab}$. But the hypersurface Bach tensor is expressible in terms of $\mathring{\bar{P}}$.

The above discussion shows that the leading transverse-order term in an odd-order conformal fundamental is expressible in terms of an operator on $\mathring{\bar{P}}$. By induction, it follows that the entirety of such an odd-order conformal fundamental form is expressible the same way, as the subleading terms are expressible in terms of lower-order conformal fundamental forms which are either even-order (which vanish) or odd-order, which are expressible in the same way.

\end{proof}

%\begin{proof}
%From Proposition~\ref{IIIo-Pb}, we have that $\IIIo = \tfrac{d-3}{d-2} \mathring{\bar{P}}$.
%
%Now consider the asymptotically hyperbolic manifold $(\mathbb{R} \times \Sigma, g^+ := \frac{ds^2 + \bar{g}}{\sigma^2})$, where $\sigma$ is the canonical defining density that solves the singular Yamabe problem as in Corollary~\ref{nonconst-Sc}. Now, by construction~\cite{}, the conformal fundamental forms are expressible in terms of conformally-invariant normal derivative operators acting on the almost Einstein tensor
%$$\IIo^{\rm e}_{ab} := [g; \nabla_{(a} \nabla_{b)\circ} \sigma + \sigma \mathring{P}_{ab}]\,.$$
%Importantly, in the asymptotically hyperbolic scale, we have that $\IIo^{\rm e}_{ab} = [g^+; \mathring{P}^{g^+}]$, so put another way, the conformal fundamental forms are expressible as the action of normal derivative operators on $\mathring{P}^{g^+}$, picking up Taylor coefficients in $s$. So it suffices to show that there exists a formula for $\mathring{P}^{g^+}$ that can be expressed solely in terms of $\IIIo$.
%
%For ease of calculation, we translate to the Ricci curvature: $\mathring{Ric} = (d-2) \mathring{P}$. Now, the conformal rescaling of the Ricci tensor is given by
%$$\mathring{Ric}^{g^+}_{ab} = \mathring{Ric}^g_{ab} + (d-2) \tfrac{\nabla_a \nabla_b \sigma}{\sigma} - \tfrac{d-2}{d} \tfrac{\Delta \sigma}{\sigma} g_{ab}\,.$$
%Because the Riemannian manifold is a product manifold, we have that
%$$\mathring{Ric}_{ab}^g = \mathring{\overline{Ric}}_{ab}^{\bar{g}} + \left(\tfrac{1}{d(d-1)} \bar{g}_{ab} - \tfrac{1}{d} n_a n_b \right) \bar{Sc}\,.$$
%
%To proceed, we must compute the relevant derivative of $\sigma$. However, we only have an explicit formula for $\sigma$ in the case of constant scalar curvatures. Nonetheless, a series expansion of $\sigma$ in this setting differs from that of the constant scalar curvature case by terms that depend only on derivatives of the scalar curvature. Thus, we have that
%\begin{align*}
%\nabla_a \nabla_b \sigma =& -B \sin(Bs) n_a n_b - 2s \sin(Bs) n_{(a} \bar{\nabla}_{b)} B + \left(\frac{Bs \cos(Bs) - \sin(Bs)}{B^2} \right) \bar{\nabla}_a \bar{\nabla}_b B \\
%&- \frac{(s^2 B^2 - 2) \sin(Bs) + 2sB \cos(Bs)}{B^3} (\bar{\nabla}_a B)(\bar{\nabla}_b B) + \mathcal{O}(\bar{\nabla} \bar{Sc})\,.
%\end{align*}
%where 
%$\mathcal{O}(\bar{\nabla} \bar{Sc})$ are terms that depend explicitly on at least one derivative of $\bar{Sc}$ and
%$$B :=  \sqrt{\frac{\bar{Sc}}{(d-1)(2-d)}}\,.$$
%
%Because $\bar{\nabla}_a \bar{Sc} = 2(d-1) \bar{\nabla}^b \mathring{\bar{P}}_{ab}$, we may express any derivatives of $B$ or $\bar{Sc}$ in terms of an operator on $\mathring{\bar{P}}$. So it suffices to check that all terms that do not depend on $\bar{\nabla} B$, $\bar{\nabla} \bar{Sc}$, or $\mathring{\overline{Ric}}$ cancel in the expression for $\mathring{Ric}^{g^+}$. Indeed, those terms are
%$$\tfrac{1}{d(d-1)} \bar{g}_{ab} \bar{Sc} - \tfrac{1}{d} n_a n_b \bar{Sc} + (d-2)(-B^2 n_a n_b) - \tfrac{d-2}{d} (-B^2) g_{ab}\,.$$
%Using that $B^2 = -\tfrac{\bar{Sc}}{(d-1)(d-2)}$, we can check that this vanishes. It thus follows that $\mathring{Ric}^{g^+}_{ab}$ can be expressed as an operator on $\mathring{\bar{P}}$, and hence so can the conformal fundamental forms. From Proposition~\ref{even-conf-FFs}, we also see that the even operators $\mathcal{O}_{2n}$ are the zero operator and that $\IIo = 0$.
%
%As $\mathring{\bar{P}} = \tfrac{d-2}{d-3} \IIIo$, it thus follows that each conformal fundamental form is in fact expressible as an operator on the jets of $\IIIo$, so we denote $\FF{k} := \mathcal{O}_k (\IIIo)$. Furthermore, when $\IIIo = 0$, we have that $(\Sigma,\bar{g})$ is Einstein from Proposition~(\ref{Einstein-IIIo}). Then, from Theorem~(\ref{E-PE}), it follows that the asymptotically hyperbolic manifold is Poincar\'e--Einstein. But then the conformal fundamental forms vanish, again by~\cite{}. Thus, we have that $\mathcal{O}_k(0) = 0$. This completes the proof.
%\end{proof}

\medskip

Given Lemma~\ref{even-conf-FFs} and Proposition~\ref{odd-conf-FFs}, it is necessary that the even conformal fundamental forms vanish and the odd conformal fundamental forms take on particular forms in order for a conformal hypersurface 
 to admit a product metric in its conformal class. However, this is not sufficient. To prove this statement, we first need an identification theorem between conformal manifolds with matching conformal fundamental forms.



\begin{theorem} \label{conformal-identification}
Let $\iota : \Sigma \hookrightarrow (M^d,\cc)$ and $\hat{\iota} : \Sigma \hookrightarrow (M^d,\hat{\cc})$ with $d$ even and $\iota(p) = \hat{\iota}(p)$ for all $p \in \Sigma$. Furthermore, let $\sigma$ be a singular Yamabe density for $\iota$. Then, $\iota^* \cc = \hat{\iota}^* \hat{\cc}$ and $\FFo{k} = \hat{\FFo{k}}$ for each $k \leq d-1$ if and only if  there exists a diffeomorphism $\phi$ fixing $\Sigma$ pointwise such that
$$\phi^* \hat{\cc} = \cc + \mathcal{O}(\sigma^{d-1})\,.$$
\end{theorem}
\begin{proof}
To prove the forward direction, it suffices to show that, for a given choice of $\bar{g} \in \iota^* \cc$, there exists metrics $g \in \cc$, $\hat{g} \in \hat{\cc}$ and an isometry $\phi$ such that $\iota^* g = \hat{\iota}^* \hat{g} = \bar{g}$ and $g - \phi^* \hat{g} = \mathcal{O}(s^{d-1})$ where $\sigma = [g;s]$, i.e.\ $s$ is the representative of $\sigma$ in the scale $g$. 
So, picking $\bar{g}$, we let $g \in \cc$, $\hat{g} \in \hat{\cc}$ be the unique Graham--Lee metrics~\cite{GrahamLee} in the conformal classes, where for $\hat{\sigma} = [\hat{g}, \hat{s}]$ we have that $|ds|_g^2 = 1 =|d\hat{s}|_{\hat{g}}^2$. That is, we pick those metric representatives $g$ and $\hat{g}$ such that the defining functions $s$ and $\hat{s}$ determined by the singular Yamabe densities $\sigma$ and $\hat{\sigma} = [\hat{g}; \hat{s}]$ (which is the singular Yamabe density of $\hat{\iota}$) are the minimal geodesic distance functions between a point $p$ and $\Sigma$ as measured by the metric representatives $g$ and $\hat{g}$ respectively.
%Note that these defining functions provide a foliation of a collar neighborhood around $\Sigma$, with leaves given by $\Sigma_t = \{p \in M | s(p) = t\} = \{p \in M | \hat{s}(p) = t\}$. Denoting coordinates on $\Sigma$ by 

As $g$ and $\hat{g}$ are Graham--Lee metrics, it follows that there always exists isometries 
$\Phi,\hat{\Phi} : \Sigma \times I \rightarrow M$ such that, in a collar neighborhood of $\Sigma$,
\begin{align*}
\Phi^* g &= dt^2 + (\bar{g}_{ij}(\vec{x}) + t g^{(1)}_{ij}(\vec{x}) + \tfrac{1}{2} t^2 g^{(2)}_{ij}(\vec{x}) + \cdots) dx^i dx^j  \\
\hat{\Phi}^* \hat{g} &= dt^2 + (\bar{g}_{ij}(\vec{x}) + t \hat{g}^{(1)}_{ij}(\vec{x}) + \tfrac{1}{2} t^2 \hat{g}^{(2)}_{ij}(\vec{x}) + \cdots) dx^i dx^j \,.
\end{align*}
Here $I$ is an open interval containing $0$, i.e.\ $ \Sigma \times I \rightarrow M$ is a two-sided $\Sigma$ in $M$.
Note that $\Phi^*(s) = \hat{\Phi}^*(\hat{s}) = t$. Thus, the existence of the desired isometry $\phi$ is equivalent to $g^{(k)}(\vec{x}) = \hat{g}^{(k)}(\vec{x})$ for all $k \leq d-2$.

%While we may use the same coordinates on $\Sigma$, on the foliations given by $s$ and $\hat{s}$ respectively, it does not follow that the parallel transport of those coordinates land on the same point. Thus, we will coordinatize by $(s,x^i)$ and $(\hat{s},\hat{x}^i)$, noting that for a point $p \in \Sigma$, $\hat{x}^i(p) = x^i(p)$. Now, we may expand these metrics in $s$ and $\hat{s}$, obtaining
%\begin{align*}
%g &= ds^2 + (\bar{g}_{ij} + s g^{(1)}_{ij} + \tfrac{1}{2} s^2 g^{(2)}_{ij} + \cdots) dx^i dx^j  \\
%\hat{g} &= d\hat{s}^2 + (\hat{\bar{g}}_{ij} + \hat{s} \hat{g}^{(1)}_{ij} + \tfrac{1}{2} \hat{s}^2 \hat{g}^{(2)}_{ij} + \cdots) d\hat{x}^i d\hat{x}^j \,.
%\end{align*}
Observe that each coefficient $g^{(k)}$ is a natural hypersurface invariant (that is tangent to the hypersurface), as it can be obtained via $\mathcal{L}_n^{k} \Phi^* g|_{t = 0}$, where $n =\partial_t$, and similarly for $\hat{g}^{(k)}$. As such, from~\cite[Theorem 2.5]{Blitz2} it follows that $g^{(k)}$ is expressible as a partial contraction polynomial in
$$\{\bar{g}, \bar{g}^{-1}, \bar{\nabla}, \bar{R}, \IIo, \,\ldots\, \FFo{k+1}, H, J|_{\Sigma}, \ldots, \nabla_n^{k-2} J|_{\Sigma}\}\,,$$
and the expression is universal, meaning that the symbolic form of the expression does not depend on the underlying choice of metric.

Note that $|n|^2 = 1$, and so from the singular Yamabe equation $|n|^2 - \tfrac{2}{d} t(\Delta t + Jt) = 1 + \mathcal{O}(t^d)$, we have that $\Delta t + Jt = 0$ in a collar neighborhood of $\Sigma$. So it follows that $\nabla_n^k (\Delta t + Jt) = 0$ for every $1 \leq k \leq d-1$. Expanding this equation, we have
\begin{align*}
0 \=& \nabla_n^k \Delta t + (k-1) \nabla_n^{k-1} J \\
\=& \nabla_n^{k-1} [\nabla_n, \nabla_a] \nabla^a t + (k-1) \nabla_n^{k-1} J \\
\=& -\nabla_n^{k-1} Ric_{nn} + (k-1) \nabla_n^{k-1} J + \text{ltots} \\
\=& -(d-1) \nabla_n^{k-1} J + (k-1) \nabla_n^{k-1} J + \text{ltots} \\
\=& -(d-k) \nabla_n^{k-1} J + \text{ltots.}
\end{align*}
In the above, $\text{ltots}$ stands for ``lower transverse order terms,'' which are those terms which involve fewer normal derivatives of the metric. Note that the fourth identity follows from~\cite[Lemma 2.4]{Blitz1}. So, it is clear that for Graham-Lee metrics, $\nabla_n^{k-1} J|_{\Sigma}$ is expressible in terms of lower transverse order terms so long as $k < d$. But because $H = 0$ (because $H \= \tfrac{1}{d} \Delta t = -\rho|_{\Sigma}$), it follows by induction that \textit{in a Graham--Lee metric representative}, the coefficients $ g^{(k)}$ have universal expressions in terms of the set
$$\{\bar{g}, \bar{g}^{-1}, \bar{\nabla}, \bar{R}, \IIo, \,\ldots\, \FFo{k+1}\}\,.$$
As these expressions are universal (for Graham--Lee metric representatives), it follows that if the intrinsic geometry and the conformal fundamental forms agree for $\iota$ and $\hat{\iota}$, then viewed as tensors on $\Sigma$, $g^{(k)} =  \hat{g}^{(k)}$ for all $k \leq d-2$ (as the conformal fundamental forms are only defined up to that order).

\medskip

The reverse direction follows by noting that if any pair of conformal fundamental forms $\FFo{k}$, $\hat{\FFo{k}}$ are not equal, then it is clear from the Graham--Lee normal form of the metrics that they are not isometric.


%Now, recall that for $k \leq d-1$, any conformal fundamental form $\FFo{k}$ may be expressed as
%$$\FFo{k}_{ab} = \delta^{k-2}(\nabla_{(a} \nabla_{b)_{\circ}} s + s \mathring{P}_{ab})\,,$$
%for some conformally-invariant operator $\delta^{k-2}$ with leading normal derivative term $\otop \nabla_n^{k-2}$. However, observe that $\nabla_{(a} \nabla_{b)_\circ} s = \tfrac{1}{2} (\mathcal{L}_n g)_{(ab)_\circ}$, where $n = g^{-1}(ds,\cdot)$. So, we may alternatively express
%$$\FFo{k}_{ab} = \delta^{k-2}(\tfrac{1}{2} (\mathcal{L}_n g)_{(ab)_{\circ}} + s \mathring{P}_{ab})\,.$$
%Because $\nabla_n$ agrees with $\mathcal{L}_n$ at leading order as a differential operator, it is clear that $\mathcal{L}_n^{k-1} 
\end{proof}

Using this theorem, we may now prove the following:
\begin{corollary} \label{classification-conformal-product}
Let $\Sigma \hookrightarrow (M^d,\cc)$ be a conformal hypersurface embedding with $d$ even
satisfying $\IIo_{ab}=0$. 
Suppose that there exists some $\tau \in \Gamma(\ce_+ \Sigma[1])$ such that
$$(\bar{\nabla}_{(a} \bar{\nabla}_{b)_{\circ}} + \mathring{\bar{P}}_{ab} - \tfrac{d-2}{d-3} \IIIo_{ab}) \tau = 0\,,$$
and that in the $\tau$ scale, $\FFo{k} = \mathcal{O}_k(\IIIo)$ with $3 \leq k \leq d-1$. Then $\Sigma \hookrightarrow (M,\cc)$ is an asymptotic conformal product manifold.
\end{corollary}
\begin{proof}
From Theorem~\ref{conformal-identification}, it suffices to show that the conformal hypersurface embedding $\iota: \Sigma \hookrightarrow (M,\cc)$ has conformal fundamental forms that agree with some asymptotic conformal product manifold. Now given a choice of representative $g_{\tau}$ determined by any extension of $\tau$ to $(M,\cc)$, we may identify an isometry $\Phi : \Sigma \times [0,\epsilon) \rightarrow M$ such that $\Phi^* g_{\tau}$ is in geodesic normal form. Now on the same collar neighborhood, we prescribe another metric $\hat{g} := dt^2 + \bar{g}_{ij}(\vec{x})dx^i dx^j + \mathcal{O}(t^{d-1})$, where $\Sigma = \mathcal{Z}(t)$ and $\bar{g} = \Phi^* \iota^* g_{\tau}$ with $\Phi^*$ acting in the natural way on $\Sigma$. This metric in turn defines a conformal class of metrics $\hat{\cc}$ in a collar neighborhood of $\Sigma$, with conformal embedding $\hat{\iota} : \Sigma \hookrightarrow (M,\hat{\cc})$, with $\hat{\iota}^* \hat{\cc} = \iota^* \cc$ and with conformal fundamental forms given by Proposition~\ref{odd-conf-FFs}: $\IIo = 0$, $\IIIo = \tfrac{d-3}{d-2} \mathring{\bar{P}}$, and $\FFo{k} = \mathcal{O}_k(\IIIo)$.

But in the scale $g_{\tau}$, we have that $\tau = [g_{\tau}; 1]$ and so by hypothesis, $\mathring{\bar{P}} = \tfrac{d-2}{d-3} \IIIo$ and $\FFo{k} = \mathcal{O}_k(\IIIo)$. Furthermore, by hypothesis $\IIo = 0$. As such, the conformal fundamental forms match the conformal fundamental forms of $\hat{\iota}: \Sigma \hookrightarrow (M,\cc)$. The corollary follows.

\end{proof}

\begin{remark}
The equation
$$(\bar{\nabla}_{(a} \bar{\nabla}_{b)_{\circ}} + \mathring{\bar{P}}_{ab} - \tfrac{d-2}{d-3} \IIIo_{ab}) \tau = 0\,,$$
in the statement of Corollary~\ref{classification-conformal-product} is a conformally invariant Jacobi-like differential equation for umbilic hypersurface embeddings $\Sigma \hookrightarrow (M^d,g)$. Solutions to this equation exist if and only if there exists an infinitesimal family of umbilic hypersurfaces generated by the vector field $\tau n$, where $n$ is the unit normal vector to $\Sigma$. This follows by demanding that the first variation of the trace-free second fundamental form with respect to the embedding vanishes.
\end{remark}



%
%\section{Hypersurface Willmore Variations}
%In this section, we produce Jacobi equations for specific classes of hypersurface embeddings in Riemannian manifolds. We use the method outlined in~\cite{https://arxiv.org/pdf/1904.08423}. Summarizing the results, we have that the following:
%
%\begin{proposition} \label{minimal-Jacobi}
%Let $\Sigma \hookrightarrow (M,g)$ be a minimal hypersurface embedding. If there exists a solution to the differential equation
%$$\bar{\Delta} \mu + (\tfrac{1}{2} K - (d-2) \bar{J} + (d-1) J) \mu = 0\,,$$
%then to first order there exists a one-parameter family of minimal hypersurfaces generated by the vector field $\mu \hat{n}$.
%\end{proposition}
%Note that this equation in general has solutions as it is simply a Helmholtz equation.
%
%\begin{proposition}
%Let $\Sigma \hookrightarrow (M^d,g)$ be an umbilic hypersurface embedding. There exists a solution to the differential equation
%$$\mathcal{J}^{(\IIo)} \mu := (\bar{\nabla}_{(a} \bar{\nabla}_{b)\circ} + \mathring{\bar{P}}_{ab} + (d-2) \Fo_{ab}) \mu = 0\,,$$
%if and only if there exists an infinitesimal family of umbilic hypersurfaces generated by the vector field $\mu \hat{n}$.
%\end{proposition}
%Note that, unlike the previous case, this system is highly-overdetermined, and hence typically there are not solutions to this equation. However, if $\Fo = 0$ and $(\Sigma,\bar{g})$ is Einstein, then the equation is trivially satisfied for the Einstein scale.
%
%\begin{theorem}
%Let $\Sigma^2 \hookrightarrow (M^3,\cc)$ be a Willmore surface embedding. Then, solutions $\mu \in \Gamma(\ce \Sigma[1])$ to
%$$(\bar{P}_4^{ext} + 4 (\hat{\bar{D}}_A K) \hat{\bar{D}}^A+ \IIo^{ab} E_{ab} - 2 \IIo \cdot \IVo  + 2 K^2) \mu = 0$$
%generate an infinitesimal family of Willmore surface embeddings generated by the vector field $\mu \hat{n} \in \Gamma(T M|_{\Sigma})$.
%\end{theorem}
%In the above theorem, $\bar{P}_4^{ext}$ is the conformally-invariant Paneitz operator on surfaces acting on weight $1$ densities from~\cite{}. Specifically, it is defined according to the map
%\begin{align*}
%f\mapsto
%\bar{\Delta}^2 f &+ 4 \bar{\nabla}^a \circ \mathcal{P}_{ab} \circ \bar{\nabla}^b f \\
%&+ \left[ 2 \bar{\nabla} \csdot \bar{\nabla} \csdot \mathcal{P} - \bar{\Delta} \mathcal{J} + 2 \mathcal{P}^2 - \mathcal{J}^2 + 2 \bar{\nabla}_a \left(\IIo^{ab} \bar{\nabla} \csdot \IIo_b \right) + 2 (\bar{\nabla} \csdot \IIo)^2  \right] f
%\end{align*}
%with $\mathcal{J} := \mathcal{P}_a^a$ where
%\begin{align*}
%\mathcal{P}_{ab} := P^\top_{ab} + H \IIo_{ab} + \tfrac{1}{2} H^2 \bar{g}_{ab} - \tfrac{1}{2} K \bar{g}_{ab}\,.
%\end{align*}
%Furthermore, the scalar $\IIo \cdot E$ is defined according to
%$$\IIo^{ab} E_{ab} := \IIo^{ab} \bar{\Delta} \IIo_{ab} - 2 \IIo^{ab} \bar{\nabla}_a \bar{\nabla} \cdot \IIo_b - K \mathcal{P}_a^a - \bar{J} K - \tfrac{1}{2} K^2\,,$$
%and for surfaces embedded in three dimensional conformal manifolds, the conformal fourth fundamental form $\IVo$ is defined by
%$$\IVo_{ab} := C_{n(ab)}^{\top}\,.$$
%
%This theorem is the first example of a more general result:
%\begin{theorem}
%Let $\Sigma^n \hookrightarrow (M^d,\cc)$ with $n \geq 4$ even be a conformal hypersurface embedding. There exists an operator
%$$J^{(n)} : \Gamma(\ce \Sigma[1]) \rightarrow \Gamma(\ce \Sigma[-d])$$
%with leading derivative term expressible in terms of $\bar{P}_{n+2}^{ext}$
%%and with an algebraic term depending on $\IIo \cdot \FF{d+1}$
%. Furthermore, when $\Sigma \hookrightarrow (M^d,\cc)$ is Willmore, then $\mu \hat{n}$ generates an infinitesimal family of Willmore hypersurfaces, where $\mu$ is a solution to $J^{(n)} \mu = 0$.
%\end{theorem}
%
%\begin{proof}
%Consider the Willmore invariant $\mathcal{B}_n$ for hypersurfaces $\Sigma^n \hookrightarrow (M^d,\cc)$ of~\cite{}, with leading term term given by $\bar{\Delta}^{n/2} H$. We will define the operator $J^{(n)}$ as the hypersurface variation of the Willmore invariant when restricted to Willmore hypersurfaces. To construct the leading order term just requires observing that $\delta \bar{\Delta}$ is an operator with derivative order, and so the leading hypersurface derivative term arises from a term of the form $\bar{\Delta}^{n/2} \delta H$. Using Proposition~\ref{minimal-Jacobi}, we find that the leading derivative term in $\delta \mathcal{B}_n$ is $\bar{\Delta}^{n/2 + 1} \mu$.
%
%Recall from~\cite{} that a normal hypersurface variation of a conformal hypersurface invariant of weight $w \neq 0$ is not itself conformally-invariant. Rather, defining $\hat{\delta}_{\mu} := \delta - w \mu$, where the variation field is given by $\mu \hat{n}$, we see that $\hat{\delta}$ maps conformal hypersurface invariants to conformal hypersurface invariants. So, we define the operator
%$$J^{(n)} \mu := \hat{\delta}_{\mu} \mathcal{B}_n\,.$$
%Clearly, this operator has leading derivative term $\bar{\Delta}^{n/2 + 1} \mu$. Furthermore, the operator is itself conformally invariant, by definition. By simple weight-counting, we see that the weights of the domain and range section spaces are $1$ and $-d$ respectively. 
%
%\end{proof}
%
%\section{Product Manifolds}
%\begin{lemma} \label{even-conf-FFs2}
%Let $\Sigma \hookrightarrow (\mathbb{R} \times \Sigma, ds^2 + \bar{g})$. Then $\FFo{2k} = 0$ for each $k \in \mathbb{Z}$.
%\end{lemma}
%\begin{proof}
%First note that all conformal fundamental forms are natural conformal hypersurface invariants, meaning they are expressible in terms of the metric $g$, the unit normal vector $\hat{n}$, the covariant derivative $\nabla$, and the curvature tensor $R$~\cite{}. Note that in any such formula, the only derivatives of $\hat{n}$ that may appear are along the hypersurface, i.e. $(\nabla^\top)^m \hat{n}$. Next, observe by direct computation of $\mathcal{L}_{\hat{n}} g|_{\Sigma}$ that $\II = 0$, as do all of its derivatives. Thus, derivatives of $\hat{n}$ may not contribute in a non-vanishing way to $\FF{2k}$. Thus, all natural conformal hypersurface invariants are expressible in terms of an undifferentiated $\hat{n}$, the metric $g$, the curvature $R$, and its derivatives.
%
%However, as the embedding is a product manifold, $\nabla = \partial_s + \bar{\nabla}$. Furthermore, $R^\top = \bar{R}$, $R_{abc \hat{n}}^\top = 0$, $R_{nabn} = 0$, and $Ric = \bar{Ric}$, and there is no $s$ dependence on any curvature quantities. Thus, every natural conformal hypersurface invariant is expressible in terms of $\bar{g}$, $\bar{\nabla}$, and $\bar{R}$, i.e. these natural conformal hypersurface invariants are in fact expressible as natural conformal invariants of $(\Sigma, \bar{g})$. Note that $\hat{n}$ cannot appear is natural conformal hypersurface invariants are tensors along $\Sigma$.
%
%Now, observe that all (even) pointwise-conformal invariants of $(\Sigma,\bar{g})$ have even weight, as $\bar{g}$ has homogeneity $2$, $\bar{\nabla}$ has homogeneity $0$, and $\bar{R}$ has homogeneity $0$ (with the natural index placement). But, $\FF{2k} \in \Gamma(\odot^2_\circ T^* \Sigma[3-2k])$ has odd weight. As no composition of the intrinsic geometric quantities can produce an odd weight, there is only one section of this space that is built from geometric quantities: the zero section. This completes the proof.
%\end{proof}
%
%\begin{proposition} \label{Einstein-IIIo2}
%Let $\Sigma \hookrightarrow (\mathbb{R} \times \Sigma,ds^2 + \bar{g})$. Then, $\IIIo = 0$ if and only if $(\Sigma,\bar{g})$ is Einstein.
%\end{proposition}
%\begin{proof}
%First, suppose that $(\Sigma,\bar{g})$ is Einstein. Then, $\bar{Sc}$ is constant. Furthermore, $\mathring{\bar{Ric}} = 0$. However, because $\bar{Ric} = Ric$ for this product manifold, we may trace this equation to find that $Sc = \bar{Sc}$. It then implies that
%$$Ric_{ab} - \tfrac{1}{d-1} \bar{g}_{ab} Sc = 0\,.$$
%Expressing this equation in terms of the Schouten tensor and taking the tangential part (and recalling that $\mathring{\bar{P}} = 0$), we have that
%$$P^\top - \bar{P} = 0\,.$$
%But for a product manifold, this implies that $\IIIo = 0$.
%
%On the other hand, if $\IIIo = 0$, then $W_{nabn} = 0$. Expressing this in terms of the Schouten tensor and the Riemann tensor and recalling that $R_{nabn} = 0$ for a product manifold, we have that
%$$P^\top - \bar{g} P_{nn} = 0\,.$$
%But also recall that because $W_{nabn} = 0$, we have that $P^\top - \bar{P}$ for geodesic manifolds (of which a product manifold is an example). So, substituting this in to the above equation and tracing, we find that
%$$\bar{P} - \tfrac{1}{d-1} \bar{J} \bar{g} = 0\,.$$
%That is, $\mathring{\bar{P}} = 0$, so $(\Sigma,\bar{g})$ is Einstein.
%\end{proof}
%
%\bigskip
%
%\begin{lemma} \label{Yamabe-soln2}
%Let $(\Sigma,\bar{g})$ be a constant scalar curvature Riemannian manifold with dimension greater than 2 and let $\Sigma \hookrightarrow (\mathbb{R} \times \Sigma,g:= ds^2 + \bar{g}_{ij} dx^i dx^j)$. Then, the canonical defining density that solves the singular Yamabe problem $\sigma$ is given by
%\begin{align} \label{SY-const-sc2}
%\sigma := \left[g;  \begin{cases}
%\sqrt{\tfrac{(d-1)(2-d)}{\bar{Sc}}} \sin\left(\sqrt{\tfrac{\bar{Sc}}{(d-1)(2-d)}} s\right) & \bar{Sc} < 0 \\
% s  & \bar{Sc}= 0 \\
% \sqrt{\tfrac{(d-1)(d-2)}{\bar{Sc}}} \sinh \left(\sqrt{\tfrac{\bar{Sc}}{(d-1)(d-2)}} s \right)  & \bar{Sc} > 0\,.
%\end{cases} \right]
%\end{align}
%Furthermore, $\sigma$ is continuous.
%\end{lemma}
%\begin{proof}
%By the same arguments found in the proof of Proposition~\ref{Einstein-IIIo}, we have that $(\mathbb{R} \times \Sigma,ds^2 + \bar{g})$ has $Sc = \bar{Sc}$. Now, we compute pointwise at $p \in \mathbb{R} \times \Sigma$.
%
%Suppose first that $\bar{Sc}(x^i(p)) = 0$. Then, it trivially follows that $|ds|^2 - \tfrac{2}{d}(\Delta s + Js) = 1$, and so in the product metric representative, $\sigma = [g; s]$.
%
%Next, suppose that $Sc < 0$. Then, suppose that $\sigma = [g; A \sin (rs)]$. The singular Yamabe equation becomes
%$$A^2 r^2 \cos (rs)^2 + \frac{2A^2(r^2 - J)}{d} \sin(rs)^2 = 1\,.$$
%Thus, we require that $A^2 r^2 = 1$ and that $2A^2 (r^2 - J)/d = 1$. Solving this system and substituting in $J = \tfrac{1}{2(d-1)} \bar{Sc}$, the result follows.
%
%Finally, suppose that $Sc > 0$. Then, suppose that $\sigma = [g; A \sinh (rs)]$. The singular Yamabe equation becomes
%$$A^2 r^2 \cosh (rs)^2 - \frac{2A^2(r^2 + J)}{d} \sinh(rs)^2 = 1\,.$$
%Again, we require that $A^2 r^2 = 1$ and that $2A^2 (r^2 + J)/d = 1$. Solving again (for positive $A$ and $r$), the result for $\sigma$ follows. Continuity follows by verifying that in the $\bar{Sc} \rightarrow 0$ limit, $\sigma \rightarrow [g;s]$ from both above and below.
%\end{proof}
%
%Trivially, then, we obtain the following corollary.
%
%\begin{corollary} \label{nonconst-Sc2}
%Let $(\Sigma,\bar{g})$ be a Riemannian manifold with dimension greater than 2 and let $\Sigma \hookrightarrow (\mathbb{R} \times \Sigma, g := ds^2 + \bar{g}_{ij} dx^i dx^j)$. Then the canonical defining density that solves the singular Yamabe problem $\sigma$ differs from that of Equation~(\ref{SY-const-sc}) by terms only involving derivatives of $\bar{Sc}$.
%\end{corollary}
%
%
%
%\begin{theorem} \label{E-PE2}
%Suppose that $(\Sigma, \bar{g})$ is Einstein and that $\Sigma \hookrightarrow (\mathbb{R} \times \Sigma,ds^2 + \bar{g})$. Then, $(\mathbb{R} \times \Sigma , g^+ := \frac{ds^2 + \bar{g}}{\sigma^2})$ is Poincar\'e--Einstein, where $\sigma$ is the singular Yamabe scale in the metric representative $ds^2 + \bar{g}$.
%\end{theorem}
%\begin{proof}
%To prove this theorem, we need to show that $\mathring{P}^{g^+} = 0$. Recall that in the metric representative $ds^2 + \bar{g}$, a conformally-invariant extension of the second fundamental form is given by
%%First, from Lemma~\ref{even-conf-FFs}, we have that $\IIo = \IVo = \cdots = 0$. Furthermore, from Proposition~\ref{Einstein-IIIo}, it follows that $\IIIo = 0$. We now need to show that all of the higher odd fundamental forms vanish as well. To do so, we will consider the canonical extension of the trace-free second fundamental form:
%$$\IIo^{\rm e}_{ab} = \nabla_{(a} \nabla_{b)\circ} \sigma + \sigma \mathring{P}_{ab}\,.$$
%Indeed, this is a conformal invariant of weight $1$, and hence in the asymptotically hyperbolic scale $g^+$, this tensor equals $\mathring{P}^{g^+}$. So it suffices to show that $\IIo^{\rm e}$ vanishes identically.
%
%Without loss of generality, we will assume that $\bar{Sc} < 0$, which is constant because $(\Sigma,\bar{g})$ is Einstein. We now calculate. First, observe that
%$$\nabla \sigma = Ar \cos (rs) \, ds\,,$$
%where $A$ and $r$ are as in Lemma~\ref{Yamabe-soln}. Next, we have that
%$$\nabla^2 \sigma = -Ar^2 \sin (rs) ds \otimes ds\,.$$
%Thus, the trace-free part of this tensor is
%$$\nabla_{(a} \nabla_{b)\circ} \sigma = -Ar^2 \sin(rs) (n_a n_b - \tfrac{1}{d} g_{ab}) =  \tfrac{d-1}{d} Ar^2 \sin(rs) (\tfrac{1}{d-1} \bar{g}_{ab} - n_a n_b)$$
%
%Now note that that
%\begin{align*}
%\mathring{P}_{ab} &= P_{ab} - \tfrac{1}{d} g_{ab} J \\
%&= P^\top +n_a n_b P_{nn} - \tfrac{1}{d} {g}_{ab} J\\
%&= \bar{P} + n_a n_b P_{nn} - \tfrac{1}{d} g_{ab} J  \\
%&= \mathring{\bar{P}} + \tfrac{1}{d-1} \bar{g}_{ab} \bar{J} + n_a n_b P_{nn} - \tfrac{1}{d} g_{ab} J  \\
%&=  \tfrac{1}{d-1} \bar{g}_{ab} \bar{J} + n_a n_b (J-\bar{J}) - \tfrac{1}{d} g_{ab} J \\
%&= \tfrac{1}{d-1} \bar{g}_{ab} \bar{J} - \tfrac{1}{d-1} n_a n_b \bar{J} - \tfrac{d-2}{d(d-1)} g_{ab} \bar{J} \\
%&= \tfrac{2}{d}\left(\tfrac{1}{d-1} \bar{g}_{ab} -  n_a n_b \right) \bar{J} \\
%&= \tfrac{1}{d(d-2)}\left(\tfrac{1}{d-1} \bar{g}_{ab} -  n_a n_b \right) \bar{Sc}\,.
%\end{align*}
%Using that $r^2 = \frac{\bar{Sc}}{(d-1)(2-d)}$, one may easily check that
%$$\nabla_{(a} \nabla_{b)\circ} \sigma + \sigma \mathring{P}_{ab} = 0\,,$$
%as required.
%%As this equation is conformally invariant, we may express it in the singular Yamabe scale, yielding $\mathring{P}_{ab}^{g^+} = 0$. That is, $(\mathbb{R} \times \Sigma,g^+)$ is Poincar\'e--Einstein.
%
%When $\bar{Sc} = 0$, the identity trivially follows, as $\nabla^2 s = 0$. A similar calculation as the above follows for the case when $\bar{Sc} > 0$.
%\end{proof}
%
%\begin{corollary} \label{Einstein-vanishing-FFs2}
%When $(\Sigma, \bar{g})$ is Einstein and $\Sigma \hookrightarrow (\mathbb{R} \times \Sigma, ds^2 + \bar{g})$ is a smooth embedding, $\FFo{k} = 0$ for each $k \in \{2, \ldots, d-1\}$.
%\end{corollary}
%\begin{proof}
%Because $(\Sigma,\bar{g})$ is Einstein, it follows that $(\mathbb{R} \times \Sigma,\tfrac{ds^2 + \bar{g}}{\sigma^2})$ is Poincar\'e--Einstein. But then the result follows trivially~\cite[Theorem ]{}.
%\end{proof}
%
%\begin{proposition} \label{odd-conf-FFs2}
%Let $\Sigma^{d-1} \hookrightarrow (\mathbb{R} \times \Sigma, ds^2 + \bar{g})$ with $d$ even. Then, $\IIIo = \tfrac{d-3}{d-2} \mathring{\bar{P}}$ and $\FFo{k} = \mathcal{O}_{k}(\IIIo)$ for each $k \in \{4 \ldots,d-1\}$, where $\mathcal{O}_{k}$ is some map on the jets of $\IIIo$ satisfying $\mathcal{O}_k(0) = 0$. Furthermore, $\IIo = 0 = \mathcal{O}_{2n}(\IIIo)$ for each $2 \leq n \leq \tfrac{d-2}{2}$.
%\end{proposition}
%\begin{proof}
%By definition, $\IIIo \= W_{nabn}$. By direct calculation using standard Riemannian hypersurface identities and the fact that for product manifolds, $R_{nabn} \= 0$, we may verify that $\IIIo = \tfrac{d-3}{d-2} \mathring{\bar{P}}$. 
%
%Now consider the asymptotically hyperbolic manifold $(\mathbb{R} \times \Sigma, g^+ := \frac{ds^2 + \bar{g}}{\sigma^2})$, where $\sigma$ is the canonical defining density that solves the singular Yamabe problem as in Corollary~\ref{nonconst-Sc}. Now, by construction~\cite{}, the conformal fundamental forms are expressible in terms of conformally-invariant normal derivative operators acting on the almost Einstein tensor
%$$\IIo^{\rm e}_{ab} := [g; \nabla_{(a} \nabla_{b)\circ} \sigma + \sigma \mathring{P}_{ab}]\,.$$
%Importantly, in the Yamabe scale, we have that $\IIo^{\rm e}_{ab} = [g^+; \mathring{P}^{g^+}]$, so put another way, the conformal fundamental forms are expressible as the action of normal derivative operators on $\mathring{P}^{g^+}$, picking up Taylor coefficients in $s$. So it suffices to show that there exists a formula for $\mathring{P}^{g^+}$ that can be expressed solely in terms of $\IIIo$.
%
%For ease of calculation, we translate to the Ricci curvature: $\mathring{Ric} = (d-2) \mathring{P}$. Now, the conformal rescaling of the Ricci tensor is given by
%$$\mathring{Ric}^{g^+}_{ab} = \mathring{Ric}^g_{ab} + (d-2) \tfrac{\nabla_a \nabla_b \sigma}{\sigma} - \tfrac{d-2}{d} \tfrac{\Delta \sigma}{\sigma} g_{ab}\,.$$
%Because the Riemannian manifold is a product manifold, we have that
%$$\mathring{Ric}_{ab}^g = \mathring{\overline{Ric}}_{ab}^{\bar{g}} + \left(\tfrac{1}{d(d-1)} \bar{g}_{ab} - \tfrac{1}{d} n_a n_b \right) \bar{Sc}\,.$$
%
%To proceed, we must compute the relevant derivative of $\sigma$. However, we only have an explicit formula for $\sigma$ in the case of constant scalar curvatures. Nonetheless, a series expansion of $\sigma$ in this setting differs from that of the constant scalar curvature case by terms that depend only on derivatives of the scalar curvature. Thus, we have that
%\begin{align*}
%\nabla_a \nabla_b \sigma =& -B \sin(Bs) n_a n_b - 2s \sin(Bs) n_{(a} \bar{\nabla}_{b)} B + \left(\frac{Bs \cos(Bs) - \sin(Bs)}{B^2} \right) \bar{\nabla}_a \bar{\nabla}_b B \\
%&- \frac{(s^2 B^2 - 2) \sin(Bs) + 2sB \cos(Bs)}{B^3} (\bar{\nabla}_a B)(\bar{\nabla}_b B) + \mathcal{O}(\bar{\nabla} \bar{Sc})\,.
%\end{align*}
%where 
%$\mathcal{O}(\bar{\nabla} \bar{Sc})$ are terms that depend explicitly on at least one derivative of $\bar{Sc}$.
%
%Note that because $B = \sqrt{\frac{\bar{Sc}}{(d-1)(2-d)}}$ and $\bar{\nabla}_a \bar{Sc} = 2(d-1) \bar{\nabla}^b \mathring{\bar{P}}_{ab}$, we may express any derivatives on $B$ or $\bar{Sc}$ in terms of an operator on $\mathring{\bar{P}}$. So it suffices to check that all terms that do not depend on $\bar{\nabla} B$, $\bar{\nabla} \bar{Sc}$, or $\mathring{\overline{Ric}}$ cancel in the expression for $\mathring{Ric}^{g^+}$. Indeed, those terms are
%$$\tfrac{1}{d(d-1)} \bar{g}_{ab} \bar{Sc} - \tfrac{1}{d} n_a n_b \bar{Sc} + (d-2)(-B^2 n_a n_b) - \tfrac{d-2}{d} (-B^2) g_{ab}\,.$$
%Using that $B^2 = -\tfrac{\bar{Sc}}{(d-1)(d-2)}$, we can check that this vanishes. It thus follows that $\mathring{Ric}^{g^+}_{ab}$ can be expressed as an operator on $\mathring{\bar{P}}$, and hence so can the conformal fundamental forms. From Proposition~\ref{even-conf-FFs}, we also see that the even operators $\mathcal{O}_{2n}$ are the zero operator and that $\IIo = 0$.
%
%As $\mathring{\bar{P}} = \tfrac{d-2}{d-3} \IIIo$, it thus follows that each conformal fundamental form is in fact expressible as an operator on the jets of $\IIIo$, so we denote $\FF{k} := \mathcal{O}_k (\IIIo)$. Furthermore, when $\IIIo = 0$, we have that $(\Sigma,\bar{g})$ is Einstein from Proposition~(\ref{Einstein-IIIo}). Then, from Theorem~(\ref{E-PE}), it follows that the asymptotically hyperbolic manifold is Poincar\'e--Einstein. But then the conformal fundamental forms vanish, again by~\cite{}. Thus, we have that $\mathcal{O}_k(0) = 0$. This completes the proof.
%\end{proof}
%
%\medskip
%
%Given a conformally compact manifold, it is established~\cite{} that there are obstructions to the interior admitting an asymptotically Poincar\'e--Einstein metric on the interior in the conformal class. On the other hand, one may ask what the obstructions there may be to a conformally compact manifold admitting an asymptotic product metric in its conformal class, where one of the leaves of the product is the boundary.
%
%Given Lemma~\ref{even-conf-FFs} and Proposition~\ref{odd-conf-FFs}, it is surely necessary that the even conformal fundamental forms vanish and the odd conformal fundamental forms take on particular values in order for a conformal manifold with boundary to admit a product metric in its conformal class. However, this is not sufficient. To prove this statement, we first need an identification theorem between conformal manifolds with matching conformal fundamental forms.
%
%
%
%\begin{theorem} \label{conformal-identification2}
%Let $\iota : \Sigma \hookrightarrow (M^d,\cc)$ and $\hat{\iota} : \Sigma \hookrightarrow (M^d,\hat{\cc})$ with $d$ even and $\iota(p) = \hat{\iota}(p)$ for all $p \in \Sigma$. Furthermore, let $\sigma$ be a singular Yamabe density for $\iota$. Then, $\iota^* \cc = \hat{\iota}^* \hat{\cc}$ and $\FFo{k} = \hat{\FFo{k}}$ for each $k \leq d-1$ if and only if 
%$$\hat{\cc} = \cc + \mathcal{O}(\sigma^{d-1})\,.$$
%\end{theorem}
%\begin{proof}
%To prove the forward direction, it suffices to show that, a given choice of $\bar{g} \in \iota^* \cc$, there exists metrics $g \in \cc$, $\hat{g} \in \hat{\cc}$ and an isometry $\phi$ such that $\iota^* g = \hat{\iota}^* \hat{g} = \bar{g}$ and $g - \phi^* \hat{g} = \mathcal{O}(s^{d-1})$ for $\sigma = [g;s]$. So, picking $\bar{g}$, we let $g \in \cc$, $\hat{g} \in \hat{\cc}$ be the unique Graham--Lee metrics~\cite{} in the conformal classes, where for $\hat{\sigma} = [\hat{g}, \hat{s}]$ we have that $|ds|_g^2 = 1 =|d\hat{s}|_{\hat{g}}^2$.
%%Note that these defining functions provide a foliation of a collar neighborhood around $\Sigma$, with leaves given by $\Sigma_t = \{p \in M | s(p) = t\} = \{p \in M | \hat{s}(p) = t\}$. Denoting coordinates on $\Sigma$ by 
%
%Because we are using Graham--Lee metric representatives, there always exists isometries $\Phi,\hat{\Phi} : \Sigma \times [0,\epsilon) \rightarrow M$ such that, in a collar neighborhood of $\Sigma$,
%\begin{align*}
%\Phi^* g &= dt^2 + (\bar{g}_{ij}(\vec{x}) + t g^{(1)}_{ij}(\vec{x}) + \tfrac{1}{2} t^2 g^{(2)}_{ij}(\vec{x}) + \cdots) dx^i dx^j  \\
%\hat{\Phi}^* \hat{g} &= dt^2 + (\bar{g}_{ij}(\vec{x}) + t \hat{g}^{(1)}_{ij}(\vec{x}) + \tfrac{1}{2} t^2 \hat{g}^{(2)}_{ij}(\vec{x}) + \cdots) dx^i dx^j \,.
%\end{align*}
%Note that $\Phi^*(s) = \hat{\Phi}^*(\hat{s}) = t$. Thus, the existence of the desired isometry $\phi$ is equivalent to $g^{(k)}(\vec{x}) = \hat{g}^{(k)}(\vec{x})$ for all $k \leq d-2$.
%
%%While we may use the same coordinates on $\Sigma$, on the foliations given by $s$ and $\hat{s}$ respectively, it does not follow that the parallel transport of those coordinates land on the same point. Thus, we will coordinatize by $(s,x^i)$ and $(\hat{s},\hat{x}^i)$, noting that for a point $p \in \Sigma$, $\hat{x}^i(p) = x^i(p)$. Now, we may expand these metrics in $s$ and $\hat{s}$, obtaining
%%\begin{align*}
%%g &= ds^2 + (\bar{g}_{ij} + s g^{(1)}_{ij} + \tfrac{1}{2} s^2 g^{(2)}_{ij} + \cdots) dx^i dx^j  \\
%%\hat{g} &= d\hat{s}^2 + (\hat{\bar{g}}_{ij} + \hat{s} \hat{g}^{(1)}_{ij} + \tfrac{1}{2} \hat{s}^2 \hat{g}^{(2)}_{ij} + \cdots) d\hat{x}^i d\hat{x}^j \,.
%%\end{align*}
%Observe that each coefficient $g^{(k)}$ is a natural hypersurface invariant (that is tangent to the hypersurface), as it can be obtained via $\mathcal{L}_n^{k} \Phi^* g|_{t = 0}$, where $n =\partial_t$, and similarly for $\hat{g}^{(k)}$. As such, from~\cite[Theorem 2.5]{blitzClassifying} it follows that $g^{(k)}$ is expressible as a partial contraction polynomial in
%$$\{\bar{g}, \bar{g}^{-1}, \bar{\nabla}, \bar{R}, \IIo, \,\ldots\, \FFo{k+1}, H, J|_{\Sigma}, \ldots, \nabla_n^{k-2} J|_{\Sigma}\}\,,$$
%and the expression is universal, meaning that the symbolic form of the expression does not depend on the underlying choice of metric.
%
%Note that $|n|^2 = 1$, and so from the singular Yamabe equation $|n|^2 - \tfrac{2}{d} t(\Delta t + Jt) = 1 + \mathcal{O}(t^d)$, we have that $\Delta t + Jt = 0$ in a collar neighborhood of $\Sigma$. So it follows that $\nabla_n^k (\Delta t + Jt) = 0$ for every $1 \leq k \leq d-1$. Expanding this equation, we have
%\begin{align*}
%0 \=& \nabla_n^k \Delta t + (k-1) \nabla_n^{k-1} J \\
%\=& \nabla_n^{k-1} [\nabla_n, \nabla_a] \nabla^a t + (k-1) \nabla_n^{k-1} J \\
%\=& -\nabla_n^{k-1} Ric_{nn} + (k-1) \nabla_n^{k-1} J + \text{ltots} \\
%\=& -(d-1) \nabla_n^{k-1} J + (k-1) \nabla_n^{k-1} J + \text{ltots} \\
%\=& -(d-k) \nabla_n^{k-1} J + \text{ltots.}
%\end{align*}
%In the above, $\text{ltots}$ stands for ``lower transverse order terms,'' which are those terms which involve fewer normal derivatives of the metric. Note that the fourth identity follows from~\cite[Lemma 2.4]{blitzClassifying}. So, it is clear that for Graham-Lee metrics, $\nabla_n^{k-1} J|_{\Sigma}$ is expressible in terms of lower transverse order terms so long as $k < d$. But because $H = 0$ (because $H \= \tfrac{1}{d} \Delta t = -\rho|_{\Sigma}$), it follows by induction that \textit{in a Graham--Lee metric representative}, the coefficients $ g^{(k)}$ have universal expressions in terms of the set
%$$\{\bar{g}, \bar{g}^{-1}, \bar{\nabla}, \bar{R}, \IIo, \,\ldots\, \FFo{k+1}\}\,.$$
%As these expressions are universal (for Graham--Lee metric representatives), it follows that if the intrinsic geometry and the conformal fundamental forms agree for $\iota$ and $\hat{\iota}$, then viewed as tensors on $\Sigma$, $g^{(k)} =  \hat{g}^{(k)}$ for all $k \leq d-2$ (as the conformal fundamental forms are only defined up to that order).
%
%\medskip
%
%The reverse direction follows by noting that if any pair of conformal fundamental forms $\FFo{k}$, $\hat{\FFo{k}}$ are not equal, then it is clear from the Graham--Lee normal form of the metrics that they are not isometric.
%
%
%%Now, recall that for $k \leq d-1$, any conformal fundamental form $\FFo{k}$ may be expressed as
%%$$\FFo{k}_{ab} = \delta^{k-2}(\nabla_{(a} \nabla_{b)_{\circ}} s + s \mathring{P}_{ab})\,,$$
%%for some conformally-invariant operator $\delta^{k-2}$ with leading normal derivative term $\otop \nabla_n^{k-2}$. However, observe that $\nabla_{(a} \nabla_{b)_\circ} s = \tfrac{1}{2} (\mathcal{L}_n g)_{(ab)_\circ}$, where $n = g^{-1}(ds,\cdot)$. So, we may alternatively express
%%$$\FFo{k}_{ab} = \delta^{k-2}(\tfrac{1}{2} (\mathcal{L}_n g)_{(ab)_{\circ}} + s \mathring{P}_{ab})\,.$$
%%Because $\nabla_n$ agrees with $\mathcal{L}_n$ at leading order as a differential operator, it is clear that $\mathcal{L}_n^{k-1} 
%\end{proof}
%
%Using this theorem, we may now prove the following:
%\begin{corollary}
%Let $\Sigma \hookrightarrow (M,\cc)$ be a conformal hypersurface embedding. Suppose that there exists some $\tau \in \Gamma(\ce_+ \Sigma[1])$ such that
%$$(\bar{\nabla}_{(a} \bar{\nabla}_{b)_{\circ}} + \mathring{\bar{P}}_{ab} + (d-2) \Fo) \tau = 0\,,$$
%and that in the $\tau$ scale, $\FFo{k} = \mathcal{O}_k(\IIIo)$. Then $\Sigma \hookrightarrow (M,\cc)$ is an asymptotically conformal product manifold.
%\end{corollary}
%\begin{proof}
%From Theorem~\ref{conformal-identification}, it suffices to show that the conformal hypersurface embedding $\iota: \Sigma \hookrightarrow (M,\cc)$ has conformal fundamental forms that agree with some asymptotically conformal product manifold. Now given a choice of representative $g_{\tau}$ determined by any extension of $\tau$ to $(M,\cc)$, we may identify an isometry $\Phi : \Sigma \times [0,\epsilon) \rightarrow M$ such that $\Phi^* g_{\tau}$ is in geodesic normal form. Now on the same collar neighborhood, we prescribe another metric $\hat{g} := dt^2 + \bar{g}_{ij}(\vec{x})dx^i dx^j + \mathcal{O}(t^{d-1})$, where $\Sigma = \mathcal{Z}(t)$ and $\bar{g} = \Phi^* \iota^* g_{\tau}$ with $\Phi^*$ acting in the natural way on $\Sigma$. This metric in turn defines a conformal class of metrics $\hat{\cc}$ in a collar neighborhood of $\Sigma$, with conformal embedding $\hat{\iota} : \Sigma \hookrightarrow (M,\hat{\cc})$, with $\hat{\iota}^* \hat{\cc} = \iota^* \cc$ and with conformal fundamental forms given by Proposition~\ref{odd-conf-FFs}: $\IIo = 0$, $\IIIo = \tfrac{d-3}{d-2} \mathring{\bar{P}}$, and $\FFo{k} = \mathcal{O}_k(\IIIo)$.
%
%But in the scale $g_{\tau}$, from the hypothesis these are also the conformal fundamental forms of $\Sigma \hookrightarrow (M,\cc)$. The corollary follows.
%
%\end{proof}
%
%
%\bigskip
%
%$${}$$
%
%\bigskip
%
%\begin{proposition}
%Let $\iota: \Sigma \hookrightarrow (M,\cc)$ and let $\tau \in \Gamma(\ce_+ \Sigma[1])$ be a solution to $\mathcal{J}^{(\IIo)} \tau = 0$. Then, for each $f \in C_+^\infty \Sigma$ such that $\bar{\nabla}^{\bar{g}_{\tau/f}} \bar{\nabla}^{\bar{g}_{\tau/f}} f = 0$,
%$$\mathcal{J}^{(\IIo)} \tau/f = 0\,.$$
%\end{proposition}
%
%\begin{proof}
%By defininition, $g_{\tau/f} := f^2 g_{\tau} = f^2 \bg/\tau^2$. Applying $\mathcal{J}^{(\IIo)}$ to $\tau/f$, 
%and using the fact that $\mathcal{J}^{(\IIo)} \tau = 0$, we have that in any scale $\tau = [\bar{g};t]$,
%\begin{align*}
%\mathcal{J}^{(\IIo)} \tau/f =& \tau \bar{\nabla}_{(a}^{\bar{g}} \bar{\nabla}_{b)_{\circ}}^{\bar{g}} f^{-1} + 2 d \tau_{(a} d(\tfrac{1}{f})_{b)_{\circ}} \\
%=& -\tau \bar{\nabla}_{(a}^{\bar{g}} f^{-2} \bar{\nabla}_{b)_{\circ}}^{\bar{g}} f - 2f^{-2} d \tau_{(a} df_{b)_{\circ}} \\
%=& 2\tau f^{-3} (\bar{\nabla}_{(a}^{\bar{g}} f)(\bar{\nabla}_{b)_{\circ}}^{\bar{g}} f) - \tau f^{-2} \bar{\nabla}_{(a}^{\bar{g}} \bar{\nabla}_{b)_{\circ}}^{\bar{g}} f- 2f^{-2} d \tau_{(a} df_{b)_{\circ}} \\
%=& 2 \tau f^{-3} (\bar{\nabla}_{(a}^{\bar{g}} f)(\bar{\nabla}_{b)_{\circ}}^{\bar{g}} f) - \tau f^{-2} \bar{\nabla}_{(a}^{\bar{g}} \bar{\nabla}_{b)_{\circ}}^{\bar{g}_{\tau/f}} f- 2f^{-2} d \tau_{(a} df_{b)_{\circ}} \\
%=& 2 \tau f^{-3} (\bar{\nabla}_{(a}^{\bar{g}} f)(\bar{\nabla}_{b)_{\circ}}^{\bar{g}} f) - \tau f^{-2} (\bar{\nabla}_{(a}^{\bar{g}_{\tau/f}} \bar{\nabla}_{b)_{\circ}}^{\bar{g}_{\tau/f}} f - 2 (\bar{\nabla}_{(a}^{\bar{g}} \log \tfrac{f}{\tau})(\bar{\nabla}_{b)_{\circ}}^{\bar{g}} f)) - 2f^{-2} d \tau_{(a} df_{b)_{\circ}} \\
%=& 4\tau  f^{-3} (\bar{\nabla}_{(a}^{\bar{g}} f)(\bar{\nabla}_{b)_{\circ}}^{\bar{g}} f) -  4f^{-2} d \tau_{(a} df_{b)_{\circ}}\,.
%\end{align*}
%Observe that in the $\tau/f$ scale, we have that $\tau = [g_{\tau/f}; f]$. Imposing that scale on the above operator action, $J^{(\IIo)} \tau/f = 0$. But this is a conformally-invariant operator acting on a scale, so it must be the case that it holds in all scales.
%\end{proof}
%
%\begin{remark}
%From the above proposition, it is clear that when $\mathring{\bar{P}}^{\bar{\bg}_{\tau}} = -(d-2) \Fo$, it is precisely those functions $f$ whose Hessians are zero in the $\tau/f$ scale that characterize additional solutions where $\mathring{\bar{P}}^{\bar{g}_{\tau/f}} = -(d-2) \Fo$.
%\end{remark}
%
%\bigskip
%
%
%\begin{theorem}
%Let $\iota: \Sigma \hookrightarrow (M,\cc)$ be a conformal hypersurface embedding. There exists a nowhere-vanishing smooth conformal Killing 1-form $\xi \in \Gamma(T^* M[2])$ such that $\xi|_{\Sigma} \in \Gamma(N^* \Sigma[2])$ and that $\xi = |\xi|^2 df$ for some function $f \in C^\infty M$ if and only if $\Sigma \hookrightarrow (M,\cc)$ is a conformal product manifold.
%\end{theorem}
%\begin{proof}
%First, suppose that $\Sigma \hookrightarrow (M,\cc)$ is a conformal product manifold. Then it follows that there exists a representative $g \in \cc$ where $g = dt^2 \oplus g_{\Sigma}$. In that representative we define $\xi = dt \in \Gamma(T^*M)$. Clearly, $\xi|_{\Sigma} \in \Gamma(N^* \Sigma)$. It is evident that $\mathcal{L}_{\xi^{\sharp}} g = 0$, and so in particular $\xi$ is a conformal Killing 1-form. Finally, in these coordinates choosing $f = t$, we find that indeed $\xi = |\xi|^2 df$. Promoting this 1-form to a weight-2 1-form on the conformal manifold completes this direction.
%
%For the other direction, suppose that such a 1-form $\xi \in \Gamma(T^* M[2])$ exists on $(M,\cc)$. By hypothesis there exists $f \in C^\infty M$ such that $\xi = |\xi|^2 df$. Further, because $\xi$ is a conformal Killing 1-form, it follows that in any metric representative $g \in \cc$,
%$$\nabla_a \xi_b = \mu_{ab} + \nu g_{ab}\,,$$
%where $\mu \in \Gamma(\wedge^2 T^* M)$ and $\nu := \tfrac{1}{d} \nabla_a \xi^a$. Note that the 1-form $\xi$ is integrable when $\xi \wedge \mu = 0$.
%
%Because $\xi$ is nowhere-vanishing, we may work in the scale $g_{\xi} = \bg/|\xi|_{\bg}^2$; in this scale, $\xi = [g_{\xi}; df]$ and $|\xi|^2_{g_{\xi}} = 1$. So, in that representative, it follows that $\mu = d^2 f = 0$, so $\xi$ is integrable and hence $\ker \xi$ foliates $(M,g_{\xi})$. As $\xi|_{\Sigma} \in \Gamma(N^* \Sigma[2])$ by hypothesis, it follows that $T \Sigma \subset \ker \xi$, and so $\xi|_{\Sigma}$ is a unit conormal for $\Sigma$ in this scale. As such, $\iota^* g_{\xi} = (g_{\xi} - \xi \otimes \xi)|_{\Sigma}$.
%
%Next, observe that in this scale $\nabla |\xi|^2 = 0$, and so
%$$0 = \nabla_a |\xi|^2 = 2\xi^b \nabla_a \xi_b = 2 \nu \xi_a \\,,$$
%meaning that $\nu = 0$ in this scale. It follows that $\xi$ is parallel. Because $\xi = [g_{\xi}; df]$ and $\xi$ is conormal to $\Sigma$, we have that $f|_{\Sigma} = c$ for some constant $c$. Defining $f' = f|_{\Sigma} - c$, it follows that $f'$ is a defining function for $\Sigma$.
%
%Observe that $\xi$ and $\bar{g}_{\xi} := g_{\xi} - \xi \otimes \xi$ are parallel; it follows that $\mathcal{L}_{\xi^{\sharp}} \bar{g}_{\xi} = 0$. In geodesic normal coordinates $(f', \vec{x})$, this implies that $\bar{g}_{\xi}$ is independent of $f'$. So we may write
%$$g_{\xi} = (df')^2 + \bar{g}_{\xi}(\vec{x})\,,$$
%explicitly a product metric.
%\end{proof}
%
%As a consequence of the above theorem, we may study the extrinsic conformal geometry of $\Sigma \hookrightarrow (M,\cc)$ to determine (at least asymptotically) when the embedding is a conformal product manifold. Indeed, one may always find a defining \textit{function} $s \in \Gamma(\ce M[0])$ for $\Sigma$ such that $ds$ is nowhere vanishing and $ds|_{\Sigma} \in \Gamma(N^* \Sigma[0])$. Many such defining functions exist---but for any choice of scale, we may associate a unique such a defining function $s$ to a choice of scale $g_s \in \cc$ by demanding that $|ds|_{g_s}^2 = 1$.  Now let $\nu = \bg^{-1}(ds,ds) \in \Gamma(\ce_+ M[-2])$ and define $\xi^s := \mu ds \in \Gamma(T^* M[2])$ with $\mu = \nu^{-1}$.
%
%Observe that $|\xi^s|^2_{\bg} > 0$ everywhere, and hence is non-vanishing. Furthermore, observe that $\xi = |\xi^s|_{\bg}^2 ds$ by construction. So showing that $\Sigma \hookrightarrow (M,\cc)$ is an asymptotically conformal product manifold is equivalent to showing that each transverse jet of $\nabla_{(a} \xi^s_{b)_{\circ}}$ vanishes on $\Sigma$.
%
%\medskip
%
%At zeroth order, this is the condition
%$$\nabla_{(a} \xi^s_{b)_{\circ}} \= 0\,.$$
%Observe that along $\Sigma$, it follows that $ds/|ds|_g \= n$, the unit conormal. (In fact, this implies that $\nabla_n s \= \mu^{-1/2}$.) Using this observation and standard hypersurface identities, the zeroth order constraint becomes
%$$\mu^{1/2} \IIo_{ab}  + \mu  n_{(a} n_{b)_{\circ}} \left(\mu^{-3/2} \nabla_n \mu + n^c \nabla_n \nabla_c s\right) \= 0\,.$$
%Therefore, from the zeroth order constraint we find that
%$$\IIo \= 0 \qquad \text{ and } \qquad n^c \nabla_n \nabla_c s \= -\mu^{-3/2} \nabla_n \mu\,,$$
%because $\mu \neq 0$. 
%
%\medskip
%
%At first order, we need to decompose the jet into transverse and tangential components. Using the operator $\delta_R : \Gamma(\odot^2 T^* M[w]) \rightarrow \Gamma(\odot^2 T^* \Sigma[w-1])$ with leading transverse operator $\otop \nabla_n$, we have a second constraint,
%$$\delta_R \nabla_{(a} \xi^s_{b)_{\circ}} = \IIo_{ab} (n^c \nabla_n \xi^s_c + H \xi^s_n)- \left(\bar{\nabla}_{(a} \bar{\nabla}_{b)_{\circ}} + \mathring{\bar{P}}_{ab} + (d-2) \Fo_{ab} - \IIo^2_{(ab)_{\circ}}  \right) \xi^s_n = 0\,.$$
%Now observe that along $\Sigma$, $n^a n^b \nabla_{(a} \xi^s_{b)_{\circ}}$ is conformally invariant. So evaluating in the scale $g_s$ and noting that $n \cdot \IIo \= 0$, it follows that $n^a n^b \nabla_{(a} \xi^s_{b)_{\circ}} \= 0$ in every scale. Now defining $K_{ab} := \nabla_{(a} \xi^s_{b)_{\circ}}$, it follows that
%$$\top(n^b \nabla_n K_{ba}- \tfrac{1}{d-1} \bar{\nabla} \csdot K_a^\top) \= 0$$
%is a conformally invariant constraint equation.
%Noting that $K_{ab}^\top \= \IIo$, the above reduces simply to $(n^b \nabla_n K_{ab})^\top = 0$. Expanding this constraint, we have that
%\begin{align*}
%0 \=& (n^b \nabla_n K_{ab})^\top \\
%\=& (n^b \nabla_n \nabla_{(a} \xi^s_{b)_{\circ}})^\top \\
%\=& (n^b n^c \nabla_{(a} \nabla_{|c|} \mu \nabla_{b)_{\circ}} s)^\top + (\mu n^b n^c R_{c(ab)_{\circ}d} \nabla^d s)^\top \\
%\=& (\tfrac{1}{2} n^b n^c \nabla_{a} \nabla_{c} \mu \nabla_{b} s + \tfrac{1}{2} n^b n^c \nabla_b \nabla_c \mu \nabla_a s - \tfrac{1}{d} g_{ab} n^b n^c \nabla_d \nabla_c \mu \nabla^d s)^\top \\
%\=& (\tfrac{1}{2} n^b n^c \nabla_{a} \nabla_{c} \mu \nabla_{b} s + \tfrac{1}{2} n^b n^c \nabla_b \nabla_c \mu \nabla_a s)^\top \\
%\=& \left(\tfrac{1}{2} n^b n^c \nabla_a^\top \left[(\nabla_c \mu)(\mu^{-1/2} n_b) + \mu \nabla_c \nabla_b s \right] + \tfrac{1}{2} n^b n^c \nabla_b \left[(\nabla_c \mu)(\nabla_a s) + \mu \nabla_a \nabla_c s\right] \right)^\top \\
%\=& \top \Big(\tfrac{1}{2} n^b\nabla_a^\top [(\nabla_n \mu)(\mu^{-1/2} n_b)] - \tfrac{1}{2} \mu^{-1/2} \II^c_a (\bar{\nabla}_c \mu)  + \tfrac{1}{2} n^b n^c \nabla_a^\top [\mu \nabla_c \nabla_b s]  \\&+ \tfrac{1}{2} n^b n^c \nabla_b \left[(\nabla_c \mu)(\nabla_a s) + \mu \nabla_a \nabla_c s\right] \Big) \\
%\=& \top \Big(\tfrac{1}{2} \nabla_a^\top (\mu^{-1/2} \nabla_n \mu)- \tfrac{1}{2} \mu^{-1/2} \II^c_a (\bar{\nabla}_c \mu)  + \tfrac{1}{2} n^b n^c \nabla_a^\top [\mu \nabla_c \nabla_b s]  \\&+ \tfrac{1}{2} n^b n^c \nabla_b \left[(\nabla_c \mu)(\nabla_a s) + \mu \nabla_a \nabla_c s\right] \Big) \\
%\=& \top \Big(-\tfrac{1}{4} \mu^{-3/2} (\bar{\nabla}_a \mu)(\nabla_n \mu) + \tfrac{1}{2} \mu^{-1/2} \bar{\nabla}_a \nabla_n \mu - \tfrac{1}{2} \mu^{-1/2} \II^c_a (\bar{\nabla}_c \mu)  + \tfrac{1}{2} n^b n^c \nabla_a^\top [\mu \nabla_c \nabla_b s]  \\&+ \tfrac{1}{2} n^b n^c \nabla_b \left[(\nabla_c \mu)(\nabla_a s) + \mu \nabla_a \nabla_c s\right] \Big) \\
%\=& \top \Big(-\tfrac{1}{4} \mu^{-3/2} (\bar{\nabla}_a \mu)(\nabla_n \mu) + \tfrac{1}{2} \mu^{-1/2} \bar{\nabla}_a \nabla_n \mu - \tfrac{1}{2} \mu^{-1/2} \II^c_a (\bar{\nabla}_c \mu)  + \tfrac{1}{2} n^b  \nabla_a^\top [\mu n^c \nabla_c \nabla_b s]  \\& - \tfrac{1}{2} n^b \II^c_a \mu \nabla_c \nabla_b s + \tfrac{1}{2} n^b n^c \nabla_b \left[(\nabla_c \mu)(\nabla_a s) + \mu \nabla_a \nabla_c s\right] \Big) \\
%\=& \top \Big(-\tfrac{1}{4} \mu^{-3/2} (\bar{\nabla}_a \mu)(\nabla_n \mu) + \tfrac{1}{2} \mu^{-1/2} \bar{\nabla}_a \nabla_n \mu - \tfrac{1}{2} \mu^{-1/2} \II^c_a (\bar{\nabla}_c \mu)  + \tfrac{1}{2}   \nabla_a^\top [\mu n^b n^c \nabla_c \nabla_b s]  \\& -  n^b \II^c_a \mu \nabla_c \nabla_b s + \tfrac{1}{2} n^b n^c \nabla_b \left[(\nabla_c \mu)(\nabla_a s) + \mu \nabla_a \nabla_c s\right] \Big) \\
%\=& \top \Big(-\tfrac{1}{4} \mu^{-3/2} (\bar{\nabla}_a \mu)(\nabla_n \mu) + \tfrac{1}{2} \mu^{-1/2} \bar{\nabla}_a \nabla_n \mu - \tfrac{1}{2} \mu^{-1/2} \II^c_a (\bar{\nabla}_c \mu)  - \tfrac{1}{2}   \nabla_a^\top [\mu^{-1/2} \nabla_n \mu]  \\& -  n^b \II^c_a \mu \nabla_c \nabla_b s + \tfrac{1}{2} n^b n^c \nabla_b \left[(\nabla_c \mu)(\nabla_a s) + \mu \nabla_a \nabla_c s\right] \Big) \\
%\=& \top \Big( - \tfrac{1}{2} \mu^{-1/2} \II^c_a (\bar{\nabla}_c \mu)   -  n^b \II^c_a \mu \nabla^\top_c \nabla_b s  + \tfrac{1}{2} n^b n^c \nabla_b \left[(\nabla_c \mu)(\nabla_a s) + \mu \nabla_a \nabla_c s\right] \Big) \\
%\=& \top \Big( - \tfrac{1}{2} \mu^{-1/2} \II^c_a (\bar{\nabla}_c \mu)   -  n^b \II^c_a \mu \nabla^\top_c \nabla_b s  + \tfrac{1}{2} n^b  (\nabla_n \mu)(\nabla^\top_a \nabla_b s)  + \tfrac{1}{2} n^b n^c \nabla_b \left[ \mu \nabla_a \nabla_c s\right]  \Big) \\
%\=& \top \Big( - \tfrac{1}{2} \mu^{-1/2} \II^c_a (\bar{\nabla}_c \mu)   -  n^b \II^c_a \mu \nabla^\top_c (\mu^{-1/2} n_b)  + \tfrac{1}{2}  (\nabla_n \mu)(\nabla^\top_a \mu^{-1/2})  + \tfrac{1}{2} n^b n^c \nabla_b \left[ \mu \nabla_a \nabla_c s\right]  \Big) \\
%\=& \top \Big( \tfrac{1}{2}  (\nabla_n \mu)(\nabla^\top_a \mu^{-1/2})  + \tfrac{1}{2} n^b n^c \nabla_b \left[ \mu \nabla_a \nabla_c s\right]  \Big) \\
%\=& \top \Big( -\tfrac{1}{4} \mu^{-3/2}  (\nabla_n \mu) \bar{\nabla}_a \mu  + \tfrac{1}{2} n^b n^c \nabla_b \left[ \mu \nabla_a \nabla_c s\right]  \Big) \\
%\=& \top \Big( -\tfrac{1}{4} \mu^{-3/2}  (\nabla_n \mu) \bar{\nabla}_a \mu  + \tfrac{1}{2}  n^c  (\nabla_n \mu) \nabla^\top_a \nabla_c s + \tfrac{1}{2} n^b n^c \mu \nabla_b \nabla_a \nabla_c s  \Big) \\
%\=& \top \Big( -\tfrac{1}{2} \mu^{-3/2}  (\nabla_n \mu) \bar{\nabla}_a \mu  + \tfrac{1}{2} n^b n^c \mu \nabla_a \nabla_b \nabla_c s + \tfrac{1}{2} n^b n^c R_{bacd} \nabla^d s \Big)  \\
%\=& \top \Big( -\tfrac{1}{2} \mu^{-3/2}  (\nabla_n \mu) \bar{\nabla}_a \mu  + \tfrac{1}{2} n^b n^c \mu \nabla_a \nabla_b \nabla_c s  \Big) \\
%\=& \top \Big( -\tfrac{1}{2} \mu^{-3/2}  (\nabla_n \mu) \bar{\nabla}_a \mu  + \tfrac{1}{2} n^c \mu \nabla^\top_a \nabla_n \nabla_c s - \tfrac{1}{2}  \mu n^c \II_a^b \nabla^\top_b \nabla_c s \Big) \\
%\=& \top \Big( -\tfrac{1}{2} \mu^{-3/2}  (\nabla_n \mu) \bar{\nabla}_a \mu  + \tfrac{1}{2} \mu \nabla^\top_a n^c \nabla_n \nabla_c s - \tfrac{1}{2} \mu \II^c_a n^b \nabla_c^\top  \nabla_b s - \tfrac{1}{2}  \mu n^c \II_a^b \nabla^\top_b \nabla_c s \Big) \\
%\=& \top \Big( -\tfrac{1}{2} \mu^{-3/2}  (\nabla_n \mu) \bar{\nabla}_a \mu  + \tfrac{1}{2} \mu \nabla^\top_a n^c \nabla_n \nabla_c s - \tfrac{1}{2} \mu \II^c_a  \nabla_c^\top  \mu^{-1/2} - \tfrac{1}{2}  \mu n^c \II_a^b \nabla^\top_b \nabla_c s \Big) \\
%\=& \top \Big( -\tfrac{1}{2} \mu^{-3/2}  (\nabla_n \mu) \bar{\nabla}_a \mu  - \tfrac{1}{2} \mu \nabla^\top_a (\mu^{-3/2} \nabla_n \mu) + \tfrac{1}{4} \mu^{-1/2} \II^c_a  \bar{\nabla}_c \mu  - \tfrac{1}{2}  \mu n^c \II_a^b \nabla^\top_b \nabla_c s \Big) \\
%\=& \top \Big( \tfrac{1}{4} \mu^{-3/2}  (\nabla_n \mu) \bar{\nabla}_a \mu  - \tfrac{1}{2} \mu^{-1/2} \bar{\nabla}_a  \nabla_n \mu + \tfrac{1}{4} \mu^{-1/2} \II^c_a  \bar{\nabla}_c \mu  - \tfrac{1}{2}  \mu n^c \II_a^b \nabla^\top_b \nabla_c s \Big) \\
%\=& \top \Big( \tfrac{1}{4} \mu^{-3/2}  (\nabla_n \mu) \bar{\nabla}_a \mu  - \tfrac{1}{2} \mu^{-1/2} \bar{\nabla}_a  \nabla_n \mu + \tfrac{1}{2} \mu^{-1/2} \II^c_a  \bar{\nabla}_c \mu  \Big) \\
%\end{align*}
%
%
%For the last irreducible component of $\nabla_n K$, we have the conformally invariant constraint equation
%$$n^a n^b \nabla_n K_{ab} - \tfrac{2}{d-2} \bar{\nabla}^a K_{na}^\top \= 0\,.$$
%
%Now, taking advantage of the zeroth order constraint, the constraint equations become
%$$\delta_R \nabla_{(a} \xi^s_{b)_{\circ}} = \left(\bar{\nabla}_{(a} \bar{\nabla}_{b)_{\circ}} + \mathring{\bar{P}}_{ab} + (d-2) \Fo_{ab} \right) \sqrt{\mu} \= 0\,,$$
%$$(n^b n^c \nabla_b \nabla_c \xi^s_a)^\top \= \mu^{-2}  (\bar{\nabla}_a \mu) \left(- \nabla_n \mu - \tfrac{3}{2} \mu^{1/2}  + \mu^{3/2} H \right) = 0\,,$$
%
%\bigskip
%
%
%
%\newpage
%$${}$$
%\newpage
%
%\section{Introduction}
%In 1867, Bonnet proved the fundamental theorem of surfaces, which, in modern language, asserts that up to rigid Euclidean motions, a surface embedded in $\mathbb{R}^3$ is uniquely determined by its first and second fundamental forms. This result easily generalizes to hypersurfaces embedded in $\mathbb{R}^d$. The naming of these tensors begs for explanation: is there a third fundamental form? A fourth? Indeed there are---but in the context of Euclidean embeddings, they either vanish or are easily expressible in terms of the first and second fundamental forms~\cite[Volume 3]{Spivak} (depending on the definition provided). However, in recent work~\cite{BGW1} it was shown that the conformally-invariant versions of these higher fundamental forms play a critical role in the geometric description of Poincar\'e--Einstein manifolds. It is the goal of this article to provide a series of results providing some geometric meaning for higher fundamental forms in the Riemannian context.
%
%To that end, consider the following definitions.
%\begin{definition}
%Let $\iota : \Sigma \hookrightarrow (M,g)$ and let $n \in \Gamma(T M)$ be any vector such that $g(n,n)|_{\Sigma} = 1$. The \textit{first fundamental form} $\bar{g} := \iota^* g = g - n \otimes n$ is the pullback of $g$ to the hypersurface; this is also often simply referred to as the induced metric. The \textit{second fundamental form} is defined according to $\II := \iota^* \nabla g(n,\cdot)$. We define the \textit{third fundamental form} according to $\III := \iota^* g(R(n,\cdot) n, \cdot)$, and finally, for $k \geq 4$, we define the \textit{$k$th fundamental form} using abstract index notation:
%$$\FF{k}_{ab} := \iota^* (n^{a_1} \cdots n^{a_{k-3}} n^c n^d \nabla_{a_1} \cdots \nabla_{a_{k-3}} R_{cabd})\,.$$
%\end{definition}
%For definition of the symbols and notations used above and in what follows, see Subsection~\ref{conventions}.
%
%While the above definition may seem ad hoc, it is a natural generalization of of the second fundamental form. To see this, note that for an outward-pointing unit normal vector $\bar{n} \in \Gamma(T M)|_{\Sigma}$, we may always~\cite{wald1984} find a function $s \in C^\infty M$ such that $g^{-1}(ds,\cdot)|_{\Sigma} = \bar{n}$ and $g^{-1}(ds,ds) = 1$ in a collar neighborhood around $\Sigma$. This function determines a canonical normal vector $g^{-1}(ds,\cdot) =: n$ which we may use to compactly express the fundamental forms. For example, because $\nabla_n n = 0$, it follows that
%$$\II_{ab} = \nabla_a n_b|_{\Sigma}\,.$$
%Furthermore, a simple calculation shows that
%$$\nabla_n \nabla_a n_b|_{\Sigma} = \III_{ab} - \II_a^c \II_{cb}$$
%and that
%$$\nabla_n^2 \nabla_a n_b|_{\Sigma} = \IV_{ab} - 2 (\II \cdot \III)_{(ab)} + 2 (\II^3)_{ab}\,.$$
%This pattern continues: the leading structure of $\nabla_{n}^{k-3} \nabla n|_{\Sigma}$ agrees with $\FF{k}$ for all $k \geq 3$. Note that the definition of the third fundamental form given in~\cite[Volume 3]{Spivak} is $-\nabla_n \nabla_a n_b|_{\Sigma}$, and the fourth fundamental form is proportional to $\nabla_n^2 \nabla_a n_b|_{\Sigma}$.
%
%\medskip
%
%The structure of this article is quite linear. Section~\ref{2ff} reviews a small subset of known results indicating the geometric significance of the second fundamental form. Section~\ref{3ff} provides several new geometric results that affirm the relationship between the third fundamental form and the existence of conjugate points on and off a hypersurface. Section~\ref{4ff} opens with a qualitative description of the geometric significance of the fourth fundamental form, and then provides a simple example of this geometric significance as well as a conjecture about more general behavior. Finally, in Section~\ref{obstructions}, we show that the fundamental forms play a role similar to those of the conformal fundamental forms in the asymptotically Poincar\'e--Einstein setting. Finally, we conclude with a discussion of generalizations of higher fundamental forms to the higher-codimension setting.
%
%\subsection{Notations and Conventions} \label{conventions}
%In this article, we denote by the pair $(M,g)$ a smooth $d$-dimensional Riemannian manifold equipped with a metric $g$. On this manifold we denote by $\nabla$ the Levi-Civita connection with Riemann curvature tensor $R$, defined via the convention
%$$R(x,y)z = \nabla_x \nabla_y z - \nabla_y \nabla_x z - \nabla_{[x,y]} z\,.$$
%Here, we have that $x,y,z \in \Gamma(TM)$ are vector fields in $M$ and $[x,y]$ is the Lie derivative of $y$ with respect to $x$.
%
%Now, for ease of computation, we use abstract index notation: we label sections of a tensor product of copies of the tangent bundle and cotangent bundle with Latin letters from the beginning of the alphabet (so $a$, $b$, $c$, etc.) to denote their tensor type. Then, a vector field $x$ would be labelled as $x^a$ and a covector (or a 1-form) $\omega$ would be labelled as $\omega_a$. With this notation, tensor contraction is indicated with repeated upper and lower indices. It follows that the above definition of the Riemann curvature tensor is equivalent to the definition
%$$R_{ab}{}^c{}_d z^d = \nabla_a \nabla_b z^c - \nabla_b \nabla_a z^c\,.$$
%We also use the convention for the Ricci curvature tensor given by $Ric_{ab} = R_{ca}{}^c{}_b$.
%
%For notational brevity, we also introduce several shorthands for various symmetrizations and products. We use round brackets $()$ surrounding the relevant indices to denote symmetrization of tensors, and we use square brackets $[]$ for antisymmetrization; for example, we would write $T_{(ab)}$ for the symmetric piece of a tensor $T_{ab}$ and $T_{[ab]}$ for the antisymmetric piece (with unit normalization). Furthermore, we use dot-products for appropriate contraction of indices, so that we might express $S_a{}^c T_{cb} =: (S \cdot T)_{ab}$. Similarly, we can write $(T^2)_{ab} := T_a{}^c T_{cb}$ and the natural generalization for higher-order products. Finally, when it is useful, we will place the symbol for a (co)vector field in the subscript of a tensor to indicate contraction; for example, if $x$ and $y$ are vector fields and $T_{ab}$ is some tensor field, we may write $T_{xy} := x^a y^b T_{ab}$.
%
%In the context of hypersurface embeddings, we label the induced tensors, connection, and curvatures with an overline $\bar{\bullet}$; for example, we would write the induced Levi-Civita connection on a hypersurface by $\bar{\nabla}$, and the corresponding Riemann curvature tensor as $\bar{R}$. Also note that we will use the same indices for tensor structures on a hypersurface because of the canonical isomorphism between $TM|_{\Sigma}$ and $T \Sigma \oplus \langle n \rangle$. We will also indicate that two expressions are equivalent along the hypersurface via the stacked equality $\=$. Furthermore, when we wish to indicate that a tensor or an operator is projected to the hypersurface, we will often decorate that tensor or operator with a $\top$; alternatively, we will use $\top$ as the projection operator.
%
%
%\section{Geometry of the second fundamental form} \label{2ff}
%Fundamentally, the second fundamental form links the connection on $TM|_{\Sigma}$ to the connection on $T \Sigma$. This is encapsulated by Gauss' formula:
%\begin{align} \label{gauss-formula}
%\bar{\nabla}_a \bar{v}^b \= \nabla^\top_a v^b + n^b \II_{ac} v^c\,,
%\end{align}
%where $\bar{v} \in \Gamma(T \Sigma)$ and $v \in \Gamma(TM)$ is any extension of $\bar{v}$ off $\Sigma$. As an aside, this result directly implies the Gauss equation, which relates the (projected) bulk curvature tensor and the induced curvature tensor:
%\begin{align} \label{gauss-equation}
%\bar{R}_{abcd} = R_{abcd}^\top + \II_{ac} \II_{bd} - \II_{ad} \II_{bc}\,.
%\end{align}
%Another useful result that follows from Equation~(\ref{gauss-formula}) is the Codazzi--Mainardi equation
%\begin{align} \label{cod}
%R_{abcn}^\top = \bar{\nabla}_a \II_{bc} - \bar{\nabla}_b \II_{ac}\,.
%\end{align}
%
%A straightforward and well-known geometric result follows directly from Equation~(\ref{gauss-formula}):
%\begin{theorem}
%Let $\Sigma \hookrightarrow (M,g)$ be a hypersurface embedding with $\II = 0$. Then, any geodesic in $(\Sigma,\bar{g})$ is pushed-forward to a geodesic in $(M,g)$.
%\end{theorem}
%\begin{proof}
%Let $\bar{\gamma}$ be a geodesic in $(\Sigma,\bar{g})$. Then by definition we have that $\bar{\nabla}_{\dot{\bar{\gamma}}} \dot{\bar{\gamma}} = 0$. Now, define $\gamma := \iota_* \bar{\gamma}$. It follows trivially that $\dot{\gamma} := \iota_* \dot{\bar{\gamma}}$ is the tangent vector to $\gamma$. But then it follows from Equation~(\ref{gauss-formula}) and the fact that $\II = 0$ that $\nabla_{\dot{\gamma}} \dot{\gamma} = 0$. Hence, we have that $\gamma$ is a geodesic in $(M,g)$. 
%\end{proof}
%It is for this reason that we say that a hypersurface is \textit{totally geodesic} when $\II = 0$. Intuitively, then, the second fundamental form characterizes the deviation of geodesics in $\Sigma$ from geodesics in $M$.
%
%However, the second fundamental form does not just inform on geodesics---it also characterizes other special curves in relation to the hypersurface. In the theory of curves in Riemannian manifolds, a curve may be characterized by its curvature and torsion functions. For a curve $\gamma : I \rightarrow M$ (where $I$ is some interval) we have that the curvature of $\gamma$ is given by the magnitude of the acceleration of its tangent vector, namely $\kappa := |\nabla_{\dot{\gamma}} \dot{\gamma}|$. In three manifolds, the torsion can be defined as the magnitude of the projection of the jerk vector onto the direction perpendicular to the osculating plane. In higher dimensions~\cite{jordan1874}, torsions can be similarly defined by a Gram-Schmidt procedure and computation of the lengths of projected components of increasing number of derivatives of the tangent vector. We will not need explicit formulae for these---to find such formulae, see~\cite{klingenberg1978}. In this language, geodesics are those curves with vanishing curvature and torsions.
%
%Given this description of curves, we can also consider the geometric significance of the second fundamental as it relates to other special curves besides just geodesics. The first example is that hypersurfaces with (symmetric) covariantly constant second fundamental forms have geodesics that push-forward to constant-curvature curves:
%\begin{proposition}
%Let $\iota : \Sigma \hookrightarrow (M^d,g)$ be a hypersurface embedding and let $\bar{\gamma}$ be a geodesic of $(\Sigma,\bar{g})$. If $\bar{\nabla}_{(a} \II_{bc)} = 0$, then $\iota_* \bar{\gamma}$ has  constant curvature. \edz{mention arclength parametrization somewhere in the intro perhaps}
%\end{proposition}
%
%\begin{proof}
%Because $\bar{\gamma}$ is a geodesic, we have that $\bar{\nabla}_{\dot{\bar{\gamma}}} \dot{\bar{\gamma}} = 0$. Pushing forward to $M$, it follows by the Gauss equation that
%\begin{align} \label{bulk-geo1}
%\nabla_{\dot{\gamma}} \dot{\gamma} + \II(\dot{\gamma},\dot{\gamma}) n = 0\,.
%\end{align}
%Now note that the curvature for an arclength parametrized curve is given by
%$$\kappa := |\nabla_{\dot{\gamma}} \dot{\gamma}|_{g}\,.$$
%So squaring Equation~(\ref{bulk-geo1}), we find that
%$$|\nabla_{\dot{\gamma}} \dot{\gamma}|^2 + 2\langle n, \nabla_{\dot{\gamma}} \dot{\gamma} \rangle \II(\dot{\gamma},\dot{\gamma}) + |\II(\dot{\gamma},\dot{\gamma})|^2_g = 0\,.$$
%Applying the Leibniz rule we have that $\langle n, \nabla_{\dot{\gamma}} \dot{\gamma} \rangle = - \II(\dot{\gamma},\dot{\gamma})$. Thus, we have that $\kappa = \pm \II(\dot{\gamma},\dot{\gamma})$.
%
%So differentiating both sides along $\dot{\gamma}$, we have that
%$$\nabla_{\dot{\gamma}} \kappa = \pm \bar{\nabla}_{\dot{\gamma}} \II(\dot{\gamma},\dot{\gamma})\,.$$
%But by hypothesis, the right hand side of this equation vanishes, so $\kappa$ is constant along $\gamma$, and hence $\gamma \hookrightarrow (M,g)$ has constant curvature.
%\end{proof}
%
%Alternatively, we may consider torsion-free curves in $M$ that can be defined by initial data on $\Sigma$:
%\begin{proposition}
%Let $\Sigma \hookrightarrow (M^d,g)$ be hypersurface embedding with $\II = 0$, and let $\gamma: I \rightarrow M$ be a curve with vanishing torsions. Further, suppose that $\gamma(0) \in \Sigma$ and $\dot{\gamma}(0), \ddot{\gamma}(0) \in T_{\gamma(0)} \Sigma$. Then, $\gamma$ lies entirely within $\Sigma$.
%\end{proposition}
%\begin{proof}
%Without loss of generality, we assume that $\gamma$ is arclength parametrized. Now, as $\gamma$ has vanishing torsion functions, it follows from the generalized Frenet--Serret formulas~\cite{} that
%\begin{align*}
%\nabla_{\dot{\gamma}} \left(\frac{1}{\kappa} \nabla_{\dot{\gamma}} \dot{\gamma} \right) = -\kappa \dot{\gamma}\,.
%\end{align*}
%Now consider instead the curve generated entirely within $\Sigma$ by the following differential equation, 
%$$\bar{\nabla}_{\dot{\bar{\gamma}}} \left(\frac{1}{\bar{\kappa}} \bar{\nabla}_{\dot{\bar{\gamma}}} \dot{\bar{\gamma}} \right) = -\kappa \dot{\bar{\gamma}}\,,$$
%with initial conditions given by $\bar{\gamma}(0) = \gamma(0)$, $\dot{\bar{\gamma}}(0) = \dot{\gamma}(0)$ and $\ddot{\bar{\gamma}}(0) = \ddot{\gamma}(0)$, and where $\kappa$ is a specified function $I \rightarrow \mathbb{R}$. Then, by the Gauss formula this curve also satisfies
%$$\nabla_{\dot{\bar{\gamma}}} \left(\frac{1}{\kappa} \nabla_{\dot{\bar{\gamma}}} \dot{\bar{\gamma}} \right) = -\kappa \dot{\bar{\gamma}}\,.$$
%But this is the same equation satisfied by $\dot{\gamma}$ and it has the same initial conditions. Because the solution is necessarily unique (as this is a system of ordinary differential equations), it follows that $\gamma = \bar{\gamma}$. But $\bar{\gamma}$ lies entirely within $\Sigma$, so the proposition holds.
%\end{proof}
%
%Thus, we see that not only does the second fundamental form characterize deviation of geodesics in $\Sigma$ from geodesics in $M$, it characterizes deviation of torsion-free curves in $\Sigma$ from torsion-free curves in $M$.
%
%Of course, the second fundamental form has a myriad of other geometric interpretations---however, we will leave those interpretations to other works. In this article, we focus (for the most part) on the geometric meaning of fundamental forms as they apply to curves.
%
%
%\section{Geometry of the third fundamental form} \label{3ff}
%
%Because the second fundamental form is, in a very explicit sense, a (Lie) derivative of the metric with respect to the normal vector, it is perhaps expected that the second fundamental form appears relating the connection of the bulk and the induced connection on the hypersurface. However, the third fundamental form is (to leading order) two Lie derivatives of the metric with respect to the normal vector. As such, we expect the third fundamental form to appear in equations that contain curvature tensors.
%
%One particularly important equation that depends on the curvature tensor is the Jacobi equation for a geodesic:
%$$\frac{D^2}{dt^2} J^c + J^a \dot{\gamma}^b \dot{\gamma}^d R_{ab}{}^c{}_{d} = 0\,.$$
%In this equation, $\frac{D}{dt}$ is the covariant derivative with respect to the arclength parameter along the geodesic, $\dot{\gamma}$ is the geodesic in question, and $J$ is the Jacobi field along that geodesic.
%
%Solutions to this equation are called Jacobi fields, and they are first-order difference vector fields between the geodesic in question and geodesics that are infinitesimally close to it. More precisely, if $\Gamma(s,t)$ is a smooth family of arclength-$t$ parametrized geodesics such that $\Gamma(0,t) = \gamma(t)$, then $J(t) = \frac{\partial \Gamma}{\partial s}|_{s=0}$. As such, when $J(t)$ has multiple zeros, we find that there are (to first order) at least two geodesics that pass through the same pair of points. Such a pair of points are called \textit{conjugate points} and there is a large body of literature investigating the existence of these points in Riemannian manifolds (for a few important examples, see~\cite{Morse1930,Schoenberg1932,morse1942,hopf1948,klingenberg1959,takagi1969,chicone1980,burago1994}). To that end, we provide new results on the relationship between conjugate points and the third fundamental form.
%
%
%\begin{proposition}
%Let $\Sigma^{d-1} \hookrightarrow (M^d,g)$ be a totally geodesic and flat hypersurface embedding with $\III$ negative definite, and let $(M,g)$ be complete. Then, for each point $p \in \Sigma$, there exists a conjugate point in $\Sigma$ to $p$.
%\end{proposition}
%\begin{proof}
%Assume $(\Sigma,\bar{g})$ is flat and totally geodesic, but $\III$ is negative definite. Now trivially, since $(\Sigma,\bar{g})$ is flat, it has no conjugate points. Furthermore, because $\II = 0$, we have that $R^\top_{abcd} = \bar{R}_{abcd}$. It follows from the definition of the third fundamental form and Equation~(\ref{gauss-equation}) that
%\begin{align*}
%\III_{ab} \= R_{nabn} \=& \top (n^c n^d R_{cabd}) \\
%\=& \top[(g^{cd} - \bar{g}^{cd}) R_{cabd}] \\
%\=& -Ric^\top_{ab} - \bar{g}^{cd} \bar{R}_{cabd} \\
%\=& -Ric^\top_{ab} + \overline{Ric}_{ab}\,.
%\end{align*}
%But $\overline{Ric} = 0$, so we have that since $\III$ is negative definite, $Ric^\top$ is positive definite. Therefore, it follows that for any geodesic $\gamma$ in $\Sigma$, we have that $Ric^\top(\dot{\gamma},\dot{\gamma}) > 0$, and hence $Ric(\dot{\gamma},\dot{\gamma}) > 0$. But then it follows~\cite{LiangZhan,MendZhou} that $\gamma$ has conjugate points in $M$. However, $\gamma$ lies within $\Sigma$, so any conjugate points lying along $\gamma$ necessarily lie within $\Sigma$. Now this is true for every geodesic starting at a point $p \in \Sigma$, and thus every point $p \in \Sigma$ has conjugate points.
%\end{proof}
%This result may be interpreted as stating that when $\III$ is negative definite and $\Sigma$ is flat, geodesics that start off near $\Sigma$ ``oscillate'' about $\Sigma$---to first order, of course. The contrapositive of this notion is given an explicit form in the following proposition.
%
%\begin{proposition}
%Let $\Sigma \hookrightarrow (M,g)$ be a totally geodesic hypersurface embedding and suppose that $\gamma$ is a geodesic in $\Sigma$ connecting points conjugate with respect to $(M,g)$. If $\III$ is positive definite, then those points are conjugate with respect to $(\Sigma,\bar{g})$.
%\end{proposition}
%
%\begin{proof}
%Because there are two conjugate points lying in $\Sigma$, there necessarily exists a non-vanishing Jacobi field along $\gamma$ that vanishes at both points. A priori, this Jacobi field satisfies
%$$\frac{D^2}{dt^2} J^c + J^a \dot{\gamma}^b \dot{\gamma}^d R_{ab}{}^c{}_{d} = 0\,.$$
%Contracting this equation with the unit conormal we find that the Jacobi field must satisfy
%$$\nabla_{\dot{\gamma}}^2 (n \csdot J) + R_{J \dot{\gamma} n \dot{\gamma}} = 0$$
%because $\frac{D}{dt} n = 0$ from the totally geodesic condition.	
%Now, decomposing the Jacobi field into a tangential and normal piece $J = q + An$, where $q$ lies within $T \Sigma$, and noting that $\II = 0$, we have that
%$$\nabla_{\dot{\gamma}}^2 A + R_{n \dot{\gamma} n \dot{\gamma}} A= 0\,.$$
%This follows because $R_{q\dot{\gamma} n \dot{\gamma}}^\top = 0$ from Equation~(\ref{cod}).
%Now because $\III(\dot{\gamma},\dot{\gamma}) = -R_{n \dot{\gamma} n \dot{\gamma}}$, we have that if $\III$ is positive definite, then $R_{n \dot{\gamma} n \dot{\gamma}}$ is strictly negative. It thus follows by standard ODE theory that if there exists $t_1 \neq t_2$ such that $A(t_1) = A(t_2) = 0$, then $A(t) = 0$.
%
%As there exist conjugate points on $\gamma$, it then follows that $A(t) = 0$ identically, and hence we have that $J \in T \Sigma$. As such, $J^a \dot{\gamma}^b \dot{\gamma}^d R_{ab}{}^c{}_d = J^a \dot{\gamma}^b \dot{\gamma}^d \bar{R}_{ab}{}^c{}_d$, because the normal component vanishes according to Equation~(\ref{cod}). Hence, we have that $J$, viewed as a vector field in $T \Sigma$, satisfies the Jacobi equation with respect to $(\Sigma,\bar{g})$. But because $J$ has two zeros (as it is a Jacobi field in $(M,g)$), it follows that those points are also conjugate with respect to $(\Sigma,\bar{g})$.
%\end{proof}
%
%Roughly speaking, this result implies that if there are conjugate points for a positive third fundamental form, those conjugate points must arise from geodesics lying in $\Sigma$ rather than from geodesics in the bulk.
%
%Finally, just as geodesics in $\Sigma$ are geodesics in $M$ if $\II = 0$, we see that Jacobi fields in $\Sigma$ are Jacobi fields in $M$ if $\II = \III = 0$:
%\begin{proposition}
%Let $\iota : \Sigma^{d-1} \hookrightarrow (M^d,g)$ be a totally geodesic hypersurface embedding. Furthermore, let $\bar{J}$ be a Jacobi field along a geodesic $\bar{\gamma}$ in $\Sigma$. It follows that if $\III = 0$, then $J := \iota_* \bar{J}$ is a Jacobi field along $\gamma := \iota_* \bar{\gamma}$.
%\end{proposition}
%
%\begin{proof}
%As $\bar{J}$ is a Jacobi field along $\bar{\gamma}$, we have that it satisfies the Jacobi equation, namely
%$$\frac{D^2}{dt^2} \bar{J}^c + \bar{J}^a \dot{\bar{\gamma}}^b \dot{\bar{\gamma}}^d \bar{R}_{ab}{}^c{}_{d} = 0\,.$$
%Now, as $\Sigma$ is totally geodesic, $\gamma := \iota_* \bar{\gamma}$ is a geodesic in $M$. As such, $\frac{D}{dt}$ is well-defined when viewed in either context.
%
%It now suffices to show that $\bar{R} = R|_{\Sigma}$. Note that from Equations~(\ref{gauss-formula}) and~(\ref{gauss-equation}), we have that $\bar{R} = R^\top$ and that $R_n^\top = 0$. Finally, the last independent component of the Riemann curvature, $R_{nabn}$, is equal to $\III$. Thus, we have that $R_{nabn}|_{\Sigma} = 0$, and hence $R_{abcd}|_{\Sigma} = \bar{R}_{abcd}$. The proposition follows.
%\end{proof}
%
%From the above results, we see that the third fundamental form is to Jacobi fields and conjugate points as the second fundamental form is to geodesics.
%
%\medskip
%
%Beyond just the relationship to conjugate points and Jacobi fields, though, the third fundamental form also shows up conspicuously in the Einstein equation:
%\begin{proposition}
%Suppose that $\Sigma \hookrightarrow (M,g)$ is totally geodesic, satisfies $\III = 0$, and $(M,g)$ is Einstein. Then, $(\Sigma,\bar{g})$ is Einstein.
%\end{proposition}
%\begin{proof}
%If $(M,g)$ is Einstein, then it satisfies
%$$Ric = \Lambda g\,,$$
%for some constant $\Lambda$. Pulling back this equation to $\Sigma$, we have that
%$$Ric^\top = \Lambda \bar{g}\,.$$
%However, because $\II = \III = 0$, it follows that $Ric^\top = \overline{Ric}$. Hence, we have that 
%$$\overline{Ric} = \Lambda \bar{g}\,,$$
%and hence $(\Sigma,\bar{g})$ is Einstein.
%
%
%
%%If $(M,g)$ is Einstein, then there exists a positive function $s \in C^\infty M$ such that
%%$$\nabla^2 s + P s = 0\,,$$
%%where $P := \tfrac{1}{d-2} (Ric - \tfrac{1}{2(d-1)} Sc\, g)$ is the Schouten tensor and $Sc = Ric_a^a$.
%%Thus, it follows that the tangential piece of this equation vanishes, namely
%%$$\bar{g}^c_{(a} \nabla^\top_{b)} \nabla_c s + P^\top_{ab} s = 0\,.$$
%%Now from the definition of $\III$, it follows that when $\II = 0$, we have that $P^\top = \bar{P}$. However, because $\Sigma \hookrightarrow (M,g)$ is totally geodesic, we have that $\II = 0$. Next, we consider the derivative terms.
%%\begin{align*}
%%\bar{g}^c_{(a} \nabla^\top_{b)} \nabla_c s \eqSig& \bar{g}^c_{(a} \nabla^\top_{b)} (\nabla^\top_c + n_c \nabla_n) s \\
%%\eqSig& \bar{g}^c_{(a} \nabla^\top_{b)}\nabla^\top_c  s \\
%%\eqSig& \bar{g}^c_{(a} \bar{\nabla}_{b)}\nabla^\top_c  s \\
%%\eqSig& \bar{\nabla}_{(a} \bar{\nabla}_{b)}  s \,.
%%\end{align*}
%%The second equality follows by the Leibniz rule, and the third and fourth equalities follow by the Gauss equation when the hypersurface is totally geodesic: $\bar{\nabla} = \nabla^\top$. Thus, we have that $\bar{s} := \iota^* s \in C^\infty \Sigma$ also satisfies
%%$$\bar{\nabla}^2 \bar{s} + \bar{P} \bar{s} \eqSig 0\,.$$
%%Hence, $(\Sigma,\bar{g})$ is conformally Einstein.
%\end{proof}
%
%
%
%
%
%
%\section{Geometry of the fourth fundamental form} \label{4ff}
%
%While the second and third fundamental forms have easily observed geometric interpretations in the form of geodesics and Jacobi fields, the geometric meaning of the fourth fundamental form is less obvious. However, by examining the geodesic equation in the context of warped product manifolds, the geometric meaning of the fourth fundamental form (and indeed, higher fundamental forms) becomes more clear.
%
%Consider the warped product $M := \Sigma  \times_{f} \mathbb{R}$ with metric $g = ds^2 + f(s) \bar{g}$. Now, defining $\Sigma \hookrightarrow (M,g)$ by $\Sigma = \{p \in M | s(p) = 0\}$, we see that $\II = \tfrac{1}{2} f'(0) \bar{g}$~\cite{oneil}. If $\II = 0$, then $\III = \tfrac{1}{2} f''(0) \bar{g}$. This pattern continues.
%
%Now, consider coordinates $(s, \vec{x})$ in a neighborhood of a point on $\Sigma$. Then, the geodesic equation in the $s$ direction is given by
%$$\frac{d^2 s}{dt^2} - \tfrac{1}{2} f'(s) \bar{g}_{ab} \frac{d x^a}{dt} \frac{dx^b}{dt} = 0\,.$$
%Thus, we see that $f(s)$ controls a potential function for a particle whose trajectory is given by $s(t)$, modulo motion in the horizontal (warped) direction. Thus, we see that when $\II = 0$, a trajectory that starts on $s = 0$ along $\Sigma$ stays on $\Sigma$---this is what it means for a hypersurface embedding to be totally geodesic. Alternatively, if $\II \neq 0$, then (in the vicinity of $\Sigma$) geodesics starting on $\Sigma$ in the direction of $\Sigma$ will move away from $\Sigma$ in a direction dictated by $\operatorname{sgn}f'(0)$ approximating a parabolic function.
%
%Similarly, if $\II = 0$ but $\III > 0$, we see that a geodesic starting on $\Sigma$ in a direction off $\Sigma$ will escape $\Sigma$ exponentially quickly. Similarly, if $\III < 0$ we see that geodesics will behave approximately like a simple harmonic oscillator as the potential is quadratic and increasing with distance from $\Sigma$. This gives intuition behind the results in Section~\ref{3ff}.
%
%Thus, we can understand the higher fundamental forms as controlling the potential energy function for geodesics. While in general, proving these sorts of results are out of reach, we can use standard ODE techniques to describe motion when $\Sigma = \mathbb{R}$. When $\II = \III = 0$, we characterize the trajectory of geodesics in the following proposition:
%
%
%%%%%\begin{theorem}
%%%%%Projected CC curves to a hypersurface?
%%%%%\end{theorem}
%%%%%\begin{proof}
%%%%%Consider an arclength parametrized curve $\gamma$ with constant curvature $\kappa$ and vanishing torsion that intersects a hypersurface $\Sigma \hookrightarrow (M,g)$ at a point $p \in \Sigma$---we shift the parametrization so that $\gamma(0) = p$. Now, with fixed $\kappa$,  $\gamma(0)$, $\dot{\gamma}(0)$ and $\ddot{\gamma}(0)$ determine $\gamma$ uniquely according to the Frenet-Serret formulas
%%%%%$$\dddot{\gamma} = -\kappa^2 \dot{\gamma}\,.$$
%%%%%Put another way, we have that
%%%%%$$\nabla_{\dot{\gamma}} \nabla_{\dot{\gamma}} \dot{\gamma} + \kappa^2 \dot{\gamma} = 0\,.$$
%%%%%Now we may equally well define a family of curves $\Gamma_s(t)$ solving the same equation with $\Gamma_s(0) = \gamma(0)$, but with 
%%%%%$$\dot{\Gamma}_s(0) =\tfrac{1}{\sqrt{1-s(n \cdot \dot{\gamma}(0))^2}} \left[(1-s) \dot{\gamma}(0) + s \bar{g}(\dot{\gamma}(0),\cdot) \right]$$
%%%%%and
%%%%%$$\ddot{\Gamma}_s(0) = \tfrac{\kappa^2}{\sqrt{\kappa^2-s(n \cdot \ddot{\gamma}(0))^2}} \left[(1-s) \ddot{\gamma}(0) + s \bar{g}(\ddot{\gamma}(0),\cdot) \right]\,.$$
%%%%%That is, $\Gamma_1(t)$ is the curve starting on $\Sigma$ with constant curvature with initial conditions given by the projection of the initial conditions of $\gamma$.
%%%%%
%%%%%Now we define a Jacobi field for these constant curvature curves by $ z:= \partial_s \Gamma_s(t)$. Note that $[z,\dot{\gamma}] = 0$. For brevity, denote $\dot{\gamma} =: T$. So, differentiating the Frenet-Serret frames with respect to $z$, we have that
%%%%%$$\nabla_z \nabla_{T} \nabla_{T} T + \kappa^2 \nabla_z T = 0\,.$$
%%%%%Rearranging this formula, we have that
%%%%%\begin{align*}
%%%%%0 =& \nabla_{T} \nabla_z \nabla_{T} T^a +  R_{z T}{}^a{}_{b} \nabla_{T} T^b + \kappa^2 \nabla_{T} z^a \\
%%%%%=& \nabla_{T} \nabla_{T} \nabla_{T} z^a + \nabla_{T} R^a{}_{TzT} + \kappa^2 \nabla_{T} z^a \\
%%%%%=& \nabla_{T}^3 z^a + (\nabla_T T^b) z^c T^d R^a{}_{bcd} + T^b (\nabla_T z^c) T^d R^a{}_{bcd} + T^b z^c (\nabla_T T^d) R^a{}_{bcd} + T^b z^c T^d \nabla_T R^a{}_{bcd} + \kappa^2 \nabla_T z^a \,.
%%%%%\end{align*}
%%%%%Now we are interested in the case where $\II = \III = 0$, and furthermore how $\IV$ impacts the normal component of $z$. In particular, we examine this equation when $t = 0$. In that case, contracting with $n$ and denoting $ a := n \cdot z$, we have that
%%%%%$$\nabla_T^3 a + a (T \cdot n) n^a T^b n^c T^d \nabla_n R^a{}_{bcd} + \kappa^2 \nabla_T a \stackrel{t=0}{=} 0\,.$$
%%%%%Or, writing this in terms of the fourth fundamental form, we have that
%%%%%$$\nabla_T^3 a + \kappa^2 \nabla_T a - (T \cdot n) \IV_{TT} a  \stackrel{t=0}{=} 0\,.$$
%%%%%\end{proof}
%
%%%%\begin{theorem}
%%%%Let $\Sigma \hookrightarrow (M^d,g)$ be a totally geodesic, separating, and oriented hypersurface embedding with $\III=0$. Furthermore, let $s$ be a unit defining function for $\Sigma$. Furthermore, let $\IV$ be positive definite. Then, for any geodesic that intersects $\Sigma$ at least once, we have that
%%%%$$\int_{-\infty}^{\infty} s(\gamma(t)) dt \geq 0\,.$$
%%%%Similarly, when $\IV$ is negative definite, we have that
%%%%$$\int_{-\infty}^{\infty} s(\gamma(t)) dt \leq 0\,.$$
%%%%\end{theorem}
%%%%IDEAS:
%%%%
%%%%Consider crossing numbers in two dimensions. If a geodesic crosses a non-zero and even number of times, then $\int s$ is sign-definite. If it crosses zero times, then $\int s$ is sign definite. The only non-$\Sigma$ geodesics that might have vanishing $\int s$ are those geodesics that cross an odd number of times (note that geodesics can't just ``kiss'' $\Sigma$ because $\Sigma$ is totally geodesic so all such geodesics would stay on $\Sigma$). Alternatively, there may be an infinite number of crossings. We examine those.
%%%%
%%%%
%%%%Suppose now that $\gamma$ crosses $\Sigma$ an odd number of times, say one time, say at the origin, and furthermore say that $\int s \leq 0$. We begin by considering the case where $\int s = 0$. Such a geodesic is generated at the origin by a unit tangent vector, say $\dot{\gamma}(0)$, specified by an angle $\theta(0)$. Now either $\gamma$ is a critical point with respect to variations of $\theta(0)$ or it is not.
%%%%
%%%%PLAN:
%%%%
%%%%Start with geodesic curve defined as $s = 0$, and a geodesic $\Gamma_{\theta}(t)$ shooting off that curve at $(s,x) = (0,0)$ with initial data given by $\dot{\Gamma}_{\theta}(0) = \cos \theta \partial_x + \sin \theta \partial_s$. Now without loss of generality suppose that $0 < \theta < \pi/2$. Then, there exists some sufficiently small $\epsilon>0$ such that $\Gamma_{\theta}$ intersects $s = \epsilon$. We would like to generate a Jacobi field for geodesics that start on the line $s = \epsilon$ and leaves....
%%%%
%%%%%Suppose that $\gamma$ crosses $\Sigma$ an odd number of times, say one time, say at the origin. Then, there is some small neighborhood around the origin that we may foliate with $\Sigma_r$ such that $\gamma$ only crosses each leaf once. Now because $\II = \III = 0$, we have that $\II_r \sim \mathcal{O}(r^2)$ and $\III_r \sim \mathcal{O}(r)$. However, because $\IV > 0$, there exists $r > 0$ such that $\II_r > 0$ and $\III_r > 0$. 
%%%%
%%%%
%%%%\bigskip
%%%%
%%%%Can we develop a heuristic argument using the volume form? In two dimensions, when $\II = \III = 0$ and $\IV > 0$, we have that the volume form shrinks as you go to $s<0$ and grows as you go to $s > 0$, at least near $s = 0$, with a critical point at $s = 0$.
%%%%
%%%%Are geodesics attracted to regions of largest volume form?
%%%%
%%%%
%%%%\bigskip
%%%%
%%%%Let's consider the simplest case where we have a metric of the form
%%%%$$g = ds^2 + (1+s^3)dx^2$$
%%%%on the piece of $\mathbb{R}^2$ given by $s > -1$.
%%%%
%%%%Note that the hypersurface given by $s = 0$ has $\II = \III = 0$ and $\IV > 0$. The geodesic equation can be reduced to
%%%%$$\frac{d^2 s}{dt^2} - \tfrac{3}{2} s(t)^2 \left(\frac{dx}{dt}\right)^2 = 0\,.$$
%%%%If we assume that $x(t)$ is invertible, we can then write this differential equation as
%%%%$$\frac{d^2 s}{dx^2} - \tfrac{3}{2} s(x)^2 = 0\,.$$
%%%%Now, multiplying both sides by $2s'$ and integrating, we find the first integral:
%%%%$$(s')^2 = s^3 + c\,.$$
%%%%By translation symmetry, let us suppose a geodesic starts at the origin. Then, we find that
%%%%$$c = (s'(0))^2\,.$$
%%%%Also note that there exists a unique location where $s' = 0$---namely, when $s^3 = -c = -(s'(0))^2$. Indeed, the fact that there exists a single extremum means that solutions cannot have vanishing total $s$. Furthermore $s$ passes through the origin and its minimum is at $s<0$, we must have that $\int_{\gamma} s > 0$ for every geodesic passing through the origin for which $x(t)$ is invertible.
%%%%
%%%%\bigskip
%%%%
%%%%$${}$$
%%%%
%%%%\bigskip
%%%%
%%%%We now consider the more general case, where
%%%%$$g = ds^2 + (1+s^3 G(s)) dx^2\,,$$
%%%%for any smooth $s^3G(s) > -1$ with $G(0) > 0$ on $\mathbb{R}^2$. Then, again assuming that $x(t)$ is invertible, we find that the geodesic equations reduce to
%%%%$$s'' - \tfrac{1}{2} s^2 (3 G(s) + sG'(s)) = 0\,.$$
%%%%Note that this can be written as
%%%%$$s'' - \tfrac{1}{2} (s^3 G)' = 0\,.$$
%%%%
%%%%Multiplying through by $2s'$ and integrating, we again have a first integral:
%%%%$$(s')^2 = s^3 G(s)  + c\,,$$
%%%%where $c$ is again a constant of integration. Because $G(s)$ is smooth and $s(0) = 0$, we therefore (again) have that $c = (s'(0))^2$. Now, we may solve for when $s' = 0$. Now, extrema of $s$ occur when
%%%%$$ s^3G(s) + (s'(0))^2 = 0\,.$$
%%%%
%%%%
%%%%CONJECTURE: When $s^3 G(s) + (s'(0))^2 = 0$ has a single real root that's positive, then there's a counterexample---so we need to show that this is impossible, or exclude this case for another reason. When there is a single real root that's negative, the result holds, and when there are at least two real roots, then the system is (probably?) oscillatory and the result is probably still true.
%%%%
%%%%%
%%%%%COUNTEREXAMPLE: $G(s) = 1-10s$, $s(0) = 0$ and $s'(0) = \tfrac{1}{2}$. This example has $s$ diverging off to negative infinity on either side of $0$ and has $t(x)$ not defined outside of a small range $x \in (-1.563,4.069)$.
%%%%%
%%%%%CONJECTURE: What we learn is that when $s^3 + s^4 G(s) = -(s'(0))^2$ has a single real root and it's positive, then we have counterexamples. When there is a single real root that's negative, then the result holds, and when there are two real roots, then the system is (probably?) oscillatory and the result is probably still true.
%%%%
%%%%In fact, there can be more than two real roots. However, only the largest negative real root, call it $s^-$, can be realized. To see this, suppose not: then, there exists a geodesic that arrives at $s^-$ from above and becomes tangent to the curve $s = s^-$ at that point. Because we can backtrack that geodesic and it leaves the curve, $s = s^-$ is not itself a geodesic. But because (without loss of generality) it is above $s = s^-$ when the curve is left of the intersection point, and because the curvature of $s^-$ is constant, the concavity must be preserved. Hence, we have that the geodesic must also leave $s = s^-$ above to the right. Hence, any approach to an extremum that is below the $x$ axis must be repulsive.
%%%%
%%%%The same argument applies for multiple positive roots. Indeed, by isotropy in the $x$ direction, if there are at least two real roots and at least one is negative and at least one is positive, then the geodesic must oscillate about $s = 0$.
%%%%
%%%%So we now have three cases: if there is one negative root, then $\int s > 0$. If there is one positive root, then $\int s < 0$. If there is at least one negative root and at least one positive root, then $s$ oscillates about $s$. We now need only show that the oscillation spends more time at $s > 0$ than at $s < 0$.
%%%%
%%%%To handle the oscillatory case, perhaps we should consider the case where $s'(0) = \epsilon \ll 1$.
%%%%
%%%%%So we consider the case where $s'(0) = \epsilon \ll 1$. For sufficiently small $\epsilon$, $G(s)$ is negligible (FORMALIZE THIS). In this case, we find that there always exists a negative root, and furthermore it coincides with the root for the same geodesic in a space where $G(s)$. Even better, this root is always the largest negative root, and it is quite close to zero.
%%%%
%%%%
%%%%\color{blue}The existence of solutions with just a single positive root might not be allowed because that would mean that $g$ is degenerate somewhere, keep this in mind. \color{black}
%%%%
%%%%Suppose there's just a single positive root, call it $s^+$. Now recall that
%%%%$$ s^3 G(s) = (s')^2 - c\,.$$
%%%%But, the manifold is only defined for $s^3 G(s) > -1$, so we must have that $(s')^2 - c > -1$, or equivalently
%%%%$$(s')^2 > c - 1\,.$$
%%%%Now we need to investigate the convexity of the geodesic to derive a contradiction.
%%%%
%%%%
%%%%\color{red}
%%%%CONTRADICTION:
%%%%
%%%%Let
%%%%$$g = ds^2 + (1+s^3 [1-\tfrac{3}{2}s - \tfrac{7}{4}s^2 + \tfrac{7}{4} s^3]) dx^2\,.$$
%%%%Note that $g$ is defined for all of $\mathbb{R}$, and $\IV >0$ for $s = 0$. Numerical solutions of the geodesic with $\gamma(0) = (0,0)$ and  $s'(0) = \sin \tfrac{\pi}{2}$ and $x'(0) = \cos \tfrac{\pi}{6}$ contradicts the conjecture, as $\int s = -\infty$.
%%%%
%%%%\color{black}
%%%%
%%%%$${}$$
%%%%
%%%%New conjecture: For any choice of $G(s)$, for geodesics that start out "sufficiently close" to $s = 0$, they spend more time at $s > 0$ than at $s < 0$.
%
%\begin{proposition}
%Let $(\mathbb{R}^2,g)$ be a warped product with a global coordinate chart $(s,x)$, and let $\Sigma \hookrightarrow (\mathbb{R}^2,g)$ be a curve embedding defined by $\Sigma = \{p \in \mathbb{R} | s(p) = 0\}$. Furthermore, let $\Sigma$ have $\II = \III = 0$ and $\IV > 0$. Then, there exists $\epsilon_0 > 0$ such that, for any unit-speed geodesic $\gamma$ with invertible $x(t)$ starting at $(0,0)$ with $0< |s'(0)|  = \epsilon < \epsilon_0$, we have that either $\gamma$ is periodic for some period $T$ and $\int_0^T s(\gamma(x)) dx > 0$, or there is only a finite interval such that $s(\gamma(x))<0$.
%
%
%%Let $g = ds^2 + (1+s^3 G(s))dx^2$ for some smooth $G(s)$ such that $G(0) > 0$ and $g$ is everywhere nondegenerate on $\mathbb{R}$. Then, there exists $\epsilon_0 > 0$ such that, for any unit-speed geodesic $\gamma$ with invertible $x(t)$ starting at $(0,0)$ with $0< |s'(0)|  = \epsilon < \epsilon_0$, we have that either $\gamma$ is periodic for some period $T$ and $\int_0^T s(\gamma(x)) dx > 0$, or there is only a finite interval such that $s(\gamma(x))<0$.
%\end{proposition}
%
%\begin{proof}
%Because $(\mathbb{R}^2,g)$ has a global coordinate chart and is a warped product, we may write $g = ds^2 + f(s) dx^2$. Furthermore, because $\II = \III = 0$ and $\IV > 0$, we have that $f(s) = 1 + s^3 G(s)$ where $G(0) > 0$. Now consider a unit-speed geodesic $\gamma$ with invertible $x(t)$ as described in the hypothesis of the proposition.
%
%Because $x(t)$ is invertible, we may write the geodesic equation for this system as
%$$s'' - \frac{1}{2} \frac{d}{ds}(s^3 G) = 0\,.$$
%Multiplying by $2s'$ and integrating once, our differential equation reduces to
%$$(s')^2 - s^3 G(s) = c\,,$$
%where $c$ is an integration constant. Note that because $s(0) = 0$, we have that $c = \epsilon^2$.
%
%%Now, because $s^3 G(s)$ is greater than $-1$, we have that for sufficiently small $s<0$, $G(s) < 0$. \color{red} IS THIS TRUE \color{black} Similarly, for sufficiently large $s>0$, we have that $G(s) > 0$.\color{red} Is this true? What about functions that are positive on the left but exponentially supressed, and functions that are negative on the right but also exponentially suppressed?\color{black}
%
%
%We study the behavior of $-s^3 G(s)$. First, consider the case where $s<0$. Because $G(0)$ is positive and continuous, we have that there exists some (small) $s_0<0$ where $G(s_0) > 0$, and hence $-s_0^3 G(s_0)>0$, and where $-s^3 G(s)$ is strictly decreasing for all $s_0 < s < 0$. Going forward, we fix such a value of $s_0$. Then, define $\delta_0^2 = -s_0^3 G(s_0)$, and define $\epsilon_0 = \delta_0/2$. 
%%Then, define $\delta_0^2 = -s_0^3 G(s_0)$, so that, if $s'(0) = \delta_0$, we have that at $s = s_0$, $s' = 0$. It follows that if $|s'(0)| < \delta_0$, then $s(x) > s_0$ for all $x$. So we define $\epsilon_0 = \delta_0/2$.
%
%Now, we consider a geodesic with $|s'(0)| = \epsilon$ for any $0<\epsilon < \epsilon_0$. Then, we may define $s_{\epsilon}$ to be the largest negative solution to the equation $-s^3 G(s) = \epsilon^2$, so that $s' = 0$ at $s_{\epsilon}$ for the prescribed initial data. Notably, because $\epsilon < \delta_0$, it follows that $s_{\epsilon} > s_0$. As such, we have that $\frac{d}{ds}(s^3 G)|_{s_{\epsilon}} > 0$. It follows that $s''|_{s_{\epsilon}} > 0$, and so $s_{\epsilon}$ is necessarily a turnaround point for $s(x)$ rather than a fixed point.
%
%Given this behavior, there exist a value of $x$, call it $x_0$, such that $s(x_0)=0$ and $s(x)<0$ for all $x_0 < x < 0$. Now for the sake of the theorem, it will be useful to compute $\int_{x_0}^{0} s(x) dx$. Note that, by construction, $s(\tfrac{x_0}{2}) = s_{\epsilon}$, so we may instead express this integral as
%$$2 \int_{\tfrac{x_0}{2}}^{0} s(x) dx = 2 \int_{s_{\epsilon}}^0 \frac{s \,ds}{s'} = \int_{s_{\epsilon}}^0 \frac{2s \,ds}{\sqrt{\epsilon^2 + s^3 G(s)}}\,.$$
%Now, we wish to show that this integral goes to zero in the $\epsilon \rightarrow 0$ limit. To do so, we will first need to produce a series expansion of this integral. To do so, first note that 
%$$\frac{2s \sqrt{s-s_{\epsilon}}}{\sqrt{\epsilon^2 + s^3 G(s)}}$$
%is smooth and nonzero as $s \rightarrow s_{\epsilon}$. This follows because $\frac{d}{ds}(s^3 G)|_{s_{\epsilon}} > 0$ and $s^3 G(s)$ is smooth, so $\epsilon^2 + s^3 G(s)$ is well approximated by a linear function around $s_{\epsilon}$. Indeed, for any nonzero $\epsilon$, this implies that
%$$\lim_{s \rightarrow s_{\epsilon}} \frac{\sqrt{s-s_{\epsilon}}}{\sqrt{\epsilon^2 + s^3 G(s)}} = \alpha \neq 0\,.$$
%By L'Hopital's rule, we therefore have that
%$$\alpha = \lim_{s \rightarrow s_{\epsilon}} \frac{\sqrt{\epsilon^2 + s^3 G(s)}}{\sqrt{s-s_{\epsilon}} (s^3 G(s))'} = \tfrac{1}{\alpha} \lim_{s \rightarrow s_{\epsilon}} \frac{1}{(s^3 G(s))'}\,.$$
%Hence, because $(s_{\epsilon}^3 G(s_{\epsilon}))' > 0$, we have that
%$$\alpha = \frac{1}{\sqrt{(s_{\epsilon}^3 G(s_{\epsilon}))'}}\,.$$
%So, it follows that
%$$\frac{2s \sqrt{s-s_{\epsilon}}}{\sqrt{\epsilon^2 + s^3 G(s)}} = \frac{2s_{\epsilon}}{\sqrt{(s_{\epsilon}^3 G(s_{\epsilon}))'}} + \mathcal{O}(s-s_{\epsilon})\,.$$
%
%Therefore, we may write the integral as
%\begin{align*}
%\int_{s_{\epsilon}}^0 \frac{2s \, ds}{\sqrt{\epsilon^2 + s^3 G(s)}} &= \int_{s_{\epsilon}}^0 \left(\frac{2 s_{\epsilon}}{\sqrt{(s_{\epsilon}^3 G(s_{\epsilon}))'}} \frac{1}{\sqrt{s - s_{\epsilon}}} + \mathcal{O}(\sqrt{s-s_{\epsilon}}) \right) ds \\
%&= \frac{4 s_{\epsilon} \sqrt{-s_{\epsilon}}}{\sqrt{(s_{\epsilon}^3 G(s_{\epsilon}))'}} + \mathcal{O}((-s_{\epsilon})^{3/2})\,.
%\end{align*}
%Because $s_{\epsilon} \rightarrow 0$ as $\epsilon \rightarrow 0^+$, at leading order we have that
%$$\int_{s_{\epsilon}}^0 \frac{2s \, ds}{\sqrt{\epsilon^2 + s^3 G(s)}} = -\frac{4 \sqrt{-3 s_{\epsilon}}}{3 \sqrt{G(0)}} + \mathcal{O}((-s_{\epsilon})^{3/2}) \stackrel{\epsilon \rightarrow 0}{\rightarrow} 0\,.$$
%
%Summarizing the above, we have that
%$$\lim_{\epsilon \rightarrow 0^+} \int_{x_0}^0 s(x) dx = 0\,.$$
%
%\bigskip
%
%We now wish to consider the time spent at $s>0$. Again because $G(0)$ is positive and $G(s)$ is continuous, we have that there exists some small $s_1$ where $G(s_1) > 0$ and hence $-s_1^3 G(s_1) < 0$.
%
%Now we have two cases: either $G(s) \geq 0$ for all $s > 0$ or there exists some $s_2$ where $G(s_2) < 0$. If $G(s) \geq 0$ for all $s >0$, then $-s^3 G(s) < 0$ for all $s >0$, and so $s(x) > 0$ for all $x > 0$. Similarly, we find that $s(x)>0$ for all $x<x_0$. Thus, there is only a finite interval $(x_0,0)$ where $s(x) < 0$. Thus either the theorem holds or we have that there exists some $s_2 >0$ such that $G(s_2) < 0$.
%
%In that case, for the reasons given in the case where $s<0$, there exists sufficiently small $\epsilon_0$ such that, for all geodesics with $s'(0) = \epsilon < \epsilon_0$, there is a turnaround point $\hat{s}_{\epsilon} > s_1$. Thus, for any such geodesic, because the differential equation is autonomous, when $s(x) = 0$, we necessarily have that $s'(x) = \pm s'(0)$, and hence the motion truly is periodic. Thus, it suffices to show that the integral of $s(x)$ over a period is positive. Note from the $s<0$ considerations, we may choose the negative contribution to be arbitrarily small. Hence it will suffice to show that regardless of how small $\epsilon_0$ is chosen to be, the positive contribution to this integral does not approach zero. That is, we would like to show that
%$$\lim_{\epsilon \rightarrow 0^+} \int_0^{\hat{s}_{\epsilon}} \frac{s \, ds}{\sqrt{\epsilon^2 + s^3 G(s)}} \not \rightarrow 0\,.$$
%But $\hat{s}_{\epsilon} > s_1 > 0$ for all $\epsilon > 0$, and hence in the limit does not approach zero. Furthermore, the integrand is strictly positive and finite over the entire interval. Thus, in the limit where $\epsilon \rightarrow 0$, the integral must be positive. The theorem follows.
%
%\end{proof}
%
%The above result does not directly generalize to the case for an arbitrary warped product manifold with the topology of $\mathbb{R} \times \mathbb{R}^{d-1}$. This is because, when the base manifold in the warped product is not flat, we may not be able to factorize the geodesic equation in the $s$ direction to remove dependence on $\frac{d \vec{x}}{dt}$. Nonetheless, we suspect that a similar result is obtainable, hence we conjecture the following:
%\begin{conjecture}
%Let $\Sigma$ be diffeomorphic to $\mathbb{R}^{d-1}$ endowed with a metric $\bar{g}$, and let $M = \mathbb{R} \times_{f} \Sigma$ be a warped product manifold with a hypersurface embedding $\Sigma_0 \hookrightarrow (M,g)$ with $\II = \III = 0$ and $\IV$ is positive definite. Then, there exists $\epsilon_0>0$ such that, for any unit-speed geodesic $\gamma$ determined by $\gamma(0) = p \in \Sigma_0$ and $0 < |n \cdot \gamma(0)| = \epsilon < \epsilon_0$, the geodesic is either periodic for some period $T$ and $\int_0^T s(\gamma(t)) dt$ or there is only a finite interval such that $s(\gamma(t)) < 0$.
%\end{conjecture}
%
%More speculatively, we suspect that a similar result holds for hypersurface embeddings that are not warped product manifolds. However, this is far from obvious.
%
%
%
%%
%%\begin{proof}
%%\color{red}
%%When this is the case, we may express the metric in a small neighborhood of a point $p \in \Sigma$ as
%%$$g = ds^2 + \bar{g} + s^3 h\,,$$
%%where $h$ is a function of all coordinates, but is otherwise smooth in a small neighborhood.
%%
%%Now consider a geodesic $\gamma$ starting at $p \in \Sigma$, pointing off $\Sigma$. Now we may define a family of hypersurfaces $\Sigma_{s_0} := \{p \in M | s(p) = s_0\}$. Assume without loss of generality that for some positive parameter, $\gamma$ passes through $\Sigma_{t}$ for some $t > 0$ at a point $q \in \Sigma$. Thus, $\gamma$ is equally well defined by $q$ and $\dot{\gamma}(q)$.
%%
%%Note that $\Sigma_t$ is \textit{not} totally geodesic, but has $\II \propto t^2$. Similarly, $\III \propto t$. Now, consider a geodesic $\rho_t$ that starts at $q$ and has initial unit velocity equal to the projection of $\dot{\gamma}(p)$ to $\Sigma_t$. Now while this geodesic does not stay on $\Sigma_t$, its deviation from $\Sigma_t$ is of order $t^2$. Now, define $\Gamma(r,T)$ as the family of geodesics linearly interpolating between $\rho_t$ and $\gamma$. In that case, we have a Jacobi field $J_t$ satisfying
%%$$\frac{D^2}{dT^2} J_T^c + J^a_t \dot{\rho}_t^b \dot{\rho}_t^d R_{ab}{}^c{}_d = 0\,.$$
%%Contracting with the normal vector field and noting that $\frac{Dn}{dT} \stackrel{q}{=} \II(\dot{\rho}_t,\cdot) \propto t^2$, we have that, at $q$,
%%$$\frac{D^2}{dT^2} (n \cdot J_t) + R_{J_t \dot{\rho}_t n \dot{\rho}_t} \stackrel{q}{=} \mathcal{O}(t^2)\,.$$
%%Expressing $J_t \stackrel{q}{=} u + An$, we have that
%%$$\frac{D^2}{dT^2} A +  R_{n \dot{\rho}_t n \dot{\rho}_t} A \stackrel{q}{=} \mathcal{O}(t^2)\,.$$
%%
%%Note that $R_{n \dot{\rho}_t n \dot{\rho}_t} = -\III_t(\dot{\rho}_t, \dot{\rho}_t)$, so
%%$$\frac{D^2}{dT^2} A - \III_t(\dot{\rho}_t, \dot{\rho}_t) A \stackrel{q}{=} \mathcal{O}(t^2)\,.$$
%% Now, if $\IV$ is positive definite, then $\III_t$ is positive definite for $t > 0$ and negative definite for $t < 0$. So when $\IV$ is positive definite, we have that the above equation (at least initially) has growing solutions for $A$ for small $t$. On the other hand, when $\IV$ is negative definite, the tendency of $A$ is to get smaller.
%%
%%We thus have that when $\IV$ is negative definite, for large parameter values geodesics tend toward $s < 0$.
%%\color{black}
%%\end{proof}
%
%
%\section{Fundamental form obstructions} \label{obstructions}
%In this section, we discuss higher fundamental forms more generally. For example, as mentioned in the introduction, depending on the definition used, fundamental forms are always expressible as powers of the seocnd fundamental form when the bulk manifold is flat. Using our definition of the fundamental forms, they trivially vanish. However, the result can be re-expressed as describing the behavior of normal derivatives of $\nabla n$, which would (typically) have terms including higher fundamental forms.
%\begin{proposition}
%Let $\Sigma \hookrightarrow (M^d,g)$ be an embedded hypersurface and let $(M,g)$ be a spaceform. Furthermore, let $n = ds$ be a unit conormal for $s$ a defining function of $\Sigma$. Then, for every positive integer $k$, $\nabla_n^k \nabla n$ is expressible as a sum of powers of the second fundamental form.
%\end{proposition}
%\begin{proof}
%We perform a simple calculation:
%\begin{align*}
%\nabla_n^k \nabla_a n_b =& \nabla_n^{k-1} n^c \nabla_a \nabla_c n_b + \nabla_n^{k-1} n^c R_{cabd} n^d \\
%=& \nabla_n^{k-1} \nabla_a \nabla_n n_b - \nabla_n^{k-1} [(\nabla_a n^c)(\nabla_c n_b)] + \nabla_n^{k-1} R_{nabn} \\
%\eqSig& \nabla_n^{k-1} R_{nabn} + \text{ products of lower order terms.}
%\end{align*}
%Because $(M,g)$ is a spaceform, it follows by induction that all such normal derivatives are expressible as products of the second fundamental form $\II$.
%%Next, observe that
%%\begin{align*}
%%\nabla_n^{k-3} R_{nbcn} \eqSig& n^a n^d n^e \nabla_e \nabla^{k-4}_n R_{abcd} + \text{ lower order terms containing $\nabla n$}\\
%%\eqSig& n^d (g^{ae} - \bar{g}^{ae}) \nabla_e \nabla^{k-4}_n R_{abcd} + \text{ lower order terms containing $\nabla n$} \\
%%\eqSig& n^d \nabla^a \nabla^{k-4}_n R_{abcd} - n^d \nabla^{\top a} \nabla^{k-4}_n R_{abcd} + \text{ lower order terms containing $\nabla n$} \\
%%\eqSig& n^d \nabla^{k-4}_n \nabla^a  R_{abcd} - n^d \nabla^{\top a} \nabla^{k-4}_n R_{abcd} + \text{ lower order terms containing $\nabla n$} \\
%%\eqSig& -n^d \nabla^{k-4}_n (\nabla_d R_{ab}{}^a{}_c + \nabla_c R_{abd}{}^a) - n^d \nabla^{\top a} \nabla^{k-4}_n R_{abcd} + \text{ lower order terms containing $\nabla n$} \\
%%\eqSig&- n^d \nabla^{\top a} \nabla^{k-4}_n R_{abcd} + \text{ lower order terms containing $\nabla n$} \\
%%\end{align*}
%
%\end{proof}
%
%\begin{remark} \label{weights2}
%More generally, when $n = ds$ is a unit conormal for $s$ a defining function of the hypersurface, a similar calculation as that above shows that $\iota^* \nabla_n^k \nabla_a n_b \= \FF{k+2} + \text{subleading}$, where the subleading terms are contractions of lower fundamental forms. More specifically, assigning a weight of $3-m$ to $\FF{m}$ and a weight of $-2$ for every contraction, we find that each summand in an expression for $\iota^* \nabla_n^k \nabla n$ has weight $1-k$.
%\end{remark}
%
%Furthermore, the fundamental forms can be viewed as an obstruction to a manifold with embedded hypersurface being isomorphic to a product manifold:
%\begin{proposition} \label{prod-vanishingFFs}
%Let $\iota : \Sigma \hookrightarrow (M,g)$ be a smooth hypersurface embedding. Then, $\FF{k} = 0$ for every $k \geq 2$ if and only if in a collar neighborhood around $\Sigma$, $(M,g)$ is locally isometric to a product manifold $((-\epsilon,\epsilon) \times \Sigma, ds^2 + \iota^* g)$ for some $\epsilon > 0$.
%\end{proposition}
%\begin{proof}
%In a collar neighborhood of $\Sigma$, we may always choose coordinates $(s,\vec{x})$ such that $g = ds^2 + \bar{g}(s,\vec{x})$ and $\bar{g}$ is a metric on slices of constant $s$. Now, because the embedding is smooth, we may Taylor series expand $\bar{g}(s,\vec{x})$ about $s = 0$, so that 
%$$g = ds^2 + \bar{g}^{(0)}(\vec{x}) + s \bar{g}^{(1)}(\vec{x}) + \cdots\,.$$
%
%Now suppose that, for each $k \geq 2$, we have that $\FF{k} = 0$. Then, it follows that $\iota^* \mathcal{L}_{\partial_s}^k g \= 0$ \edz{need to prove this... maybe put it earlier?} for all $k \geq 1$. It therefore follows that $\bar{g}^{(k)}(\vec{x}) =0$ for each $k \geq 1$. Hence, we have that
%$$g = ds^2 + \bar{g}^{(0)}(\vec{x}) = ds^2 + \iota^* g\,.$$
%Thus, the claim of isometry follows.
%
%On the other hand, if $g = ds^2 + \bar{g}(\vec{x})$, then for the same reason, it follows that $\II = \tfrac{1}{2} \iota^* \mathcal{L}_n g = 0$. Furthermore, because the manifold is a product and the $\mathbb{R}$ prodand is flat, it follows that the curvature tensor $R = \bar{R}$. Thus, it follows that $R_{nabn} = \III = 0$. Finally, because the manifold is a product, the connection satisfies $\nabla_a = n_a \partial_s + \bar{\nabla}_a$. But $\partial_s R = 0$, so it follows that all higher fundamental forms vanish as well.
%\end{proof}
%
%\begin{remark}
%It trivially follows that, if each fundamental form vanishes, then in a collar neighborhood of $\Sigma \hookrightarrow (M,g)$, all geodesics are determined by geodesics on $\Sigma$. In particular, every geodesic on $M$ is given by $t \mapsto (\alpha t, \bar{\gamma}(t)) \in M$, where $\bar{\gamma}$\edz{do we need to make this arclength parametrization?} is a geodesic on $\Sigma$ and $\alpha$ is some constant. \edz{make this bidirectional}
%\end{remark}
%
%
%
%%
%%\color{red}
%%Willmore surfaces?
%%\color{black}
%%
%%LOOK FOR OTHER INTERESTING CLASSES OF TOTALLY GEODESIC HYPERSURFACES... zero locus of Killing vector fields? What else? Geodesic hypersurfaces in spaceforms have non-vanishing third fundamental form but vanishing second fundamental form?
%%
%%Redo calculation with conformally Einstein---right everything in terms of $\IIo$, $H$, $\mathring{P}$, and $J$. Also, what about Bach-flat to Bach-flat?
%%
%%Conformally Einstein:
%%$$\nabla_{a} \nabla_{b} s + P_{ab} s + g_{ab} \rho = 0\,.$$
%%Pullback to the hypersurface with $\II = \III = 0$:
%%\begin{align*}
%%0 \=& \bar{\nabla}_a \bar{\nabla}_b \bar{s} + P^\top_{ab} \bar{s} + \bar{\rho} \bar{g}_{ab} \\
%%\=& \bar{\nabla}_a \bar{\nabla}_b \bar{s} + \tfrac{1}{d-2}\left(Ric^\top_{ab} - \frac{Sc}{2(d-1)} \bar{g}\right) \bar{s}  + \bar{\rho} \bar{g}_{ab} \\
%%\=& \bar{\nabla}_a \bar{\nabla}_b \bar{s} + \tfrac{1}{d-2} \left(\bar{Ric}_{ab} - \frac{2 Ric_{nn} + \bar{Sc}}{2(d-1)} \bar{g} \right) \bar{s}   + \bar{\rho} \bar{g}_{ab}\\
%%\=&\bar{\nabla}_a \bar{\nabla}_b \bar{s} + \tfrac{1}{d-2} \left((d-3) \bar{P}_{ab} + (\bar{J} - \frac{\bar{Sc}}{2(d-1)}) \bar{g} \right) \bar{s}  + \bar{\rho} \bar{g}_{ab} \\
%%\=&\bar{\nabla}_a \bar{\nabla}_b \bar{s} + \tfrac{1}{d-2} \left((d-3) \bar{P}_{ab} +\tfrac{1}{2(d-1)(d-2)}  \bar{Sc} \bar{g} \right) \bar{s}  + \bar{\rho} \bar{g}_{ab} \\
%%\=&\bar{\nabla}_a \bar{\nabla}_b \bar{s} + \tfrac{1}{d-2} \left((d-3) \bar{P}_{ab} +\tfrac{1}{d-1}  \bar{J} \bar{g} \right) \bar{s}  + \bar{\rho} \bar{g}_{ab} \\
%%\end{align*}
%%Note that because $\III = 0$, we have that $Ric_{nn} = 0$. Taking the trace-free piece, we see that the conformally-Einstein condition is not satisfied.
%%
%%
%%When does symmetric space pullback to symmetric space?
%%
%
%
%
%
%\section{Warped Product Manifolds and Fundamental Forms}
%Warped product manifolds with a hypersurface as the base manifold are of particular interest and can be characterized using higher fundamental forms. In fact, all even fundamental forms vanish for warped product manifolds:
%\begin{proposition} \label{vanishing-evens2}
%Let $(M,g)$ be a warped product $\mathbb{R} \times_f \Sigma$ for some $f : \Sigma \rightarrow \mathbb{R}$. Furthermore, let $\Sigma \hookrightarrow (M,g)$ be a hypersurface embedding. Then, $\FF{k} = 0$ for all $k = 2n$ for positive integers $n$.
%\end{proposition}
%\begin{proof}
%As $(M,g)$ is a warped product, in a coordinate system adapted to the warped product $(t, x^i)$, the metric is expressible as
%$$g = f(x^1, \ldots, x^{d-1}) dt^2 + h_{ij}(x^1, \ldots, x^{d-1}) dx^i dx^j\,.$$
%%Now, we may explicitly compute the second fundamental form in this metric, noting that $n = \sqrt{f} dt$:
%%$$\iota^* \nabla_a n_b = \iota^* [(dt)_b \partial_a \sqrt{f}] - \sqrt{f} \iota^* \Gamma^t_{ab}\,.$$
%%An explicit computation of the required Christoffel symbols show that they vanish, so we have that $\II = 0$.
%
%Now we change coordinates so that $g$ is asymptotically in normal form, i.e. so that
%$$g = F^{(n+2)}(s,y^j) ds^2 + G_i^{(n+1)}(s,y^j) ds dy^i + \gamma_{ij}(s,y^1,\ldots,y^{d-1}) dy^i dy^j\,,$$
%where $n \geq 0$ is even and $F^{(n+2)}(s,y^j) = 1 + \mathcal{O}(s^{n+2})$ is an even function of $s$. Similarly, we require that $G_i^{(n+1)} = \mathcal{O}(s^{n+1})$ is an odd function of $s$.
%
%To find such a coordinate transformation, we solve order-by-order inductively. The base case is fairly easy, with $n=0$: define $t = \frac{s}{\sqrt{f(y^i)}}$ and $x^i = y^i$. Then, we have that
%$$g = ds^2 - \frac{s}{f^{3/2}} (\partial_{y^i} f) dy^i ds + \gamma_{ij} dy^i dy^j \,,$$
%for some $\gamma_{ij}$. Clearly, here $F^{(2)}(s,y^j) = 1$ is even and $G_i^{(1)}(s,y^j) = \frac{s}{f^{3/2}} \partial_{y^i} f$ is odd.
%We now use induction to proceed. 
%
%Suppose that $n \geq 0$ is even, that there exists $T_{n+1}(s,y^j) = \frac{s}{\sqrt{f(y^j)}} + \cdots + Q_{n+1}(y^j) s^{n+1}$ an odd polynomial in $s$, and $X^i_{n}(s,y^j) = y^i + \cdots + P^i_{n}(y^j) s^{n}$ an even function of $s$ such that
%$$g = F^{(n+2)}(s,y^j) ds^2 + G_i^{(n+1)}(s,y^j) ds dy^i + \gamma_{ij} dy^i dy^j$$
%when we use coordinates $t = T_{n+1}(s,y^j)$ and $x^i = X^i_{n}(s,y^j)$. When $n=0$, this is the base case.
%
%We now increment $n \rightarrow n+2$. Let $t = T_{n+1}(s,y^j) + s^{n+3} A_{n+3}(y^j)$ and $x^i = X^i_{n}(s,y^j) + s^{n+2} B^i_{n+2}(y^j)$, with $A_{n+3}$ and $B^i_{n+2}$ arbitrary functions of $\vec{y}$. We would like to show that there exist unique $A_{n+3}$ and $B^i_{n+2}$ such that $g_{ss} = F^{(n+4)}$ and $g_{si} = G^{(n+3)}_i$. To do so, we may compute components of the metric:
%\begin{align*}
%g_{si} =& G^{(n+1)}_i +  (n+2)  s^{n+1}\gamma_{jk} (\partial_{y^i} X_n^j) B^k_{n+2} \\
%&+ s^{n+2} \left((n+3) g_{tt} A_{n+3} \partial_{y^i} T_{n+1} + \gamma_{jk} (\partial_s X^j_n)\partial_{y^i} B^k_{n+2} \right) \\
%&+ s^{n+3} g_{tt} (\partial_s T_{n+1}) \partial_{y^i} A_{n+3} \\
%&+ (n+2) s^{2n+3} \gamma_{jk} B^j_{n+2} \partial_{y^i} B^k_{n+2} + (n+3) s^{2n+5} g_{tt} A_{n+3}  \partial_{y^i} A_{n+3} \,.
%\end{align*}
%Because $\gamma_{jk}$ is invertible, we may choose $B^i_{n+2}$ so that the leading term of second term cancels with the leading term of $G^{(n+1)}_i$. Because $G^{(n+1)}_i$ is odd and $X^j_n$ is even, the remaining higher-order terms are shunted to order $n+3$ and remain odd. Then, note that because both $T_{n+1}$ and $\partial_s X^j_n$ are odd in $s$, we see that the second line is actually odd and is of order $s^{n+3}$. Indeed, it follows that $g_{si}$ is odd and of order $s^{n+3}$. Now having chosen $B^i_{n+2}$, we may consider $g_{ss}$:
%\begin{align*}
%g_{ss} =& 1 + 2(n+2) s^{n+1} \gamma_{jk} B^j_{n+2} \partial_s X^k_n \\
%&+ F^{(n+2)}  + 2(n+3) s^{n+2} g_{tt} A_{n+3} \partial_s T_{n+1} \\
%&+ (n+2)^2 s^{2n+2} \gamma_{jk} B^j_{n+2} B^k_{n+2} + (n+3)^2 s^{2n+4} g_{tt} A_{n+3}^2 \,.
%\end{align*}
%Now because $B^i_{n+2}$ is fixed and $\partial_s X^k_n$ is odd, the second term on the first line is shunted to the second line (and remains even). Thus, we may choose $A_{n+3}$ uniquely so that the second line vanishes to leading order. But again, $\partial_s T_{n+1}$ is even as is $F^{(n+2)}$, so all higher-order terms remain even. Thus, the induction holds.
%
%
%Given such a coordinate transformation, we can now determine the parity of $\gamma_{ij}$. However, this is trivially the case: $\frac{dt}{dy^i}$ is odd, so $\frac{dt}{dy^i} \frac{dt}{dy^j}$ is even. Similarly, $\frac{dx^a}{dy^i} \frac{dx^b}{dy^j}$ is even. Therefore, we find that $\gamma_{ij}$ is an even function of $s$. Taking stock of this coordinate transform, we have that
%$$g = ds^2 + (\gamma^{(0)}_{ij} + s^2 \gamma^{(2)}_{ij} + \cdots) dy^i dy^j + \mathcal{O}(s^m)$$
%for any choice of $m$. Therefore, we have that in these coordinates, $n^a = \partial_s + \mathcal{O}(s^m)$, and so $\mathcal{L}_{n} T = \partial_s T + \mathcal{O}(s^{m-1})$ for any tensor $T$. So, for any $k \leq m$, we have that $\iota^* \mathcal{L}_n^k g = \iota^* \partial_s^k g$. Notably, for $k$ odd, we have that we may always find a coordinate transform so that
%$$\iota^* \partial_s^k g = 0\,.$$
%
%Now we would like to show that $\iota^* \mathcal{L}_n^k g$ is related to the $(k+2)$th fundamental form. By direct computation, we find that $\iota^* \mathcal{L}_n g = \II = 0$, and $\iota^* \mathcal{L}_n^2 g = \iota^* \nabla_n \nabla n = \III$ (which follows because $\II = 0$). Now, we would like to show via induction that $\mathcal{L}_n^k g = \nabla_n^{k-1} \nabla n + \text{subleading}$, where the subleading terms are products composed of fewer normal derivatives of $\nabla n$.
%% In particular, in the spirit of Remark~\ref{weights}, we assign $\nabla n$ a weight of $1$, $\nabla_n$ a weight of $-1$, and $g^{-1}$ a weight of $-2$.
%Now, we compute:
%\begin{align*}
%\iota^* \mathcal{L}_n^{k+1} g_{ab} =& \iota^* \mathcal{L}_n (\nabla_n^{k} \nabla_a n_b + \text{subleading}) \\
%=& \iota^* \left[ \nabla_n^{k+1} \nabla_a n_b + 2 (\nabla n)^c_{(a} \nabla_n^{k-1} (\nabla n)_{b)c} + \nabla_n (\text{subleading})_{ab} + 2 (\nabla n)^c_{(a} (\text{subleading})_{b)c} \right]\,.
%\end{align*}
%This proves the induction. Note that this induction also proves that, if we assign a weight of $1$ to $\nabla n$, a weight of $-1$ to $\nabla_n$, and a weight of $-2$ to $g^{-1}$, we also have that every summand has weight $2-k$. 
%
%But because $n = ds$ to arbitrarily high order, from Remark~\ref{weights}, it follows that
%$$\iota^* \mathcal{L}_n^{k} g \= \FF{k+1} + \text{subleading,}$$
%where the subleading terms are now those referred to in that remark. However, when $k$ is odd, we also have that the left hand side vanishes, so in particular it follows that for odd $k$, $\FF{k+1} = \text{subleading}$ and the subleading terms have weight $2-k$. But $2-k$ is odd, and odd fundamental forms always have even weight. Thus, it follows that the subleading terms cannot be composed entirely of odd fundamental forms. Hence, it follows that if $\FF{m}$ vanishes for each $m \leq k-1$ even, then it also follows that $\FF{k+1}$ vanishes. This completes the proof.
%\end{proof}
%
%It is interesting to note that, for a warped product manifold with metric $g = f^2 dt^2 + \bar{g}$, the third fundamental form of $\mathcal{Z}(t)$ satisfies
%$$\III = \frac{\bar{\nabla}^2 f}{f}\,.$$
%Furthermore, one can show \color{red} DO IT \color{black}
%\begin{align*}
%\V_{ab} =& \frac{\bar{\nabla}^c f}{f^2} \left(-3 \bar{\nabla}_c \bar{\nabla}_a \bar{\nabla}_b f + 2 \bar{\nabla}_{(a} \bar{\nabla}_{b)} \bar{\nabla}_c f \right) \\
%&+ \frac{1}{f^3} \left(3 |\bar{\nabla}f|^2 \bar{\nabla}_a \bar{\nabla}_b f - 2 (\bar{\nabla}^c f)(\bar{\nabla}_{(a} f) \bar{\nabla}_{b)} \bar{\nabla}_c f \right) \\
%=& \frac{\bar{\nabla}^c f}{f}(-3 \bar{\nabla}_c \III_{ab} + 2 \bar{\nabla}_{(a} \III_{b)c})\,.
%\end{align*}
%In particular, we find that $\III = 0$, then so does $\V$. This suggests that for a warped product manifold, the odd fundamental forms are fixed by linear differential operators on $\III$. To show this requires several results. First, we show that the vanishing of the third fundamental form is a particularly strong constraint.
%\begin{proposition} \label{warped-product-isometry2}
%Let $(M,g)$ be a warped product $M = \mathbb{R} \times_f \Sigma$ for some $f: \Sigma \rightarrow \mathbb{R}$ where $(\Sigma,\bar{g})$ is a connected complete Riemannian manifold. Furthermore, let $\III = 0$. Then, $M$ is isometric to $\mathbb{R} \times_f \left(\mathbb{R} \times \bar{\Sigma} \right)$, and there exists a choice of coordinates such that
%$$g = f^2 dt^2 + df^2 + \bar{\gamma}_{ij} dx^i dx^j\,.$$
%\end{proposition}
%\begin{proof}
%First, note that if $\III = 0$, then $\bar{\nabla}^2 f = 0$. Then, the isometry between $\Sigma$ and $\mathbb{R} \times \bar{\Sigma}$ follows directly from~\cite[Theorem 6.2]{wu2014}. The remainder of the result follows from the definition of $M$. \color{red} CHECK THIS \color{black}
%\end{proof}
%
%Note that from this proposition we easily get the following corollary.
%\begin{corollary}
%Let $(M^3,g)$ be a warped product $M = \mathbb{R} \times_f \Sigma^2$ for some $f: \Sigma \rightarrow \mathbb{R}$ with vanishing Hessian on $\Sigma$, a connected complete Riemannian manifold. Then, $(M,g)$ is isometric to either $\mathbb{R}^2 \times S^1$ or $\mathbb{R}^3$ equipped with the flat metric.
%\end{corollary}
%\begin{proof}
%From Proposition~\ref{warped-product-isometry}, it follows that $(M,g)$ is isometric to $\mathbb{R}^2 \times \Lambda$. However, because $\Sigma$ is connected and complete, $\Lambda$ is either topologically a circle or the real line. Either way, in some small neighborhood, the isometric condition implies that there exists a choice of coordinates such that
%$$g = f^2 dt^2 + df^2 + dy^2\,.$$
%However, this metric is just the flat metric in cylindrical coordinates, so the corollary follows.
%%
%%However, $f(x,y)$ is hypersurface Hessian, so in particular it follows that $\partial_{x}^2 f = \partial_x \partial_y f = \partial_y^2 f = 0$, so in these coordinates we have that $f(x,y) = Ax + By + C$. By explicit calculation, it follows that $R = 0$, and hence $g$ is locally flat. Furthermore, because the topology of $M$ is $\mathbb{R}^3$ and hence trivial, $g$ is globally flat.
%\end{proof}
%
%Next, it follows from Proposition~\ref{warped-product-isometry} that a vanishing third fundamental form means that a warped product is a product:
%\begin{proposition} \label{III0-prod}
%Let $(M,g)$ be a warped product $M = \mathbb{R} \times_f \Sigma$ where $\Sigma$ is a connected complete Riemannian manifold and $\III = 0$. Then, $(M,g)$ is a product manifold $\mathbb{R} \times \Sigma$.
%\end{proposition}
%\begin{proof}
%From Proposition~\ref{warped-product-isometry}, we may write
%$$g = f^2 dt^2 + df^2 + \bar{\gamma}_{ij} dx^i dx^j\,.$$
%But then the standard coordinate transformation to cartesian coordinates $f = \sqrt{u^2 + v^2}$ and $t = \arctan \tfrac{v}{u}$ yields a product metric (at least locally):
%$$g = dv^2 + du^2 + \bar{\gamma}_{ij} dx^i dx^j\,,$$
%and here $\Sigma = \mathcal{Z}(v)$. The proposition follows.
%\end{proof}
%
%But from this proposition, it follows that each fundamental form vanishes, and hence the odd fundamental forms must be expressible as linear operators on $\III$. We thus have the following theorem:
%\begin{proposition} \label{diff-op-FFs2}
%Let $(M,g)$ be a warped product $M = \mathbb{R} \times_f \Sigma$ where $\Sigma$ is a connected complete Riemannian manifold. Then, there exists an infinite family $\{O^f_n\}_{n=1}^{\infty}$ of differential operators on $\Sigma$ depending on $f$ such that
%$$\FF{2n+3} = O^f_n (\III)\,.$$
%Furthermore, this operator satisfies $O^f_n(0) = 0$.
%\end{proposition}
%\begin{proof}
%\color{red} check this carefully \color{black}
%First, note that $\FF{2n+3}$ has a local, differential formula expressible entirely in terms of $f$, $\bar{g}$, their partial derivatives, and $\bar{g}^{-1}$. Now, as such objects are tensorial, we  can instead write such objects using formulas constructed entirely from $f$, its (hypersurface) covariant derivatives, the metric $\bar{g}$ and its inverse, and the curvature tensor $\bar{R}$.
%
%Now in the case where $f$ is constant, we trivially have from Propositions~\ref{prod-vanishingFFs} and~\ref{III0-prod} that $\FF{2n+3} = 0$. Therefore, any such formula for $\FF{2n+3}$ cannot contain terms that only depend on $\bar{R}$ and $f$---i.e. all summands in such a formula must contain one or more derivatives of $f$. Thus, we may decompose any such formula into the following pieces:
%$$\FF{2n+3} = O(\bar{\nabla}^2 f) + P(\bar{\nabla} f)\,,$$
%where $O$ is some homogeneous differential operator and $P$ is a (potentially non-linear) homogeneous algebraic operator composed of the curvature tensor $\bar{R}$, its derivatives, the metric, its inverse, and $f$.
%
%Now consider the case where $\III = 0$, and thus $\bar{\nabla}^2 f = 0$. Because $\FF{2n+3} = 0$, it therefore follows that
%$$P(\bar{\nabla} f) = 0\,.$$
%However, $\bar{\nabla} f$ does not necessarily vanish (as it is independent of $\bar{\nabla}^2 f$). But because $P$ is an algebraic operator that vanishes for \textit{any} choice of $f$ satisfying $\bar{\nabla}^2 f$, we must have that $P = 0$ identically. It thus follows that $\FF{2n+3} = O(\bar{\nabla}^2 f) = O^f_n(\III)$. As $O$ is homogeneous, it trivially follows that $O^f_n(0) = 0$.
%%
%%Now, we suppose for contradiction that there exists some odd fundamental form $\FF{2n+3}$ such that there exists no differential operator $O^f_n$ such that $\FF{2n+3} = O^f_n(\III)$. But $\III = f^{-1} \bar{\nabla} \bar{\nabla} f$, and so in particular any differential formula for  $\FF{2n+3}$ that contains in each term at least one instance of $\bar{\nabla}^2 f$ necessarily can be written as an operator on $\III$ that depends on $f$. So, there must exist terms in this formula for $\FF{2n+3}$ that are expressible just in terms of $\bar{R}$ (and its covariant derivatives) and at most one derivative of $f$.
%%
%%Now consider the case where $\III = 0$. In that case, from Propositions~\ref{prod-vanishingFFs} and~\ref{III0-prod}, it follows that $\FF{2n+3} = 0$. However, $\FF{2n+3}$ depends only  on (derivatives of) $\bar{R}$ and at most a single derivative of $f$ (in a way where at least one term cannot be written in terms of $\bar{\nabla}^2 f$). But, $f$ and $\bar{R}$ are independent, so it must be the case that $\FF{2n+3}$ vanishes identically. Hence, it follows that $\FF{2n+3}$ is expressible as a differential operator on $\III$. The second claim follows as well.
%\end{proof}
%
%From this result, we finally have the following theorem.
%\begin{theorem}
%Let $\Sigma \hookrightarrow (M,g)$ be an analytic hypersurface embedding with $(\Sigma,\iota^* g)$ a connected complete Riemannian manifold. Further, suppose that $\FF{2n} = 0$ for each $n$, that there exists a positive analytic function $f : \Sigma \hookrightarrow \mathbb{R}$ satisfying $\bar{\nabla}^2 f = f \III$, and that $\FF{2n+3} = O^f_n (\III)$. Then, in a neighborhood $U$ of a point in $\Sigma$, $(U,g|_U)$ is diffeomorphic to a warped product manifold with base manifold $U \cap \Sigma$.
%\end{theorem}
%\begin{proof}
%Because $\Sigma \hookrightarrow (M,g)$ is a hypersurface embedding, in a neighborhood $U$ around a point in $\Sigma$ there exists a coordinate system $(t,\vec{x})$ such that
%$$g = ds^2 + \bar{g}(s,\vec{x})_{ij} dx^i dx^j\,.$$
%Furthermore, because the embedding is analytic, by shrinking $U$ if necessary, a series expansion of $\bar{g}(s,\vec{x})_{ij}$ about $s = 0$ converges. Now, in this normal form, each coefficient in the series expansion (except the zeroth order term) uniquely determines the fundamental forms. Equivalently, the metric $\bar{g}(0,\vec{x})$ and the set of fundamental forms uniquely (up to diffeomorphism) fixes $\bar{g}(s,\vec{x})$ in the neighborhood $U$, and hence the metric $g|_U$.
%
%Now, consider a warped product metric on $U$, given by
%$$\hat{g} = f(\vec{y})^2 dt^2 + \bar{\gamma}(\vec{y})_{ij} dy^i dy^j\,,$$
%where $\mathcal{Z}(t) = \Sigma \cap U$, $\{y^i\}$ are coordinates extended off $\mathcal{Z}(t)$ by parallel transport, $f$ is the function in the hypothesis in the theorem, and $\bar{\gamma} = \bar{g}|_{\Sigma}$. From Proposition~\ref{vanishing-evens}, it follows that $\hat{\FF{2n}} = 0$ for each positive $n$. By direct computation, it follows that $\hat{\III} = f^{-1} \bar{\nabla}^2 f = \III$. Furthermore, from Proposition~\ref{diff-op-FFs} it follows that each $\hat{\FF{2n+3}}$ is uniquely determined by $\hat{\III}$ and $\{O_k^f\}$, so $\hat{\FF{2n+3}} = \FF{2n+3}$.
%
%Now, we have two analytic metrics, $\hat{g}$ and $g$, on $U$ that have matching fundamental forms. But because they are analytic, the metric is uniquely (up to diffeomorphism) determined (in $U$) by the metric on $U \cap \Sigma$ and the higher fundamental forms. It therefore follows by uniqueness up to diffeomorphism that there exists a diffeomorphism $F : M \rightarrow M$ that preserves $\Sigma$ and that $\hat{g} = F^* g$, as required by the theorem.
%\end{proof}
%
%\begin{remark}
%The above theorem shows that warped product manifolds are strongly constrained---in particular, all higher fundamental forms vanish or are determined by $\III$ and the intrinsic geometric structure $(\Sigma,\bar{g})$.
%\end{remark}
%
%
%\section{Conclusion}
%In this article, we have provided several new results which suggest geometric interpretations of the higher fundamental forms. Evidently, the second fundamental form is an obstruction to curves with vanishing torsion that start on a hypersurface, staying on the hypersurface. This includes geodesics as a special case. The third fundamental form, on the other hand, controls conjugate points on hypersurfaces, or in other words obstructs hypersurface Jacobi fields from being bulk Jacobi fields. The fourth fundamental form is less obvious---but as we can frame the fundamental forms as a Taylor series expansion of the potential energy for the geodesic trajectory, a non-vanishing fourth fundamental form makes one side (or the other) of a hypersurface an attractor for geodesics. For the sake of simplicity, we decide to halt at the fourth fundamental form, but in principle one could continue to interpret the geometric meaning of the higher fundamental forms.
%
%While these interpretations make sense for hypersurfaces embedded in Riemannian manifolds, the very \textit{definitions} of higher fundamental forms are less clear for higher codimension submanifolds. The second fundamental form is easily defined---for a given orthonormal frame for the normal bundle to a submanifold $\{n_{\alpha}\}$, we can define $\II_{\alpha} := \iota^* \nabla n_{\alpha}$. As such, we propose that, for a warped product manifold with coordinates $(x^1, \cdots, x^m, y^1, \cdots, y^n)$, a higher fundamental form $\FF{k}_{\alpha_1 \cdots \alpha_{k-1}}$ is half the coefficient of $x^{\alpha_1} \cdots x^{\alpha_{k-1}}$ in the Taylor series expansion of the warping function times the metric on the base manifold. It is not yet clear how to generalize this notion to arbitrary submanifolds in an invariant way.
%
%%\begin{proposition}
%%Let $\Sigma^{d-1} \hookrightarrow (M^d,g)$ be an embedded hypersurface with every fundamental form vanishing. Then, for the conformal embedding $\Sigma^{d-1} \hookrightarrow (M^d,[g])$, we have that
%%$$\IIo = \IIIo = \cdots = \FFo{d-1} = 0\,.$$
%%\end{proposition}
%%\begin{proof}
%%Recall that a $k$th fundamental form is defined (for $k \geq 3$) according to
%%$$\FF{k}_{ab} = \top \nabla_{n}^{k-3} P_{ab}\,.$$
%%However, from~\cite{} it follows that, when the $k$th conformal fundamental form exists,
%%$$\FFo{k} \propto \otop \nabla_{n}^{k-3} P_{ab} + \text{lower order terms.}$$
%%Now, because $\II = 0$, it trivially follows that $\IIo = 0$. Thus, 
%%\end{proof}
%%
%%
%%\begin{proposition}
%%Let $\Sigma^{d-1} \hookrightarrow (M^d,g)$ be an embedded hypersurface with every fundamental form vanishing. Then, for any neighborhood $U$ containing a point $p \in \Sigma$, there exists a coordinate system such that
%%$$g = dr^2 + \bar{g} + r^2 \iota_* \bar{P} + \cdots\,.$$
%%\end{proposition}
%%\begin{proof}
%%For this embedding, we may always construct a defining function $r$ such that $|dr|_g = 1$ in a collar neighborhood of $\Sigma$. Because every fundamental form vanishes, it follows that the trace-free pieces of these fundamental forms vanish. Now  for any metric $\tilde{g} \in [g]$ the conformal class generated by $g$, the trace-free conformal fundamental forms vanish, $\IIo = \IIIo = \cdots = \FFo{d-1} = 0$. It follows from~\cite{} that such a conformal embedding $\Sigma \hookrightarrow (M,[g])$ corresponds to an asymptotically Poincar\'e--Einstein manifold, $\Sigma \hookrightarrow (M,g^o)$ where $g^o = \bg/\sigma^2$ for $\sigma = [g; r]$.
%%
%%Now, for any asymptotically Poincar\'e--Einstein metric $g^o$, there always exists a coordinate system such that (CITE RESULT!!!)
%%$$g^o = \frac{dr^2 + \bar{g} + r^2 \bar{P} + \cdots}{r^2}\,.$$
%%But then it follows easily that $g = r^2 g^o = dr^2 + \bar{g} + r^2 \bar{P} + \cdots$.
%%\end{proof}
%
%\section{To be sorted}
%
%
%
%
%
%%
%%\begin{proposition}
%%Let $\Sigma^{d-1} \hookrightarrow (M^d,g)$ with $d$ even be a totally geodesic hypersurface embedding with induced metric $\bar{g}$ such that $(\Sigma,\bar{g})$ is geodesically complete. Then, if $\III$ is negative definite, then for each point $p \in \Sigma$, there exists a conjugate point to $p$.
%%%and let $\gamma  : (-1,1) \rightarrow \Sigma$ be a unit speed geodesic of $(\Sigma,\bar{g})$ (and thus of $(M,g)$). If $\III$ is negative definite
%%%with $\delta$ the maximum eigenvalue of $\III$, maximizing over all points $p \in \Sigma$,
%%% then there exists a conjugate point to $\gamma(0)$ in $\Sigma$.
%%
%%
%%%then there exists some $\epsilon > 0$ and a geodesic $\rho: (-1,1) \rightarrow M$ satisfying $\rho(0) = \gamma(0)$, $\bar{g}^a_b \dot{\rho}^b(0) = \dot{\gamma}^a(0)$, and that $g(\dot{\gamma}(0),\dot{\rho}(0)) < \epsilon$ such that $\rho$ intersects $\Sigma$ at least twice.
%%\end{proposition}
%%
%%\begin{proof}
%%Note that if $\III$ is negative definite, then we must have any sectional curvature using the normal vector $n$ is positive definite. Alternatively, because $\III_{ab} = R_{nabn}$ when the embedding is totally geodesic, we have that
%%$$\III = \overline{Ric} - Ric^\top\,.$$
%%As $\III$ is negative definite, it follows that for any geodesic $\gamma$, we have that
%%$$ \overline{Ric}(\dot{\gamma},\dot{\gamma}) < Ric(\dot{\gamma},\dot{\gamma})\,.$$
%%
%%Now suppose that, viewed intrinsically, $(\Sigma,\bar{g})$ has no conjugate points and is not flat. Then, by a theorem of (Liang--Zhan 1996) we have that $\overline{Ric}$ is negative definite.
%%
%%
%%
%%If $\III$ is positive, then $Ric^\top(\dot{\gamma},\dot{\gamma}) < \overline{Ric}(\dot{\gamma},\dot{\gamma})$. If there are no conjugate points on $(\Sigma,\bar{g})$, then $(\Sigma,\bar{g})$ is flat.
%%
%%
%%STATEMENT: If $M$ has no conjugate points and $Ric \geq 0$, then $R = 0$.
%%
%%CONTRAPOSITIVE: If $R \neq 0$, then either $M$ has conjugate points or $Ric < 0$.
%%
%%\bigskip
%%
%%Suppose $\III$ is negative definite and $(\Sigma,\bar{g})$ is Ricci flat but not flat. Then, because $\III = \overline{Ric} - Ric^\top$, it follows that $Ric^\top$ is positive definite on $T \Sigma$, and in particular $R \neq 0$. It thus follows from the contrapositive of the above that points $p \in \Sigma \subset M$ have conjugate points.
%%
%%Now, because $(\Sigma,\bar{g})$ is not flat, we have that either $(\Sigma,\bar{g})$ has conjugate points or $\overline{Ric} < 0$. However, $(\Sigma,\bar{g})$ is Ricci flat, so there must exist conjugate points in $(\Sigma,\bar{g})$. 
%%
%%
%%
%%\bigskip
%%
%%
%%Assume $(\Sigma,\bar{g})$ is flat and that $\III \leq 0$. Now trivially, since $(\Sigma,\bar{g})$ is flat, it has no conjugate points. Furthermore, because $\III \neq 0$, we have that $(M,g)$ is not flat. Then, it follows from the contrapositive that either $M$ has conjugate points or $Ric < 0$. However, when $(\Sigma,\bar{g})$ is flat and geodesic, we have that $\III = -Ric^\top$, and so $Ric \geq 0$. But then it must be the case that $M$ has conjugate points.
%%
%%\bigskip
%%
%%Assume $(\Sigma,\bar{g})$ is flat and totally geodesic, but $\III$ is negative definite. Now trivially, since $(\Sigma,\bar{g})$ is flat, it has no conjugate points. Furthermore, because $\overline{Ric} = 0$, we have that $\III = -Ric^\top$, and so in particular, for any geodesic $\gamma$ in $\Sigma$, we have that $Ric^\top(\dot{\gamma},\dot{\gamma}) > 0$. But then by the contrapositive of a theorem proved by Liang and Zhan in 1996, we have that $\gamma$ has a conjugate point in $M$.
%%
%%\bigskip
%%
%%
%%
%%It then follows that $Ric^\top(\dot{\gamma},\dot{\gamma}) < 0$, and so either there are conjugate points or Ric
%%
%%https://www.scielo.br/j/aabc/a/mLGm9qJDCzrPHXSymqMTkPt/?format=pdf\&lang=en???
%%
%%%"if γ : R → M is a geodesic  without conjugate points and if Ric(γ(t), γ(t)≥ 0, then Ric(γ(t), γ(t)≡ 0."
%%
%%\bigskip
%%
%%We begin by picking $p \in \Sigma$. We must show that there exists a point $q \in \Sigma$ so that there is a geodesic connecting $p$ and $q$ and a non-vanishing Jacobi field along that geodesic that vanishes at both $p$ and $q$.
%%
%%First, pick any unit speed geodesic originating from $p$ lying in $\Sigma$, call it $\gamma$. Next, we define a family of geodesics (for fixed $s$) by $\Gamma(s,t) := \operatorname{exp}_{\gamma(0)}(tv_s)$, where $v_{s} \in T_{\gamma(0)} M$ is a unit vector satisfying $v_s \cdot \dot{\gamma}(0) = s > 0$ and $\bar{g}^a_b v_s^b = \alpha \dot{\gamma}(0)$ for some positive constant $\alpha$.
%%%Unless $\rho \in \Sigma$, we have that $\operatorname{span}(\dot{\gamma}(0), \dot{\rho}(0))$ generates a subspace $V_{\gamma(0)}$ of the tangent space at $\gamma(0)$ that is orthogonal to $\Sigma$. Then, we define a family of vectors $v_\theta = \dot{\gamma}(0) \cos \theta + \dot{\rho}(0) \sin \theta$. Now define $\Gamma(t,\theta) := \operatorname{exp}_{\gamma(0)}(t v_{\theta})$, so that $\Gamma(t,0) = \gamma(t)$ and $\Gamma(t,\pi/2) = \rho(t)$, and more generally $\Gamma(t,\theta)$ is a geodesic starting at $\gamma(0)$ with initial direction making an angle $\theta$ with respect to $\dot{\gamma}(0)$.
%%
%%Now, $\partial_{s} \Gamma|_{s=0} =: J(t)$ is a Jacobi field (Spivak IV.8.4) along $\gamma$. We now examine the projection of the Jacobi field to the normal bundle of the hypersurface.
%%The Jacobi equation is
%%$$\frac{D^2}{dt^2} J^c + J^a \dot{\gamma}^b \dot{\gamma}^d R_{ab}{}^c{}_{d} = 0\,.$$
%%Now, contracting with $n$, we find that
%%$$n \cdot \frac{D^2}{dt^2} J^c + R_{J \dot{\gamma} n \dot{\gamma}} = 0\,.$$
%%Integrating by parts we find that
%%$$\nabla_{\dot{\gamma}} g(n, \frac{D}{dt} J) - g(\frac{D}{dt} n, \frac{D}{dt} J) + R_{J,\dot{\gamma}n \dot{\gamma}} = 0\,.$$
%%
%%As $\Sigma$ is totally geodesic, we have that $\frac{D}{dt} n = 0$, so integrating by parts again, we have that
%%$$\nabla_{\dot{\gamma}}^2 (n \cdot J) + R_{J \dot{\gamma} n \dot{\gamma}} = 0\,.$$
%%
%%Now, suppose that we may decompose so that $J = q + A n$ such that $u$ is a Jacobi field along $\gamma$ lying entirely within $T \Sigma$. Then, we have that
%%$$\nabla_{\dot{\gamma}}^2 A + R_{u \dot{\gamma} n \dot{\gamma}} + A R_{n \dot{\gamma} n \dot{\gamma}} = 0\,.$$
%%However, as $u$ lies entirely within $T \Sigma$ and $\Sigma$ is totally geodesic, we have that $R_{u \dot{\gamma} n \dot{\gamma}} = 0$. Hence, we have the following equation controlling $A_s$:
%%$$\partial_t^2 A(t) + R_{n \dot{\gamma} n \dot{\gamma}} A(t) = 0\,.$$
%%Note that $\III(\dot{\gamma}, \dot{\gamma}) = R_{n \dot{\gamma} \dot{\gamma} n}$, so we may rewrite this equation as
%%$$\partial_t^2 A(t) - \III(\dot{\gamma}(t),\dot{\gamma}(t)) A(t) = 0\,.$$
%%
%%By negative definiteness of $\III$, we have that there is some $\delta < 0$ such that $0 < - \delta \leq -\III(\dot{\gamma}(t),\dot{\gamma}(t))$. Now the solutions of the differential equation
%%$$\partial_t^2 A(t) - \delta A(t) = 0$$
%% oscillate with period $T = \frac{2\pi}{\sqrt{-\delta}}$. By the Sturm--Picone comparison theorem, because $-\delta \leq \III(\dot{\gamma}(t),\dot{\gamma}(t))$, it follows that given an initial condition $A(0)=0$, we have that there exists $0<t_1 < \frac{2 \pi}{\sqrt{-\delta}}$ such that $A(t_1) = 0$.
%%
%%Now this is not sufficient to pr111111111ove that $\gamma(t_1)$ is a conjugate point, as $q(t_1)$ does not necessarily vanish. However, because $(\Sigma,\bar{g})$ is geodesically complete, there exists a geodesic connecting every point in $\Sigma$ to $\gamma(t_0)$. Now at $\gamma(0)$ there is a $d-2$ parameter family of unit tangent vectors, which define a $(d-2)$-dimensional family of unit-speed geodesics originating at $\gamma(0)$ parametrized by $\{e_i\}_{i=1}^{d-2} \in S^{d-2}$ and denoted by $\gamma_{\{e_i\}}$. We then define $G(\{e_i\},s,q,t)$ to be the family of geodesics linearly interpolating (between $q=0$ and $q=1$) from $\gamma_{\{e_i\}}$ to $\Gamma(t,s)$. As this is a family of geodesics, we have that $Q_{\{e_i\},s} := \partial_q G|_{q = 0}$ is a Jacobi field for each geodesic parametrized by $\{e_i\}$ and each inner product parametrized by $s$. Without loss of generality, we subtract off the component of the Jacobi field that is parallel to $\dot{\gamma}_{\{e_i\}}$.
%%
%%Note that by the same computation as above, $Q_{\{e_i\},s}$ has a vanishing normal component at some point $p$, where $p := \gamma_{\{e_i\}}(t_{\{e_i\},s})$ and $0< t_{\{e_i\},s} < \frac{2 \pi}{\sqrt{-\delta}}$. Now at such a point, we are interested in the behavior of the tangent vectors $\dot{\gamma}_{\{e_i\}}(t_{\{e_i\},s}) \in T_p \Sigma$ and the tangential component of the Jacobi field $q_{\{e_i\},s} \in T_p \Sigma / \langle \dot{\gamma}_{\{e_i\}}(t_{\{e_i\},s})  \rangle$. In particular, we would like to study the map
%%$$j : S^{d-2} \times \mathbb{R}_+ \rightarrow T \Sigma \times T \Sigma$$
%%that sends the parametrizing data to the tangent vector and tangential component of the Jacobi field when the normal component of the Jacobi field vanishes.
%%Note that the tangent vectors are arclength parametrized and the Jacobi field lies orthogonal to the tangent vector, so in particular the image of this map actually lies in $TS^{d-2}$.
%%
%%Our goal now is to show that $\pi \circ j$ is surjective, as the existence of a vanishing $q$ then follows from the hairy ball theorem. To do so, suppose not: then, there exists some point $\theta \in S^{d-2}$ such that for every $(\{e_i\},s)$, we have that $\dot{\gamma}_{\{e_i\}}(t_{\{e_i\},s}) \neq \theta$.
%%
%%
%%
%%For each $s$, the Borsuk-Ulam theorem then implies that there exists $\{\hat{e}_i\}$ such that $(\pi_1 \circ j)(\{\hat{e}_i\},s) = (\pi_1 \circ j)(\{-\hat{e}_i\},s)$. 
%%
%% and so at this point the tangential component of the Jacobi field $q_{\{e_i\},s} \in \mathbb{R}^{d-2}$ IS THIS TRUE???. We thus have a map
%%$$S^{d-2} \times \mathbb{R} \ni (\{e_i\},s) \mapsto q_{\{e_i\},s}(t_{\{e_i\},s}) \in \mathbb{R}^{d-2}\,.$$
%%As $d$ is even, the hairy ball theorem CHECK THE STATEMENT implies that, for any $s$ there exists at least one point $\{e_i\}$ that is sent to the zero vector. It thus follows that for any $s$ there exists $\{e_i\}$ such that the Jacobi field vanishes at $t_{\{e_i\},s}$. Hence, we have that $G(\{e_i\},s,0,t_{\{e_i\},s})$ is a conjugate point to $\gamma(0)$.
%%
%%%
%%% Now the tangential component $q_{\{e_i\},s}$  of $Q_{\{e_i\},s}$ obeys
%%%$$\bar{g} \left( q_{\{e_i\},s}, \frac{D^2}{dt^2} q_{\{e_i\},s} \right) + \bar{R}_{q_{\{e_i\},s} \dot{\gamma}_{\{e_i\}} q_{\{e_i\},s} \dot{\gamma}_{\{e_i\}} } = 0\,.$$
%%%Our goal is to show that there exists some $\{e_i\}$ such that $q_{\{e_i\},s}(t_{\{e_i\},s}) = 0$ for all $s$.
%%
%%
%%\end{proof}
%
%
%
%
%
%
%
%QUESTION:
%
%-- What do Jacobi fields look like for minimal and/or umbilic hypersurfaces? What do they say about the third fundamental form?
%
%Jacobi fields for minimal hypersurfaces: https://arxiv.org/pdf/1904.08423, Equation 2.31
%
%-- What do Jacobi fields look like for generalizations of circles? Does this say anything about the third fundamental form?
%
%-- If all fundamental forms vanish, is the manifold a product manifold between the hypersurface and a line?
%
%-- Normal form for Riemannian manifold when fundamental forms vanish, via PE?
%
%-- Equivalent of $\bar{J}$ and $J$ when $\III = 0$? What about higher codimension?
%
%-- Einstein $\Rightarrow$ higher fundamental forms are powers of second fundamental form.
%
%-- Einstein scale equation in bulk implies Einstein scale equation in hypersurface $\Leftrightarrow$ $\II = \III = 0$.
%
%-- Obstruction flat $\Rightarrow$ boundary obstruction flat iff higher fundamental forms vanish
%
%-- Jacobi fields for constant curvature zero torsion curves related to fourth fundamental form? What about projection of those curves intersecting geodesic hypersurface... does the projection of the Jacobi field also satisfy the relevant jacobi equation
%
%
%\bigskip
%
%-- Warped product... vanishing second fundamental form, but third fundamental form is not necessarily zero. Does vanishing third fundamental form imply all higher fundamental forms vanish in warped product setting? (With hypersurface as the base space)
%
%


\bibliographystyle{siam}
\bibliography{products-RiemConf-MOFM}


\end{document}

